% Réduction des conséquences néfastes des activités humaines sur les ressources naturelles — SVT 1ère C — Cours signé OnBuch+
% Programme officiel MINESEC. Compiler avec : tectonic NCT11-reduction-consequences-activites-humaines-ressources.tex
\newcommand{\DOCMATIERE}{SVT}
\newcommand{\DOCNIVEAU}{1ère C}
\newcommand{\DOCMODULE}{Module 3 : L'Éducation à l'Environnement et au Développement Durable}
\newcommand{\DOCLECON}{NCT11}
\newcommand{\DOCTITRE}{Réduction des conséquences néfastes des activités humaines sur les ressources naturelles}
% OnBuch+ — préambule « cours signés » ENRICHI (compilé avec tectonic / XeLaTeX)
% Système visuel « Pop » d'OnBuch+ : fond crème (jamais blanc), bordures encre
% épaisses, ombres « dures » décalées, titres Archivo Black en capitales.
% Version MATHÉMATIQUES : graphes (pgfplots/tikz), boîtes pédagogiques
% (définition, propriété, méthode, attention, le savais-tu, exercice résolu,
% exercice, TP, bilan), page de garde et signature OnBuch+ ancrée en bas.
%
% Macros attendues dans main.tex AVANT \input{preamble} :
%   \DOCMATIERE  \DOCNIVEAU  \DOCMODULE  \DOCLECON (numéro)  \DOCTITRE
\documentclass[11pt]{article}

\usepackage{fontspec}
\usepackage[french]{babel}
\usepackage[a4paper,margin=2cm,top=3.4cm,bottom=2.8cm,headheight=1.4cm,footskip=1.3cm]{geometry}
\usepackage{amsmath,amssymb,amsfonts}
\usepackage{mathtools} % \xmapsto, \coloneqq... souvent utilisés par les modèles
\usepackage{cancel} % simplifications barrées (bilans de forces)
% \abs{z} (module, valeur absolue) et \norm{u} (norme), utilisés par les modèles
\providecommand{\abs}{}\let\abs\relax\DeclarePairedDelimiter{\abs}{\lvert}{\rvert}
\providecommand{\norm}{}\let\norm\relax\DeclarePairedDelimiter{\norm}{\lVert}{\rVert}
\usepackage[table]{xcolor} % [table] : fournit aussi \rowcolors (alternance de lignes)
\usepackage{pagecolor}
\usepackage{tikz}
\usetikzlibrary{arrows.meta,positioning,calc,shapes.geometric,shapes.misc,shapes.arrows,decorations.pathmorphing,decorations.pathreplacing,decorations.markings,patterns,shadows,mindmap,trees,fit,backgrounds,matrix,intersections,angles,quotes}
\usepackage{pgfplots}
\usepackage[version=4]{mhchem} % équations nucléaires \ce{^{235}_{92}U}
\usepackage[european,siunitx]{circuitikz} % schémas électriques
\pgfplotsset{compat=1.18}
\usepgfplotslibrary{fillbetween,groupplots} % aires entre courbes (calcul intégral)
\usepackage{enumitem}
\usepackage{fancyhdr}
% [most] charge toutes les bibliothèques tcolorbox (listings, minted, raster,
% documentation...) et gonfle inutilement la mémoire TeX sur un document long
% (~300 boîtes) : on ne charge que ce qui est réellement utilisé.
\usepackage[skins,breakable]{tcolorbox}
\usepackage{graphicx}
\usepackage{colortbl}
\usepackage{booktabs}
\usepackage{tabularx}
\usepackage{diagbox} % case d'en-tête diagonale des tableaux à double entrée
% Colonnes extensibles centrée / gauche / droite (C, L, R), souvent utilisées par les modèles
\newcolumntype{C}{>{\centering\arraybackslash}X}
\newcolumntype{L}{>{\raggedright\arraybackslash}X}
\newcolumntype{R}{>{\raggedleft\arraybackslash}X}
\usepackage{array}
\usepackage{multicol}
\usepackage{siunitx}
% parse-numbers=false : accepte aussi \SI{2,5 \times 10^{-3}}{...}
\sisetup{locale=FR,output-decimal-marker={,},group-minimum-digits=4,parse-numbers=false}
\DeclareSIUnit\ohm{\ensuremath{\symup{\Omega}}} % U+2126 absent de la police
\DeclareSIUnit\division{div} % réglages d'oscilloscope (ms/div, V/div)
\DeclareSIUnit\tr{tr} % tours (vitesse de rotation, ex. \SI{1500}{\tr\per\minute})
\DeclareSIUnit\ch{ch} % cheval-vapeur (puissance moteur, unité usuelle hors SI)
\DeclareSIUnit\franccfa{FCFA} % franc CFA (coûts, contexte camerounais)
\DeclareSIUnit\dioptre{\ensuremath{\delta}} % vergence d'une lentille (dioptrie)
\usepackage{qrcode}
% \degree : commande fréquemment utilisée par les modèles pour un angle ou une
% température en dehors de \SI{}{} (non fournie par les paquets chargés).
\providecommand{\degree}{^{\circ}}
% \qed (fin de démonstration), utilisé par les modèles sans amsthm
\providecommand{\qed}{\hfill$\square$}
% \barre : double barre des valeurs interdites dans un tableau de variations (tkz-tab)
\providecommand{\barre}{\ensuremath{\|}}
% environnement proof (amsthm non chargé) utilisé par les modèles
\ifdefined\proof\else\newenvironment{proof}[1][Démonstration]{\par\noindent\textit{#1.}\ }{\qed\par}\fi
\usepackage{lastpage}
\usepackage{hyperref}
\hypersetup{hidelinks}

% ============================================================
% Palette « Pop » OnBuch+ — mêmes hex que la classe P (plus_ui.dart)
% ============================================================
\definecolor{poporange}{HTML}{F59321}
\definecolor{popdark}{HTML}{E07A0C}
\definecolor{popcream}{HTML}{FFF6E8}
\definecolor{popink}{HTML}{1C1714}
\definecolor{poppurple}{HTML}{7A5AE0}
\definecolor{popgreen}{HTML}{2E9E6A}
\definecolor{poppink}{HTML}{E0457B}
\definecolor{popblue}{HTML}{2F7DE1}
\definecolor{popgold}{HTML}{C98A12}
\definecolor{popmuted}{HTML}{8A7F76}
\definecolor{popline}{HTML}{EADFD2}
\definecolor{popgray}{HTML}{8A7F76}
\colorlet{popgrey}{popgray}
% Alias tolérants (le modèle nomme parfois les couleurs différemment)
\colorlet{popwhite}{white}
\colorlet{popblack}{popink}
\colorlet{popred}{poppink}
\colorlet{popyellow}{popgold}
% Teintes claires pour fonds de tableaux / zones
\colorlet{popblueL}{popblue!10!white}
\colorlet{poporangeL}{poporange!14!white}
\colorlet{popgreenL}{popgreen!12!white}
\colorlet{poppurpleL}{poppurple!12!white}
\colorlet{poppinkL}{poppink!12!white}
\colorlet{popgoldL}{popgold!14!white}

\pagecolor{popcream}

% ============================================================
% Typographie
% ============================================================
\defaultfontfeatures{Ligatures=TeX}
\setmainfont{PlusJakartaSans-Medium.ttf}[Path=../../../fonts/,BoldFont=PlusJakartaSans-Bold.ttf]
% Maths en sans-serif (Fira Math) avec chiffres et lettres droites de Plus
% Jakarta, pour que les formules se fondent dans le texte.
\usepackage[math-style=ISO,bold-style=ISO]{unicode-math}
\setmathfont{FiraMath-Regular.otf}[Path=../../../fonts/]
\setmathfont{PlusJakartaSans-Medium.ttf}[Path=../../../fonts/,range={up/num,up/{Latin,latin},\mathup}]
\usepackage{icomma} % virgule décimale sans espace en mode math
% Caractères mathématiques Unicode absents des polices de texte (Plus Jakarta,
% Archivo Black) : rendus par leur équivalent mathématique, en texte comme en
% formule — sinon ils s'affichent en carré vide (ex. « Arithmétique dans ℤ »).
\usepackage{newunicodechar}
\newunicodechar{ℝ}{\ensuremath{\mathbb{R}}}
\newunicodechar{ℤ}{\ensuremath{\mathbb{Z}}}
\newunicodechar{ℂ}{\ensuremath{\mathbb{C}}}
\newunicodechar{ℕ}{\ensuremath{\mathbb{N}}}
\newunicodechar{ℚ}{\ensuremath{\mathbb{Q}}}
\newunicodechar{𝔻}{\ensuremath{\mathbb{D}}}
\newunicodechar{α}{\ensuremath{\alpha}}
\newunicodechar{β}{\ensuremath{\beta}}
\newunicodechar{γ}{\ensuremath{\gamma}}
\newunicodechar{δ}{\ensuremath{\delta}}
\newunicodechar{ε}{\ensuremath{\varepsilon}}
\newunicodechar{θ}{\ensuremath{\theta}}
\newunicodechar{λ}{\ensuremath{\lambda}}
\newunicodechar{μ}{\ensuremath{\mu}}
\newunicodechar{σ}{\ensuremath{\sigma}}
\newunicodechar{τ}{\ensuremath{\tau}}
\newunicodechar{φ}{\ensuremath{\varphi}}
\newunicodechar{ψ}{\ensuremath{\psi}}
\newunicodechar{ω}{\ensuremath{\omega}}
\newunicodechar{Σ}{\ensuremath{\Sigma}}
% majuscules produites par \MakeUppercase dans les titres (α-aminés -> Α)
\newunicodechar{Α}{\ensuremath{\alpha}}
\newunicodechar{Β}{\ensuremath{\beta}}
\newunicodechar{Γ}{\ensuremath{\Gamma}}
\newunicodechar{Δ}{\ensuremath{\Delta}}
\newunicodechar{Ε}{\ensuremath{\varepsilon}}
\newunicodechar{Θ}{\ensuremath{\Theta}}
\newunicodechar{Λ}{\ensuremath{\Lambda}}
\newunicodechar{Μ}{\ensuremath{\mu}}
\newunicodechar{Τ}{\ensuremath{\tau}}
\newunicodechar{Φ}{\ensuremath{\Phi}}
\newunicodechar{Ψ}{\ensuremath{\Psi}}
\newunicodechar{Ω}{\ensuremath{\Omega}}
\newunicodechar{↦}{\ensuremath{\mapsto}}
\newunicodechar{∈}{\ensuremath{\in}}
\newunicodechar{∉}{\ensuremath{\notin}}
\newunicodechar{∧}{\ensuremath{\wedge}}
\newunicodechar{⇔}{\ensuremath{\Leftrightarrow}}
\newunicodechar{⟺}{\ensuremath{\Longleftrightarrow}}
\newunicodechar{⇒}{\ensuremath{\Rightarrow}}
\newunicodechar{⟹}{\ensuremath{\Longrightarrow}}
\newunicodechar{∘}{\ensuremath{\circ}}
\newunicodechar{∪}{\ensuremath{\cup}}
\newunicodechar{∩}{\ensuremath{\cap}}
\newunicodechar{≡}{\ensuremath{\equiv}}
\newunicodechar{⊂}{\ensuremath{\subset}}
\newunicodechar{⊥}{\ensuremath{\perp}}
\newunicodechar{∃}{\ensuremath{\exists}}
\newunicodechar{∀}{\ensuremath{\forall}}
\newunicodechar{∛}{\ensuremath{\sqrt[3]{\,}}}
\newunicodechar{∜}{\ensuremath{\sqrt[4]{\,}}}
\newunicodechar{⃗}{\ensuremath{{}^{\to}}}
\newunicodechar{ⁿ}{\ifmmode^{n}\else\textsuperscript{\ensuremath{n}}\fi}
\newunicodechar{ˣ}{\ifmmode^{x}\else\textsuperscript{\ensuremath{x}}\fi}
\newunicodechar{ᵇ}{\ifmmode^{b}\else\textsuperscript{\ensuremath{b}}\fi}
\newunicodechar{ʳ}{\ifmmode^{r}\else\textsuperscript{\ensuremath{r}}\fi}
\newunicodechar{ᵘ}{\ifmmode^{u}\else\textsuperscript{\ensuremath{u}}\fi}
\newunicodechar{ʸ}{\ifmmode^{y}\else\textsuperscript{\ensuremath{y}}\fi}
\newunicodechar{ᵃ}{\ifmmode^{a}\else\textsuperscript{\ensuremath{a}}\fi}
\newunicodechar{ᵉ}{\ifmmode^{e}\else\textsuperscript{\ensuremath{e}}\fi}
\newunicodechar{ᵏ}{\ifmmode^{k}\else\textsuperscript{\ensuremath{k}}\fi}
\newunicodechar{ᵖ}{\ifmmode^{p}\else\textsuperscript{\ensuremath{p}}\fi}
\newunicodechar{ᴺ}{\ifmmode^{N}\else\textsuperscript{\ensuremath{N}}\fi}
\newunicodechar{ᶜ}{\ifmmode^{c}\else\textsuperscript{\ensuremath{c}}\fi}
\newunicodechar{ʰ}{\ifmmode^{h}\else\textsuperscript{\ensuremath{h}}\fi}
\newunicodechar{ᵠ}{\ifmmode^{q}\else\textsuperscript{\ensuremath{q}}\fi}
\newunicodechar{ₙ}{\ifmmode_{n}\else\textsubscript{\ensuremath{n}}\fi}
\newunicodechar{ₐ}{\ifmmode_{a}\else\textsubscript{\ensuremath{a}}\fi}
\newunicodechar{ᵢ}{\ifmmode_{i}\else\textsubscript{\ensuremath{i}}\fi}
\newunicodechar{ᵣ}{\ifmmode_{r}\else\textsubscript{\ensuremath{r}}\fi}
\newunicodechar{ₖ}{\ifmmode_{k}\else\textsubscript{\ensuremath{k}}\fi}
\newunicodechar{ᵦ}{\ifmmode_{\beta}\else\textsubscript{\ensuremath{\beta}}\fi}
\newunicodechar{ₓ}{\ifmmode_{x}\else\textsubscript{\ensuremath{x}}\fi}
\newfontfamily\popdisplay{ArchivoBlack-Regular.ttf}[Path=../../../fonts/]
\newfontfamily\poplogo{SpaceGrotesk-Bold.ttf}[Path=../../../fonts/]
\newfontfamily\popsemi{PlusJakartaSans-SemiBold.ttf}[Path=../../../fonts/]
\newfontfamily\popxbold{PlusJakartaSans-ExtraBold.ttf}[Path=../../../fonts/]
\newfontfamily\popxboldls{PlusJakartaSans-ExtraBold.ttf}[Path=../../../fonts/,LetterSpace=3.0]

\setlist[enumerate]{leftmargin=1.7em,itemsep=5pt,topsep=4pt}
\setlist[itemize]{leftmargin=1.7em,itemsep=4pt,topsep=4pt,label={\color{poporange}\textbullet}}
\setlength{\parindent}{0pt}
\setlength{\parskip}{5pt}
\renewcommand{\arraystretch}{1.25}

% pgfplots : style Pop par défaut
\pgfplotsset{
  every axis/.append style={axis line style={popink,line width=1pt},tick style={popink},
    grid style={popline,line width=0.5pt},label style={font=\small},tick label style={font=\footnotesize}},
  popcurve/.style={poporange,line width=1.8pt,no markers,smooth},
}
% Même style disponible dans un \draw TikZ simple (hors pgfplots)
\tikzset{popcurve/.style={poporange,line width=1.8pt}}
\tikzset{popgrid/.style={popline,very thin}}
\colorlet{popcurve}{poporange} % le nom du style est parfois utilisé comme couleur

% ============================================================
% Pill / squiggle
% ============================================================
\newcommand{\poppill}[3]{%
  \tikz[baseline=-0.55ex]\node[draw=popink,line width=1.2pt,rounded corners=2.6pt,%
    fill=#2,inner xsep=6.5pt,inner ysep=3pt]{%
    \popxboldls\fontsize{7}{7}\selectfont\color{#3}\MakeUppercase{#1}};%
}
\newcommand{\popsquiggle}[1]{%
  \tikz[baseline=-0.4ex]{
    \draw[#1,line width=2.6pt,line cap=round]
      plot [smooth] coordinates {(0,0.09) (0.34,0.01) (0.68,0.21) (1.02,0.01) (1.36,0.10)};
  }%
}
% Mot-clé surligné (marqueur orange sous le texte)
\newcommand{\cle}[1]{{\popxbold\color{popdark}#1}}

% ============================================================
% En-tête / pied
% ============================================================
\pagestyle{fancy}
\fancyhf{}
\renewcommand{\headrulewidth}{0pt}
\renewcommand{\footrulewidth}{0pt}
\lhead{\raisebox{-0.15ex}{\poplogo\fontsize{13}{13}\selectfont\color{popink}OnBuch\color{poporange}+}}
\rhead{\poppill{\DOCMATIERE\ \textbullet\ \DOCNIVEAU\ \textbullet\ Leçon \DOCLECON}{white}{popink}}
\lfoot{\popxbold\fontsize{7.6}{7.6}\selectfont\color{popmuted}\MakeUppercase{Cours signé OnBuch+}\ \textbullet\ {\popsemi\DOCTITRE}}
\rfoot{\tikz[baseline=-0.6ex]\node[rounded corners=8.5pt,fill=popink,minimum height=17pt,minimum width=17pt,inner xsep=5pt,inner sep=0]{\popxbold\fontsize{8}{8}\selectfont\color{popcream}\thepage\,/\,\pageref*{LastPage}};}

% ============================================================
% Titres
% ============================================================
\newcommand{\chaptitre}[2]{%
  \begin{tcolorbox}[enhanced,colback=poporange,colframe=popink,boxrule=2.6pt,arc=11pt,
      drop shadow=popink,left=15pt,right=15pt,top=12pt,bottom=11pt,before skip=3pt,after skip=18pt]
    \poppill{#2}{popink}{popcream}\par\vspace{9pt}
    {\raggedright\hyphenpenalty=10000\exhyphenpenalty=10000\popdisplay\fontsize{22}{24}\selectfont\color{popcream}\MakeUppercase{#1}\par}
  \end{tcolorbox}
}
% Section numérotée : pastille orange avec le numéro + titre Archivo
\newcounter{popsec}
\newcommand{\coursec}[1]{%
  \stepcounter{popsec}\setcounter{popsub}{0}%
  \par\vspace{16pt}\noindent
  \begin{minipage}{\linewidth}
  \tikz[baseline=-0.7ex]\node[draw=popink,line width=1.6pt,fill=poporange,rounded corners=4pt,minimum size=22pt,inner sep=0]{\popdisplay\fontsize{12}{12}\selectfont\color{popcream}\thepopsec};%
  \hspace{8pt}{\popdisplay\fontsize{15}{17}\selectfont\color{popink}\MakeUppercase{#1}}\par
  \vspace{4pt}\noindent{\color{poporange}\rule{36pt}{3.6pt}}
  \end{minipage}\par\nopagebreak\vspace{8pt}
}
\newcounter{popsub}
\newcommand{\courssub}[1]{%
  \stepcounter{popsub}%
  \par\vspace{9pt}\noindent{\popxbold\fontsize{11.5}{13}\selectfont\color{popink}{\color{poporange}\thepopsec.\thepopsub}\hspace{6pt}#1}\par\nopagebreak\vspace{4pt}
}

% ============================================================
% Boîtes pédagogiques — même recette : fond blanc, bordure épaisse colorée,
% ombre dure encre, badge plein. Titre optionnel [#1].
% ============================================================
\tcbset{popbase/.style={breakable,enhanced,colback=white,boxrule=2.3pt,arc=9pt,
  shadow={3pt}{-3pt}{0pt}{popink},left=14pt,right=14pt,top=11pt,bottom=11pt,before skip=11pt,after skip=15pt,
  fontupper=\popsemi\fontsize{10.6}{14.5}\selectfont\color{popink},
  fontlower=\popsemi\fontsize{10.6}{14.5}\selectfont\color{popink}}}
% Coche dessinée (le glyphe ✓ n'existe pas dans les polices du document)
\newcommand{\popcheck}[1]{\tikz[baseline=-0.5ex]\draw[#1,line width=1.6pt,line cap=round,line join=round](-0.13,0.0)--(-0.04,-0.09)--(0.13,0.1);}
\providecommand{\coche}{\popcheck{popgreen}}
\providecommand{\resultat}[1]{\par\smallskip\noindent\fcolorbox{popgreen}{popgreenL}{\parbox{\dimexpr\linewidth-2\fboxsep-2\fboxrule\relax}{\textbf{Résultat.} #1}}\par\smallskip}
\newcommand{\popboxhead}[3]{\poppill{#1}{#2}{white}\ifx\relax#3\relax\else\hspace{6pt}{\popxbold\fontsize{10.6}{12}\selectfont\color{popink}#3}\fi\par\vspace{7pt}}

\newtcolorbox{aretenir}[1][]{popbase,colframe=popgreen,before upper={\popboxhead{À retenir}{popgreen}{#1}}}
\newtcolorbox{exemplebox}[1][]{popbase,colframe=poporange,before upper={\popboxhead{Exemple}{poporange}{#1}}}
\newtcolorbox{definition}[1][]{popbase,colframe=popblue,before upper={\popboxhead{Définition}{popblue}{#1}}}
\newtcolorbox{propriete}[1][]{popbase,colframe=poppurple,before upper={\popboxhead{Propriété}{poppurple}{#1}}}
\newtcolorbox{methode}[1][]{popbase,colframe=popgold,before upper={\popboxhead{Méthode}{popgold}{#1}}}
\newtcolorbox{attention}[1][]{popbase,colframe=poppink,before upper={\popboxhead{Attention}{poppink}{#1}}}
\newtcolorbox{experience}[1][]{popbase,colframe=popblue,colback=popblueL,before upper={\popboxhead{Expérience}{popblue}{#1}}}
% « Le savais-tu ? » : carte violette pleine (contexte, culture, Cameroun)
\newtcolorbox{savaistu}[1][]{popbase,colback=poppurple,colframe=popink,
  fontupper=\popsemi\fontsize{10.4}{14.2}\selectfont\color{white},
  before upper={\poppill{Le savais-tu\,?}{white}{poppurple}\ifx\relax#1\relax\else\hspace{6pt}{\popxbold\color{white}#1}\fi\par\vspace{7pt}}}
% Exercice résolu : énoncé en haut, \tcblower puis la solution sur fond crème
\newcounter{popexo}\newcounter{popexr}
\newtcolorbox{exoresolu}[1][]{popbase,colframe=poporange,colbacklower=popcream,
  segmentation style={popink,line width=1.2pt,dashed},
  before upper={\stepcounter{popexr}\popboxhead{Exercice résolu \thepopexr}{poporange}{#1}},
  before lower={\poppill{Solution}{popink}{popcream}\par\vspace{6pt}}}
% Exercice (non résolu en place) — la correction est dans la partie Corrigés
\newtcolorbox{exercice}[1][]{popbase,colframe=popink,
  before upper={\stepcounter{popexo}\popboxhead{Exercice \thepopexo}{popink}{#1}}}
% Corrigé
\newtcolorbox{corrige}[1][]{popbase,colframe=popgreen,colback=popgreenL,
  before upper={\popboxhead{Corrigé}{popgreen}{#1}}}
% Objectifs / prérequis / bilan
\newtcolorbox{objectifs}{popbase,colframe=popink,colback=poporangeL,
  before upper={\popboxhead{Objectifs de la leçon}{popink}{}}}
\newtcolorbox{prerequis}{popbase,colframe=popink,colback=popblueL,
  before upper={\popboxhead{Prérequis}{popblue}{}}}
\newtcolorbox{bilan}{popbase,colframe=popink,colback=white,boxrule=2.8pt,
  before upper={\popboxhead{Fiche bilan}{poporange}{L'essentiel à connaître pour l'examen}}}
% Figure encadrée avec légende
\newtcolorbox{popfigure}[1][]{enhanced,breakable=false,colback=white,colframe=popline,boxrule=1.4pt,arc=8pt,
  left=10pt,right=10pt,top=10pt,bottom=8pt,before skip=10pt,after skip=12pt,halign=center,
  lower separated=false,halign lower=center,
  fontlower=\popsemi\fontsize{9}{11.5}\selectfont\color{popmuted},
  #1}
\newcommand{\legende}[1]{\par\vspace{6pt}{\popsemi\fontsize{9}{11.5}\selectfont\color{popmuted}#1}}

% ============================================================
% Signature OnBuch+
% ============================================================
\newcommand{\poplogoinline}[1]{{\poplogo\fontsize{#1}{#1}\selectfont\color{popink}OnBuch\color{poporange}+}}
\newcommand{\signatureonbuch}{%
  \begin{tcolorbox}[enhanced,colback=popink,colframe=popink,boxrule=2.2pt,arc=10pt,
      drop shadow=poporange,left=14pt,right=14pt,top=10pt,bottom=10pt,before skip=0pt,after skip=0pt]
    \tikz[baseline=-0.7ex]\node[circle,fill=popgreen,draw=popcream,line width=1.3pt,minimum size=22pt,inner sep=0]{\popcheck{white}};%
    \hspace{9pt}\begin{minipage}[c]{0.62\linewidth}
      {\popxbold\fontsize{12}{13}\selectfont\color{popcream}Signé\ \textbullet\ {\poplogo OnBuch\color{poporange}+}}\par\vspace{2pt}
      {\raggedright\popsemi\fontsize{8}{10.5}\selectfont\color{popline}\MakeUppercase{Équipe pédagogique OnBuch+}\\\MakeUppercase{Conforme au programme officiel MINESEC}\par}
    \end{minipage}\hfill
    {\poplogo\fontsize{22}{22}\selectfont\color{popcream}OnBuch\color{poporange}+}
  \end{tcolorbox}%
}

% ============================================================
% Page de garde
% #1 titre · #2 matière · #3 niveau · #4 module · #5 n° leçon · #6 durée ·
% #7 description / objectifs (texte ou itemize)
% ============================================================
\newcommand{\pagegarde}[7]{%
  \thispagestyle{empty}
  \newgeometry{a4paper,margin=1.7cm,top=1.6cm,bottom=1.4cm}%
  \begin{tikzpicture}[remember picture,overlay]
    \fill[poporange!18] ([xshift=-3.2cm,yshift=-3.6cm]current page.north east) circle (4.6cm);
    \fill[poppurple!14] ([xshift=2.4cm,yshift=7.6cm]current page.south west) circle (3.8cm);
  \end{tikzpicture}%
  {\poplogo\fontsize{30}{30}\selectfont\color{popink}OnBuch\color{poporange}+}\hfill
  \raisebox{0.6ex}{\poppill{Cours signé \textbullet\ Premium}{popink}{popcream}}\par
  \vspace{1pt}\popsquiggle{poppurple}\par
  \vspace{8pt}
  \begin{tcolorbox}[enhanced,colback=poporange,colframe=popink,boxrule=2.8pt,arc=13pt,
      drop shadow=popink,left=18pt,right=18pt,top=12pt,bottom=13pt,after skip=10pt]
    \poppill{#2 \textbullet\ #3}{popink}{popcream}\hspace{5pt}\poppill{#4}{white}{popink}\par\vspace{8pt}
    {\popdisplay\fontsize{14}{14}\selectfont\color{popink}LEÇON #5}\par\vspace{4pt}
    {\raggedright\hyphenpenalty=10000\exhyphenpenalty=10000\popdisplay\fontsize{25}{27}\selectfont\color{popcream}\MakeUppercase{#1}\par}
    \vspace{8pt}
    {\popxbold\fontsize{9.5}{9.5}\selectfont\color{popink}\MakeUppercase{Durée conseillée : #6}}
  \end{tcolorbox}
  \begin{tcolorbox}[enhanced,colback=white,colframe=popink,boxrule=2.2pt,arc=10pt,
      drop shadow=popink,left=16pt,right=16pt,top=10pt,bottom=10pt,after skip=8pt]
    \poppill{Dans cette leçon}{poppurple}{white}\par\vspace{5pt}
    \popsemi\fontsize{9.3}{12.6}\selectfont\color{popink}#7
  \end{tcolorbox}
  \noindent\begin{minipage}[t]{0.487\linewidth}
    \begin{tcolorbox}[enhanced,colback=popgreen,colframe=popink,boxrule=2.2pt,arc=10pt,
        drop shadow=popink,left=11pt,right=11pt,top=10pt,bottom=10pt,height=3.1cm,valign=center]
      \tcbox[colback=white,colframe=popink,boxrule=1.3pt,arc=5pt,boxsep=0pt,nobeforeafter,
        left=3pt,right=3pt,top=3pt,bottom=3pt,valign=center]{\qrcode[height=1.9cm]{https://play.google.com/store/apps/details?id=com.luvvix.onbuch&hl=fr}}%
      \hspace{8pt}\begin{minipage}[c]{0.5\linewidth}
        {\popdisplay\fontsize{11}{12.5}\selectfont\color{white}\MakeUppercase{Télécharge OnBuch}}\par\vspace{4pt}
        {\raggedright\popsemi\fontsize{8.4}{11}\selectfont\color{white}Gratuit \textbullet\ Google Play\\Tuteur IA, annales, concours, résultats\par}
      \end{minipage}
    \end{tcolorbox}
  \end{minipage}\hfill
  \begin{minipage}[t]{0.487\linewidth}
    \begin{tcolorbox}[enhanced,colback=poppurple,colframe=popink,boxrule=2.2pt,arc=10pt,
        drop shadow=popink,left=11pt,right=11pt,top=10pt,bottom=10pt,height=3.1cm,valign=center]
      \tikz[baseline=-0.7ex]\node[circle,fill=white,draw=popink,line width=1.3pt,minimum size=1.6cm,inner sep=0]{\popdisplay\fontsize{24}{24}\selectfont\color{poppurple}+};%
      \hspace{8pt}\begin{minipage}[c]{0.6\linewidth}
        {\popdisplay\fontsize{11}{12.5}\selectfont\color{white}\MakeUppercase{Passe à OnBuch+}}\par\vspace{4pt}
        {\raggedright\popsemi\fontsize{8.4}{11}\selectfont\color{white}Tous les cours signés, Tuteur IA illimité, fiches PDF — onglet OnBuch+ de l'app\par}
      \end{minipage}
    \end{tcolorbox}
  \end{minipage}
  \vspace{8pt}
  \popcontenu
  \vfill
  \signatureonbuch
  \restoregeometry
  \newpage
}

% Rangée « Ce cours contient » (page de garde)
\newcommand{\poptile}[3]{% #1 couleur #2 gros texte #3 légende
  \begin{tcolorbox}[enhanced,colback=#1,colframe=popink,boxrule=2pt,arc=9pt,shadow={2.5pt}{-2.5pt}{0pt}{popink},
      left=6pt,right=6pt,top=9pt,bottom=9pt,halign=center,valign=center,height=1.75cm,width=\linewidth,nobeforeafter]
    {\popdisplay\fontsize{12.5}{13}\selectfont\color{white}\MakeUppercase{#2}}\par\vspace{3pt}
    {\centering\popsemi\fontsize{7.8}{9.5}\selectfont\color{white}#3\par}
  \end{tcolorbox}}
\newcommand{\popcontenu}{%
  \par\noindent{\popxboldls\fontsize{7.5}{7.5}\selectfont\color{popmuted}CE COURS SIGNÉ CONTIENT}\par\vspace{7pt}
  \noindent\begin{minipage}[t]{0.235\linewidth}\poptile{popblue}{Cours}{complet, illustré, pas à pas}\end{minipage}\hfill
  \begin{minipage}[t]{0.235\linewidth}\poptile{poporange}{Méthodes}{et exercices résolus}\end{minipage}\hfill
  \begin{minipage}[t]{0.235\linewidth}\poptile{poppink}{Exercices}{type examen + corrigés}\end{minipage}\hfill
  \begin{minipage}[t]{0.235\linewidth}\poptile{popgreen}{Fiche bilan}{l'essentiel à retenir}\end{minipage}\par}

% Sommaire
\newtcolorbox{sommaire}{enhanced,colback=white,colframe=popink,boxrule=2.2pt,arc=10pt,
  drop shadow=popink,left=18pt,right=18pt,top=13pt,bottom=13pt,before skip=6pt,after skip=20pt,
  before upper={\poppill{Sommaire}{poppurple}{white}\par\vspace{9pt}\popsemi\fontsize{10.6}{14.5}\selectfont\color{popink}}}

% Signature de fin de cours — poussée tout en bas de la dernière page
\newcommand{\findecours}{%
  \par\vfill\nopagebreak
  \begin{center}\popsquiggle{poporange}\end{center}\vspace{4pt}
  \signatureonbuch
}
\AtBeginDocument{\def\llbracket{\mathopen{[\mkern-3.5mu[}}\def\rrbracket{\mathclose{]\mkern-3.5mu]}}} % intervalles d'entiers (glyphe ⟦ absent de la police)
% Glyphes absents de Fira Math : redéfinis à partir de glyphes disponibles
\AtBeginDocument{%
  \def\setminus{\mathbin{\backslash}}\let\smallsetminus\setminus
  \def\vdots{\mathord{\vbox{\baselineskip3.2pt\lineskiplimit0pt\kern4pt\hbox{.}\hbox{.}\hbox{.}}}}%
  \def\ddots{\mathinner{\mkern1mu\raise6.5pt\hbox{.}\mkern2mu\raise3.6pt\hbox{.}\mkern2mu\raise0.7pt\hbox{.}\mkern1mu}}%
  \def\triangle{\mathord{\scalebox{0.9}{$\symup{\Delta}$}}}%
  \def\nabla{\mathord{\raisebox{\depth}{\scalebox{1}[-1]{$\symup{\Delta}$}}}}%
  \def\checkmark{\popcheck{popdark}}%
  \def\star{\mathbin{\ast}}\def\bigstar{\mathord{\ast}}%
  \def\diamond{\mathbin{\scalebox{0.6}{$\Diamond$}}}%
  \def\bigcup{\mathop{\mathchoice{\raisebox{-0.3em}{\scalebox{1.7}{$\cup$}}}{\scalebox{1.25}{$\cup$}}{\cup}{\cup}}\displaylimits}%
  \def\bigcap{\mathop{\mathchoice{\raisebox{-0.3em}{\scalebox{1.7}{$\cap$}}}{\scalebox{1.25}{$\cap$}}{\cap}{\cap}}\displaylimits}%
  \def\lesseqqgtr{\mathrel{\lessgtr}}\def\bigoplus{\mathop{\oplus}\displaylimits}%
}
\providecommand{\ssi}{si et seulement si} % abréviation fréquente
\AtBeginDocument{\providecommand{\wideparen}{\overparen}} % arc de cercle

%% FIGFIX-BEGIN v5 — correctifs globaux des figures (docs/onbuch-generation/tools/figfix)
\usepackage{adjustbox}\usepackage{etoolbox}\tcbuselibrary{raster}
% toute figure TikZ/pgfplots plus large que la ligne est réduite à la largeur disponible
\BeforeBeginEnvironment{tikzpicture}{\begin{adjustbox}{max width=\linewidth}}
\AfterEndEnvironment{tikzpicture}{\end{adjustbox}}
% légendes pgfplots lisibles par-dessus les courbes
\pgfplotsset{every axis/.append style={legend style={fill=white,fill opacity=0.92,text opacity=1,draw=black!15,inner sep=3pt}}}
% étiquettes d'arêtes (syntaxe edge["..."]) avec fond blanc
\tikzset{every edge quotes/.append style={fill=white,inner sep=1.2pt}}
%% FIGFIX-END
\begin{document}

\pagegarde{Réduction des conséquences néfastes des activités humaines sur les ressources naturelles}{SVT}{1ère C}{Module 3 : L'Éducation à l'Environnement et au Développement Durable}{NCT11}{3 h}{Cette leçon étudie les principales ressources naturelles du Cameroun (eau, sols, forêts, biodiversité) et les pressions que les activités humaines exercent sur elles. Elle introduit la notion de développement durable et propose des démarches concrètes de gestion raisonnée illustrées par des exemples camerounais.
\vspace{4pt}
\begin{multicols}{2}\small
\begin{itemize}
\item À la fin de cette leçon, je sais définir une ressource naturelle et distinguer ressources renouvelables et non renouvelables
\item À la fin de cette leçon, je sais citer les principales ressources naturelles du Cameroun et leur répartition
\item À la fin de cette leçon, je sais identifier les pressions exercées par les activités humaines sur l'eau, les sols, les forêts et la biodiversité
\item À la fin de cette leçon, je sais analyser un document (carte, tableau, texte) portant sur l'exploitation d'une ressource naturelle au Cameroun
\item À la fin de cette leçon, je sais définir le développement durable et ses trois piliers
\item À la fin de cette leçon, je sais expliquer le principe de gestion raisonnée d'une ressource renouvelable
\end{itemize}\end{multicols}}

\begin{sommaire}
\begin{multicols}{2}
\begin{enumerate}
\item Les ressources naturelles : définition et typologie
\item Les pressions humaines sur l'eau et les sols
\item Les pressions humaines sur les forêts et la biodiversité
\item Le développement durable et la gestion raisonnée des ressources
\item Exemples camerounais d'exploitation et de conservation
\item Synthèse : analyser une situation et proposer des solutions durables
\item Activité d'intégration
\item Méthodes et exercices résolus
\item Exercices
\item Corrigés des exercices
\item Fiche bilan
\end{enumerate}
\end{multicols}
\end{sommaire}

\begin{objectifs}
\begin{itemize}
\item À la fin de cette leçon, je sais définir une ressource naturelle et distinguer ressources renouvelables et non renouvelables
\item À la fin de cette leçon, je sais citer les principales ressources naturelles du Cameroun et leur répartition
\item À la fin de cette leçon, je sais identifier les pressions exercées par les activités humaines sur l'eau, les sols, les forêts et la biodiversité
\item À la fin de cette leçon, je sais analyser un document (carte, tableau, texte) portant sur l'exploitation d'une ressource naturelle au Cameroun
\item À la fin de cette leçon, je sais définir le développement durable et ses trois piliers
\item À la fin de cette leçon, je sais expliquer le principe de gestion raisonnée d'une ressource renouvelable
\item À la fin de cette leçon, je sais proposer des mesures de gestion durable adaptées à une situation locale camerounaise
\item À la fin de cette leçon, je sais citer des exemples camerounais de conservation (parcs nationaux, forêts communautaires, reboisement)
\end{itemize}
\end{objectifs}

\begin{prerequis}
\begin{itemize}
\item Notions de base sur les écosystèmes et les chaînes alimentaires (classes antérieures)
\item Notion de pollution et ses grandes catégories (eau, air, sol)
\item Connaissance générale de la géographie physique du Cameroun (fleuves, zones forestières, savanes)
\item Notion de population et de croissance démographique
\item Lecture et interprétation de tableaux et de graphiques simples
\end{itemize}
\end{prerequis}

\newpage

\coursec{Les ressources naturelles : définition et typologie}

Dans la région de l'Est-Cameroun, un village voit sa rivière se tarir chaque saison sèche depuis que des exploitants forestiers ont déboisé les collines voisines. Les pêcheurs de la Sanaga constatent la raréfaction du capitaine, et à Yaoundé, l'eau de forage est parfois impropre à la consommation. Comment expliquer ces dégradations ? Quelles activités humaines en sont responsables ? Et surtout, quelles solutions concrètes peuvent permettre aux communautés camerounaises d'exploiter leurs ressources sans les épuiser ? C'est pour répondre à ces questions que nous commençons par poser les bases : qu'est-ce qu'une ressource naturelle, comment la classer, et que possède concrètement le Cameroun ?

\courssub{Définition et classification des ressources naturelles}

\begin{definition}[Ressource naturelle]
Une \cle{ressource naturelle} est un élément du milieu naturel (milieu physique ou vivant) que l'Homme prélève, utilise ou transforme pour satisfaire ses besoins matériels, énergétiques, culturels ou spirituels. Elle n'existe en tant que « ressource » que par la capacité technique et le besoin social de l'Homme à l'exploiter.
\end{definition}

Observons notre environnement immédiat : l'eau que nous buvons, le bois de nos tables, le sable de nos routes, le pétrole qui alimente nos groupes électrogènes, l'air que nous respirons. Tous ces éléments sont fournis par la nature sans intervention humaine initiale. Cependant, leur disponibilité dans le temps diffère radicalement. C'est le critère de \cle{renouvelabilité} qui structure la classification fondamentale.

\begin{popfigure}
\begin{tikzpicture}[
    level distance=1.8cm,
    sibling distance=3.5cm,
    every node/.style={draw, rounded corners, fill=popcream, thick, inner sep=4pt, font=\small},
    edge from parent/.style={draw, thick, popdark, ->},
    level 1/.style={sibling distance=5cm},
    level 2/.style={sibling distance=2.5cm}
]
\node [fill=popblue!20, font=\bfseries] {Ressources naturelles}
    child { node [fill=popgreen!20] {Renouvelables}
        child { node[align=center]{Eau douce\\ (Sanaga, Nyong, nappes)} }
        child { node[align=center]{Forêts\\ (Bassin du Congo, mangroves)} }
        child { node[align=center]{Sols\\ (Sols volcaniques Ouest, ferralitiques Sud)} }
        child { node[align=center]{Biodiversité\\ (Faune/Flore, parcs nationaux)} }
    }
    child { node [fill=poporange!20] {Non renouvelables (stock)}
        child { node[align=center]{Minerais\\ (Bauxite Minim-Martap, Fer Mbalam, Cobalt)} }
        child { node[align=center]{Pétrole\\ (Bassin Rio del Rey, Logone-Birni)} }
        child { node[align=center]{Gaz naturel\\ (Logbaba, Sanaga Sud)} }
        child { node[align=center]{Roches industrielles\\ (Calcaire Figuil, Pouma)} }
    };
\end{tikzpicture}
\legende{Classification des ressources naturelles avec exemples camerounais. Les ressources renouvelables se régénèrent à l'échelle humaine si le rythme de prélèvement est respecté ; les ressources non renouvelables constituent un stock fini.}
\end{popfigure}

\begin{definition}[Ressource renouvelable]
Une ressource est dite \cle{renouvelable} si elle est capable de se régénérer naturellement à l'échelle d'une vie humaine (quelques décennies), \textbf{à condition} que son rythme d'exploitation ne dépasse pas son rythme de régénération naturelle. Exemples : l'eau douce (cycle de l'eau), les forêts (photosynthèse, croissance), les sols (pédogenèse lente), la biodiversité (reproduction).
\end{definition}

\begin{attention}[Piège classique : « L'eau est inépuisable »]
\emph{Faux.} Une ressource renouvelable mal gérée peut disparaître \textbf{localement} ou devenir inutilisable.
\begin{itemize}
    \item \textbf{Eau :} Surexploitation des nappes (forages sauvages à Yaoundé/Douala) $\rightarrow$ abaissement de la nappe, intrusion saline côtière, puits asséchés.
    \item \textbf{Forêt :} Coupe rase sans reboisement $\rightarrow$ érosion, perte d'habitat, désertification locale (ex. : pourtour du Bénoué).
    \item \textbf{Sol :} Agriculture intensive sans amendement $\rightarrow$ lessivage, perte de fertilité (sols « fatigués » de l'Ouest).
\end{itemize}
Le caractère renouvelable est une \emph{potentialité}, pas une garantie absolue.
\end{attention}

\courssub{Condition de renouvelabilité : le rythme d'exploitation vs régénération}

Pour comprendre pourquoi une ressource renouvelable s'épuise, menons une démarche d'investigation simplifiée sur le cas de la forêt communautaire.

\begin{experience}[Modélisation simplifiée : stock forestier et prélèvement]
\textbf{Question :} Comment évolue le volume de bois d'une forêt si l'on prélève plus que l'accroissement naturel ?
\textbf{Hypothèses :}
\begin{enumerate}
    \item Si prélèvement $<$ accroissement : le stock augmente.
    \item Si prélèvement $=$ accroissement : le stock est stable (exploitation durable).
    \item Si prélèvement $>$ accroissement : le stock diminue jusqu'à l'épuisement.
\end{enumerate}
\textbf{Données (simulées pour 1 ha de forêt dense humide) :}
\begin{itemize}
    \item Accroissement naturel moyen (biomasse) : $\approx 5\ \text{m}^3/\text{an}$.
    \item Scénario A (Gestion durable) : Prélèvement $= 4\ \text{m}^3/\text{an}$.
    \item Scénario B (Surexploitation) : Prélèvement $= 12\ \text{m}^3/\text{an}$.
\end{itemize}
\textbf{Interprétation (calcul sur 20 ans) :}
\begin{align*}
\text{Scénario A :}\quad & \text{Stock final} = \text{Stock initial} + 20 \times (5 - 4) = \text{Stock initial} + 20\ \text{m}^3 \quad (\text{Augmentation}) \\
\text{Scénario B :}\quad & \text{Stock final} = \text{Stock initial} + 20 \times (5 - 12) = \text{Stock initial} - 140\ \text{m}^3 \quad (\text{Effondrement})
\end{align*}
\textbf{Conclusion :} La durabilité d'une ressource renouvelable repose sur l'inégalité fondamentale :
\[
\textbf{Prélèvement annuel} \le \textbf{Régénération annuelle}
\]
Toute politique forestière, hydraulique ou halieutique doit garantir cette inégalité.
\end{experience}

\begin{aretenir}[Condition de durabilité pour une ressource renouvelable]
Une ressource renouvelable ne reste renouvelable que si le \textbf{taux d'exploitation} (prélèvement par unité de temps) est \textbf{inférieur ou égal au taux de régénération naturelle} (accroissement, recharge, reproduction). Au-delà, le capital naturel s'amenuise : c'est de la \cle{consommation de capital} et non de l'exploitation de revenu.
\end{aretenir}

\courssub{Panorama des ressources naturelles du Cameroun}

Le Cameroun, souvent qualifié d'« Afrique en miniature » par sa diversité géologique, climatique et biologique, dispose d'un patrimoine naturel exceptionnel. Cette diversité est le socle de son économie (secteur primaire $\approx 15-20\%$ du PIB, emploi $\approx 50\%$ de la population active).

\begin{popfigure}
\begin{tikzpicture}[scale=0.85, every node/.style={font=\small}]
    % Contour approximatif Cameroun (simplifié)
    \draw[thick, popdark] (0,0) .. controls (1,1) and (3,3) .. (4,5) .. controls (5,6) and (6,6) .. (7,5) .. controls (8,4) and (9,3) .. (9.5,1) .. controls (9,0) and (7,-0.5) .. (5,0) .. controls (3,-0.5) .. (1,0) .. controls (0.5,0.5) .. (0,0);
    
    % Villes repères
    \node[circle, fill=popblue, text=white, inner sep=1.5pt] (Yde) at (3.5, 2.8) {};
    \node[above] at (Yde) {Yaoundé};
    \node[circle, fill=popblue, text=white, inner sep=1.5pt] (Dla) at (5.5, 1.5) {};
    \node[right] at (Dla) {Douala};
    \node[circle, fill=popblue, text=white, inner sep=1.5pt] (Gar) at (7.5, 4.2) {};
    \node[above] at (Gar) {Garoua};
    \node[circle, fill=popblue, text=white, inner sep=1.5pt] (Mar) at (8.5, 2.5) {};
    \node[right] at (Mar) {Maroua};
    \node[circle, fill=popblue, text=white, inner sep=1.5pt] (Ber) at (2.5, 4.5) {};
    \node[left] at (Ber) {Bertoua};

    % Fleuves (bleu)
    \draw[popblue, thick] (3.5, 2.8) .. controls (4.5, 2) and (5.5, 1.5) .. (5.5, 1.5); % Sanaga vers Douala
    \node[popblue, rotate=-20] at (4.5, 2.2) {Sanaga};
    \draw[popblue, thick] (2.5, 4.5) .. controls (3, 3.5) and (4, 2.5) .. (5.5, 1.5); % Nyong
    \node[popblue, rotate=-30] at (3.5, 3.5) {Nyong};
    \draw[popblue, thick] (7.5, 4.2) -- (8.5, 2.5); % Bénoué
    \node[popblue, rotate=-30] at (8, 3.5) {Bénoué};
    \draw[popblue, thick] (8.5, 2.5) -- (9.2, 1.2); % Logone
    \node[popblue, rotate=-40] at (8.8, 2) {Logone};
    \draw[popblue, thick] (1.5, 0.5) -- (3, 1.5); % Wouri
    \node[popblue] at (2, 1) {Wouri};

    % Zone Forestière (Vert - Sud/Est)
    \fill[popgreen!15, pattern=north east lines, pattern color=popgreen!50] 
        (0.5,0.5) .. controls (1,2) and (2,4) .. (3,5) .. controls (4,5.5) and (5,5) .. (6,4) .. controls (6.5,3) and (6,2) .. (5.5,1.5) .. controls (5,1) and (4,0.5) .. (3,0.5) .. controls (2,0.5) .. (1,0.5) .. controls (0.5,0.5) .. (0.5,0.5);
    \node[popgreen!70!black, align=center, font=\small\bfseries] at (3, 3.5) {Forêt dense\\ humide\\ (Bassin du Congo)};

    % Zone Soudanienne / Agriculture (Jaune - Nord/Centre)
    \fill[popgold!30, pattern=dots, pattern color=poporange] 
        (6,3.5) .. controls (7,4.5) and (8,4.5) .. (9,3.5) .. controls (9,2) and (8,1.5) .. (7,2) .. controls (6.5,2.5) .. (6,3.5);
    \node[poporange!80!black, align=center, font=\small\bfseries] at (7.5, 3) {Zone soudanienne\\ Agriculture coton,\\ céréales, élevage};

    % Zone Volcanique / Haute altitude (Brun - Ouest)
    \fill[popmuted!40, pattern=crosshatch dots, pattern color=popdark] 
        (2.5, 2) ellipse (1.2 and 0.8);
    \node[popdark, align=center, font=\small\bfseries] at (2.5, 2) {Hauts plateaux\\ Sols volcaniques\\ (Café, Arabica)};

    % Ressources minières / Pétrole (Symboles)
    \node[draw, diamond, fill=poppurple!20, inner sep=2pt] at (1.5, 4.8) {}; \node[left, poppurple] at (1.5, 4.8) {Fer (Mbalam)};
    \node[draw, diamond, fill=poppurple!20, inner sep=2pt] at (2.5, 1.5) {}; \node[below, poppurple] at (2.5, 1.5) {Bauxite (Minim-Martap)};
    \node[draw, star, star points=5, star point ratio=2.25, fill=poppink!50, inner sep=2pt] at (6.5, 1.2) {}; \node[below, poppink] at (6.5, 1.2) {Pétrole (Rio del Rey)};
    \node[draw, star, star points=5, star point ratio=2.25, fill=poppink!50, inner sep=2pt] at (8.5, 3.5) {}; \node[above, poppink] at (8.5, 3.5) {Pétrole (Logone-Birni)};
    \node[draw, star, star points=5, star point ratio=2.25, fill=poppink!50, inner sep=2pt] at (4.5, 3.5) {}; \node[above, poppink] at (4.5, 3.5) {Gaz (Logbaba)};

    % Légende manuelle
    \begin{scope}[shift={(10.5, 4.5)}]
        \draw[popblue, thick] (0,0.5) -- (1,0.5); \node[right] at (1,0.5) {Fleuves principaux};
        \fill[popgreen!15, pattern=north east lines, pattern color=popgreen!50] (0,-0.5) rectangle (1,0); \node[right] at (1,-0.25) {Forêt dense (40\% terr.)};
        \fill[popgold!30, pattern=dots, pattern color=poporange] (0,-1.5) rectangle (1,-1); \node[right] at (1,-1.25) {Zone agricole nord};
        \fill[popmuted!40, pattern=crosshatch dots, pattern color=popdark] (0,-2.5) rectangle (1,-2); \node[right] at (1,-2.25) {Sols volcaniques Ouest};
        \node[diamond, draw, fill=poppurple!20, inner sep=2pt] at (0.5, -3.5) {}; \node[right] at (1, -3.5) {Minerais};
        \node[star, star points=5, star point ratio=2.25, draw, fill=poppink!50, inner sep=2pt] at (0.5, -4.5) {}; \node[right] at (1, -4.5) {Pétrole / Gaz};
    \end{scope}
\end{tikzpicture}
\legende{Carte simplifiée du Cameroun : localisation des grandes ressources naturelles. La forêt dense humide couvre le Sud et l'Est (bassin du Congo). Les fleuves (Sanaga, Nyong, Bénoué, Logone, Wouri) drainent le territoire. Le Nord est une zone agricole et pastorale (coton, céréales). L'Ouest possède des sols volcaniques fertiles. Les ressources minières (fer, bauxite) et hydrocarbures (pétrole, gaz) sont localisées sur la carte.}
\end{popfigure}

\begin{savaistu}[Le Cameroun, poumon du bassin du Congo]
Le Cameroun abrite une partie significative du \textbf{bassin du Congo}, deuxième massif forestier tropical au monde après l'Amazonie. On y recense plus de \textbf{8 000 espèces végétales} (dont de nombreuses endémiques), près de 900 espèces d'oiseaux et plus de 300 espèces de mammifères (gorilles, chimpanzés, éléphants de forêt, bongos). Cette biodiversité exceptionnelle confère au pays une responsabilité planétaire en matière de régulation climatique (puits de carbone) et de conservation du patrimoine génétique mondial.
\end{savaistu}

\begin{exemplebox}[Chiffres clés : la ressource forestière camerounaise]
\begin{itemize}
    \item \textbf{Superficie forestière :} $\approx 20\ \text{millions d'hectares}$ (soit près de \textbf{40 \% du territoire national}).
    \item \textbf{Forêt de production permanente (UFA) :} $\approx 7,5\ \text{millions d'hectares}$ attribués en concessions forestières.
    \item \textbf{Forêt communautaire / communale :} Outil de gestion décentralisée en développement.
    \item \textbf{Contribution économique :} Le secteur forêt-bois représente le 2\textsuperscript{e} employeur après l'État et $\approx 6-8\%$ du PIB hors pétrole. Exportations : grumes, sciages, placages (vers Chine, UE, Vietnam).
\end{itemize}
\end{exemplebox}

\courssub{Importance économique et sociale : capital naturel et services écosystémiques}

Au-delà de la matière première (bois, eau, minerais), les ressources naturelles fournissent des \cle{services écosystémiques} indispensables à la vie et à l'économie, souvent invisibles dans les comptes nationaux classiques.

\begin{propriete}[Services écosystémiques (Classification MEA - Millennium Ecosystem Assessment)]
On distingue quatre catégories de services rendus par les écosystèmes (capital naturel) :
\begin{enumerate}
    \item \textbf{Services d'approvisionnement (Provisioning) :} Biens matériels prélevés. Ex : Eau potable (Sanaga, forages), bois d'œuvre/énergie, aliments (poisson, gibier, plantes médicinales), ressources génétiques.
    \item \textbf{Services de régulation (Regulating) :} Bénéfices issus de la régulation des processus naturels. Ex : \emph{Régulation du climat} (stockage carbone forêt du Sud), \emph{Purification de l'eau} (zones humides, sols forestiers filtrent la Sanaga), \emph{Pollinisation} (abeilles, chauves-souris pour café/cacao), \emph{Protection contre l'érosion} (couvert forestier sur pentes de l'Ouest).
    \item \textbf{Services culturels (Cultural) :} Bénéfices non matériels. Ex : \emph{Écotourisme} (Parcs de Waza, Korup, Bouba Ndjida, Lobéké), valeurs spirituelles (forêts sacrées), éducation, recherche scientifique.
    \item \textbf{Services de support (Supporting) :} Nécessaires à la production de tous les autres. Ex : \emph{Formation des sols} (pédogenèse), \emph{Cycle des nutriments}, \emph{Production primaire} (photosynthèse), \emph{Fourniture d'habitat}.
\end{enumerate}
La dégradation d'une ressource (ex. déforestation) entraîne la perte \textbf{simultanée} de plusieurs de ces services.
\end{propriete}

\begin{exoresolu}[Analyse des services écosystémiques d'une zone humide urbaine : le lac municipal de Yaoundé]
\textbf{Énoncé :} Le lac municipal de Yaoundé (créé artificiellement) est envahi par les jacinthes d'eau (\emph{Eichhornia crassipes}), reçoit des eaux usées non traitées et voit ses berges bétonnées. Identifiez les services écosystémiques perdus ou dégradés et proposez une mesure de restauration pour chacun.

\textbf{Solution détaillée :}
\begin{center}
\begin{tabularx}{\linewidth}{|l|X|X|}\hline
\textbf{Service écosystémique} & \textbf{État dégradé observé} & \textbf{Mesure de restauration / gestion durable} \\ \hline
\textbf{Épuration de l'eau (Régulation)} & Eutrophisation, odeurs, mortalité piscicole (manque $O_2$). & Création de zones tampons végétalisées (phragmites, typha) en amont pour filtrer les nitrates/phosphates ; station d'épuration en amont. \\ \hline
\textbf{Régulation des crues (Régulation)} & Berges bétonnées $\rightarrow$ écoulement rapide, risque inondation aval. & Renaturation des berges (pentes douces, végétation ripisylve) pour stocker l'eau et ralentir l'écoulement. \\ \hline
\textbf{Habitat / Biodiversité (Support)} & Monoculture de jacinthe $\rightarrow$ disparition poissons, oiseaux. & Élimination mécanique/biologique des jacinthes ; réintroduction espèces natives ; création d'îlots de repos. \\ \hline
\textbf{Loisirs / Éducation (Culturel)} & Insalubrité, interdiction baignade/pêche, paysage dégradé. & Aménagement paysager (promenades, observatoires), programmes éducatifs scolaires, valorisation esthétique. \\ \hline
\textbf{Approvisionnement (Provisioning)} & Eau impropre à tout usage (même arrosage). & Traitement permettant réutilisation pour arrosage espaces verts (économie eau potable). \\ \hline
\end{tabularx}
\end{center}
\textbf{Conclusion :} La restauration intégrée (hydrologique, écologique, paysagère) restaure un \textbf{bouquet de services} bien plus large que la simple « dépollution ».
\end{exoresolu}

\begin{methode}[Analyser la pression sur une ressource locale (Démarche type Probatoire)]
Pour répondre à une question du type « Identifier les impacts d'une activité sur une ressource » :
\begin{enumerate}
    \item \textbf{Identifier la ressource ciblée :} Eau (nappe/rivière), Sol, Forêt, Espèce, Paysage.
    \item \textbf{Qualifier l'activité humaine :} Type (agriculture, industrie, urbanisation, exploitation forestière, minière), échelle (artisanale/industrielle), localisation.
    \item \textbf{Analyser le mécanisme d'impact :} Prélèvement direct ? Pollution (chimique, thermique, sédimentaire) ? Fragmentation d'habitat ? Introduction d'espèces invasives ? Modification du régime hydrologique ?
    \item \textbf{Distinguier impact direct / indirect / cumulé :} Ex. Route forestière (direct: déforestation) $\rightarrow$ accès braconnage (indirect) $\rightarrow$ érosion + perte biodiversité (cumulé).
    \item \textbf{Évaluer la réversibilité :} Réversible à court terme (pollution organique) / long terme (métaux lourds) / irréversible (extinction espèce, perte sol).
    \item \textbf{Proposer des mesures :} \emph{Éviter} (changer localisation), \emph{Réduire} (meilleures techniques), \emph{Compenser} (reboisement, restauration zone humide), \emph{Réglementer} (quotas, aires protégées, plans de gestion).
\end{enumerate}
\end{methode}

\begin{exercice}[Application : Le cas de la bauxite de Minim-Martap (Adamaoua)]
Le projet d'exploitation de la bauxite de Minim-Martap (l'un des plus grands gisements mondiaux) prévoit l'ouverture d'une mine à ciel ouvert, la construction d'un chemin de fer vers le port de Kribi (environ 600 km) et d'une centrale hydroélectrique. 
\begin{enumerate}
    \item Classer la ressource « bauxite » (renouvelable / non renouvelable) et justifier.
    \item Identifier trois ressources naturelles \textbf{renouvelables} qui risquent d'être impactées par les infrastructures (mine, voie ferrée, barrage) dans cette zone de savane boisée et de galerie forestière.
    \item Pour l'une de ces ressources renouvelables, expliquer le mécanisme de l'impact et proposer une mesure d'atténuation conforme à la séquence « Éviter - Réduire - Compenser ».
\end{enumerate}
\end{exercice}

\begin{corrige}[Exercice n°1 : Bauxite de Minim-Martap]
\begin{enumerate}
    \item \textbf{Bauxite = Ressource non renouvelable (stock).} Justification : C'est un minerai (hydroxyde d'aluminium) formé par altération géochimique sur des millions d'années (paléosols). Son taux de formation géologique est nul à l'échelle humaine. Tout prélèvement diminue le stock définitif.
    \item \textbf{Ressources renouvelables impactées :}
        \begin{itemize}
            \item \textbf{Eau de surface / nappes :} Rivière Lom, affluents, nappes alluviales (pollution par boues, hydrocarbures, acidité ; modification hydrologie par barrage).
            \item \textbf{Sols / Couvert végétal :} Déforestation emprise mine/voie ferrée/barrage (érosion, perte fertilité, fragmentation).
            \item \textbf{Biodiversité (Faune/Flore) :} Perte d'habitat (grands mammifères, primates, oiseaux), barrière écologique (voie ferrée), braconnage facilité (accès).
            \item \textbf{Pâturages / Ressources fourragères :} Perte de surfaces pour l'élevage transhumant (conflits éleveurs/agriculteurs/miniers).
        \end{itemize}
    \item \textbf{Exemple : L'eau (Rivière Lom).}
        \begin{itemize}
            \item \textbf{Mécanisme :} Décapage stérile $\rightarrow$ mise à nu sols latéritiques $\rightarrow$ ruissellement intense $\rightarrow$ apport massif de matières en suspension (MES) et métaux (Fe, Al) dans la rivière $\rightarrow$ colmatage lits, toxicité faune aquatique, coût potabilisation aval.
            \item \textbf{Mesures :}
                \begin{itemize}
                    \item \emph{Éviter :} Ne pas stocker les stériles en zone inondable / près du lit mineur.
                    \item \emph{Réduire :} Bassins de décantation dimensionnés pour crues centennales ; traitement eaux de mine (neutralisation pH) ; revêtement voies d'accès pour limiter érosion.
                    \item \emph{Compenser :} Reboisement bassins versants amont ; restauration zones humides aval ; financement station traitement eau potable pour populations riveraines.
                \end{itemize}
        \end{itemize}
\end{enumerate}
\end{corrige}

Cette première section pose le vocabulaire et le cadre conceptuel indispensable. Nous avons vu que le Cameroun possède un capital naturel diversifié (eau, forêts, sols, biodiversité, minerais), mais que la distinction renouvelable / non renouvelable ne suffit pas : une ressource renouvelable n'est durable que si l'on respecte sa \textbf{vitesse de régénération}. La notion de \textbf{services écosystémiques} élargit la valeur de la nature bien au-delà de la matière première. Dans la section suivante, nous analyserons concrètement comment les activités humaines (agriculture, urbanisation, industrie) exercent leurs \textbf{pressions} sur l'eau et les sols, premières victimes de l'anthropisation.

\coursec{Les pressions humaines sur l'eau et les sols}

Dans la section précédente, nous avons défini les ressources naturelles et établi leur typologie : ressources renouvelables (eau, sols, forêts, biodiversité) et ressources non renouvelables (minerais, pétrole). Nous avons insisté sur un point essentiel : une ressource dite \emph{renouvelable} ne le reste que si le rythme de son exploitation n'excède pas le rythme de sa régénération. C'est précisément là que le bât blesse. Au Cameroun comme ailleurs, la croissance démographique (environ 2,6 \% par an), l'urbanisation rapide et l'intensification des activités agricoles et industrielles exercent sur l'eau et les sols des \cle{pressions} croissantes. Cette section a pour but d'identifier ces pressions, d'en expliquer les mécanismes et d'en mesurer les conséquences, afin de préparer le terrain aux solutions de gestion durable que nous étudierons ensuite.

\courssub{Les pressions humaines sur la ressource en eau}

\textbf{Une observation pour commencer.}\par
En avril 2017, puis à plusieurs reprises ensuite, la ville de Yaoundé a connu des flambées de choléra ; dans l'Extrême-Nord, les épidémies de choléra sont presque cycliques autour de Mayo-Danay et de Logone-et-Chari. Parallèlement, dans certaines zones de Douala, les forages qui donnaient de l'eau à 15 m de profondeur il y a vingt ans doivent aujourd'hui descendre bien plus bas. Comment expliquer simultanément une eau qui \emph{manque} et une eau qui \emph{rend malade} ? La réponse tient en deux types de pressions distinctes qu'il faut apprendre à distinguer.

\begin{attention}[Pollution ou surexploitation : ne pas confondre]
Il faut distinguer rigoureusement deux types de pression sur une ressource :
\begin{itemize}
\item la \cle{pollution} : dégradation de la \textbf{qualité} de la ressource (l'eau existe, mais elle est impropre à l'usage) ;
\item la \cle{surexploitation} : dégradation de la \textbf{quantité} de la ressource (on prélève plus vite que la ressource ne se renouvelle, les réserves s'épuisent).
\end{itemize}
Au Probatoire, confondre ces deux notions dans l'analyse d'un document est une erreur fréquente et pénalisée. Une nappe peut être à la fois surexploitée (son niveau baisse) \emph{et} polluée (par les nitrates agricoles, par exemple).
\end{attention}

\courssub{Les prélèvements excessifs et la baisse des nappes phréatiques}

L'eau douce utilisée par l'Homme provient des eaux de surface (rivières, lacs) et des eaux souterraines (\cle{nappes phréatiques}), alimentées par l'infiltration des eaux de pluie à travers le sol. Tant que les prélèvements (boisson, irrigation, industrie) restent inférieurs à la recharge naturelle de la nappe, l'équilibre est maintenu. Mais lorsque la population croît, que l'agriculture irriguée s'étend (par exemple les périmètres irrigués de la SEMRY à Yagoua pour la culture du riz) et que l'industrie pompe massivement, les prélèvements dépassent la recharge : le niveau de la nappe s'abaisse, les puits s'assèchent, et il faut forer toujours plus profondément, à des coûts croissants. Dans les zones côtières, la baisse des nappes peut provoquer une intrusion d'eau de mer salée dans les aquifères, rendant l'eau définitivement impropre à la consommation.

\courssub{Les pollutions de l'eau : domestique, industrielle et agricole}

\textbf{La pollution domestique.}\par
Dans les quartiers non desservis par un réseau d'assainissement (nombreux à Douala, Yaoundé ou Maroua), les \cle{eaux usées} domestiques (eaux de toilette, de lessive, de cuisine) sont rejetées directement dans les rigoles, les marigots ou le sol. Elles transportent des matières organiques, des détergents et surtout des micro-organismes pathogènes : vibrion cholérique (\emph{Vibrio cholerae}), bacille de la typhoïde (\emph{Salmonella typhi}), amibes, œufs de vers parasites. Lorsque ces eaux contaminent les puits, les sources ou les rivières utilisées pour la boisson, elles déclenchent des \cle{maladies hydriques} : choléra, fièvre typhoïde, dysenteries. Le choléra, par la diarrhée profuse qu'il provoque, peut tuer en quelques heures par déshydratation : c'est une urgence de santé publique directement liée à la qualité de l'eau.

\textbf{La pollution industrielle.}\par
Les rejets industriels (brasseries, tanneries, huileries, raffineries, ateliers mécaniques) contiennent des métaux lourds (plomb, mercure, cadmium), des hydrocarbures, des acides et des solvants. Rejetés sans traitement dans les cours d'eau — le cas des berges du Wouri à Douala est emblématique —, ils tuent la faune aquatique, rendent l'eau toxique et contaminent les poissons consommés par les populations riveraines (bioaccumulation des métaux lourds dans la chaîne alimentaire).

\textbf{La pollution agricole.}\par
L'intensification de l'agriculture repose sur les \cle{engrais} (apports d'azote et de phosphore) et les \cle{pesticides} (insecticides, herbicides, fongicides). Sous l'effet des pluies, une partie de ces produits n'est pas retenue par le sol : elle ruisselle vers les rivières ou s'infiltre vers les nappes. Les nitrates des nappes polluées sont dangereux pour les nourrissons (méthémoglobinémie), et les pesticides, persistants, s'accumulent dans les organismes aquatiques.

\begin{definition}[L'eutrophisation]
L'\cle{eutrophisation} est l'enrichissement excessif d'un milieu aquatique en nutriments (nitrates, phosphates), provoquant la prolifération d'algues puis l'asphyxie du milieu.
\end{definition}

\textbf{Mécanisme de l'eutrophisation, étape par étape.}\par
\begin{enumerate}
\item Les engrais agricoles (nitrates $\ce{NO3-}$, phosphates $\ce{PO4^{3-}}$) et les eaux usées domestiques enrichissent le lac ou la rivière en nutriments.
\item Ces nutriments provoquent une prolifération explosive des algues microscopiques : c'est le \emph{bloom} algial, qui donne à l'eau une couleur verte.
\item Les algues de surface font de l'ombre : les algues profondes et les plantes aquatiques ne reçoivent plus de lumière et meurent.
\item Les bactéries décomposent cette masse de matière organique morte en consommant l'oxygène dissous : la teneur en oxygène de l'eau s'effondre.
\item Poissons et autres organismes aquatiques meurent par asphyxie : le milieu devient quasiment mort, nauséabond et impropre à tout usage.
\end{enumerate}
Le phénomène est observable dans plusieurs plans d'eau urbains camerounais (par exemple le lac municipal de Yaoundé) et menace les zones côtières riches en nutriments.

\begin{methode}[Identifier l'impact d'une activité humaine sur une ressource]
Face à un document ou une situation au Probatoire, raisonne toujours selon la chaîne de causalité complète :
\begin{enumerate}
\item \textbf{Activité} : quelle action humaine ? (ex. épandage d'engrais)
\item \textbf{Modification du milieu} : que change-t-elle physiquement ou chimiquement ? (ex. enrichissement de l'eau en nitrates)
\item \textbf{Conséquence sur la ressource} : comment la ressource est-elle dégradée ? (ex. eutrophisation, mort des poissons)
\item \textbf{Conséquence pour l'Homme} : quel impact sanitaire, économique ou social ? (ex. perte de la pêche, eau impropre à la consommation)
\end{enumerate}
Une réponse qui saute directement de l'activité à la conséquence finale, sans expliciter le mécanisme intermédiaire, est une réponse incomplète.
\end{methode}

\courssub{Les pressions humaines sur les sols}

\textbf{Une observation pour commencer.}\par
Dans la nuit du 28 au 29 octobre 2019, après plusieurs jours de pluies diluviennes, un pan entier de colline s'est effondré sur le quartier Gouache à Bafoussam, dans la région de l'Ouest. Plus de 40 personnes ont perdu la vie, ensevelies dans leur sommeil. Les enquêtes ont montré que les maisons avaient été construites sur des pentes raides, déboisées et fragilisées. Cette catastrophe n'est pas une simple « fatalité naturelle » : elle illustre de manière tragique les conséquences des pressions humaines sur les sols.

\textbf{Les principales pressions sur les sols.}\par
\begin{itemize}
\item La \cle{déforestation} : abattage des arbres pour le bois de chauffe (première source d'énergie domestique au Cameroun), le bois d'œuvre, l'extension des champs et des habitations. Or les racines des arbres retiennent le sol et leur canopée amortit l'impact des gouttes de pluie.
\item L'\cle{agriculture sur brûlis} : le paysan défriche une parcelle, brûle la végétation, cultive quelques années jusqu'à épuisement du sol, puis abandonne la parcelle pour en défricher une autre. Pratique adaptée autrefois, lorsque les jachères duraient 15 à 20 ans, elle devient destructrice lorsque la pression démographique réduit la jachère à 2 ou 3 ans : le sol n'a plus le temps de reconstituer sa fertilité.
\item Le \cle{surpâturage} : dans le Grand Nord (régions de l'Extrême-Nord, du Nord et de l'Adamaoua), la concentration de troupeaux trop importants sur des parcours limités détruit le couvert herbacé ; le sol nu, piétiné et tassé, est livré au vent et aux pluies.
\item L'\cle{urbanisation non maîtrisée} : constructions anarchiques sur les pentes et les zones marécageuses, imperméabilisation des sols par le béton et le bitume (l'eau ne s'infiltre plus, elle ruisselle), carrières de sable et de latérite non réhabilitées.
\end{itemize}

\courssub{Le mécanisme de l'érosion}

\begin{definition}[L'érosion]
L'\cle{érosion} est l'enlèvement et le transport des particules de sol par l'eau ou le vent, accélérée par la destruction du couvert végétal.
\end{definition}

Décomposons le mécanisme en chaîne de causalité, conformément à la méthode vue plus haut :

\begin{enumerate}
\item \textbf{Disparition du couvert végétal} (déforestation, brûlis, surpâturage) : le sol est mis à nu. Les gouttes de pluie frappent directement la terre et éclatent les agrégats ; il n'y a plus de racines pour retenir les particules ni d'humus pour absorber l'eau.
\item \textbf{Ruissellement} : l'eau ne s'infiltre plus, elle s'écoule en surface en emportant les particules détachées. Sur les pentes, ce ruissellement creuse des rigoles puis des ravins (érosion en nappes, puis érosion en ravines, spectaculaire dans certaines zones de l'Adamaoua et du Nord).
\item \textbf{Perte de la couche fertile} : c'est l'horizon supérieur du sol, riche en humus et en éléments nutritifs, qui part en premier. Il ne reste qu'un sous-sol pauvre, compact, parfois induré (cuirasses latéritiques), sur lequel plus rien ne pousse.
\end{enumerate}

\begin{popfigure}
\begin{center}
\begin{tikzpicture}[
  box/.style={draw=popblue, fill=popblueL, rounded corners=3pt, align=center, font=\small, minimum height=1.1cm, text width=3.1cm},
  bad/.style={draw=poporange, fill=poporangeL, rounded corners=3pt, align=center, font=\small, minimum height=1.1cm, text width=3.1cm},
  arr/.style={-{Stealth[length=3mm]}, thick, popdark}
]
\node[box, align=center] (a) at (0,0){\textbf{Déforestation}\\ destruction du couvert végétal};
\node[bad, align=center] (b) at (4.6,0){\textbf{Érosion}\\ ruissellement, sol mis à nu};
\node[bad, align=center] (c) at (9.2,0){\textbf{Perte de fertilité}\\ départ de la couche humifère};
\node[bad, align=center] (d) at (13.8,0){\textbf{Baisse des rendements}\\ pauvreté, insécurité alimentaire};
\draw[arr] (a) -- (b);
\draw[arr] (b) -- (c);
\draw[arr] (c) -- (d);
\end{tikzpicture}
\end{center}
\legende{Chaîne de causalité reliant la déforestation à la baisse des rendements agricoles. Chaque maillon est un mécanisme explicitable : c'est cette chaîne complète qu'il faut restituer dans une analyse de document.}
\end{popfigure}

\textbf{Une investigation guidée : mesurer l'effet du couvert végétal sur l'érosion.}\par
Comment prouver scientifiquement que c'est bien la disparition de la végétation qui accélère l'érosion, et non une simple coïncidence ?

\begin{experience}[Effet du couvert végétal sur le ruissellement (expérience de cours)]
\textbf{Hypothèse} : la destruction du couvert végétal augmente le ruissellement et l'érosion du sol.\par
\textbf{Protocole} : on remplit deux bacs identiques avec le même sol, disposés sur la même pente. Le bac A est planté d'herbe dense ; le bac B est laissé nu. On arrose les deux bacs avec la même quantité d'eau, déversée au même rythme (pluie simulée), et on recueille l'eau de ruissellement dans des récipients gradués.\par
\textbf{Observations} : le bac B (sol nu) produit un volume d'eau de ruissellement nettement plus grand, et cette eau est fortement chargée en particules de terre (eau boueuse). Le bac A (sol couvert) produit peu de ruissellement, et l'eau recueillie est presque claire : une grande partie de l'eau s'est infiltrée.\par
\textbf{Interprétation} : à pente, sol et pluie identiques, seule la présence de végétation varie. C'est donc bien le couvert végétal qui favorise l'infiltration, freine le ruissellement et retient les particules du sol.\par
\textbf{Conclusion} : l'hypothèse est validée. Toute activité qui détruit le couvert végétal (déforestation, brûlis, surpâturage) accélère mécaniquement l'érosion.
\end{experience}

Ce même raisonnement s'applique à l'échelle d'un bassin versant réel. Lorsqu'une colline dominant une rivière est déboisée, l'eau qui s'infiltrait autrefois ruisselle désormais en entraînant des particules fines : la \cle{turbidité} de la rivière (sa teneur en matières en suspension, mesurée en NTU) augmente fortement, surtout lors des pluies.

\begin{popfigure}
\begin{center}
\begin{tikzpicture}
\begin{axis}[
  width=0.95\linewidth, height=6.5cm,
  xlabel={Temps (mois)},
  ylabel={Turbidité de l'eau (NTU)},
  xmin=0, xmax=24, ymin=0, ymax=55,
  xtick={0,6,12,18,24},
  ytick={0,10,20,30,40,50},
  legend style={at={(0.02,0.97)}, anchor=north west, font=\small, draw=popline},
  grid=major, grid style={popline!40},
]
\addplot[popcurve, domain=0:24, samples=200, color=popblue] {8 + 2*sin(deg(x*1.05)) + 1.5*sin(deg(x*3.7))};
\addlegendentry{Bassin versant forestier (témoin)}
\addplot[popcurve, domain=0:12, samples=100, color=poporange] {9 + 2*sin(deg(x*1.05)) + 1.5*sin(deg(x*3.7))};
\addplot[popcurve, domain=12:24, samples=150, color=poporange, forget plot] {9 + 2.4*(x-12) + 4*sin(deg(x*1.05)) + 2*sin(deg(x*3.7))};
\addlegendentry{Bassin versant déboisé au mois 12}
\draw[dashed, popdark] (axis cs:12,0) -- (axis cs:12,55);
\node[font=\small\itshape, text=popdark, anchor=west] at (axis cs:12.3,50) {déboisement};
\end{axis}
\end{tikzpicture}
\end{center}
\legende{Évolution (données simulées réalistes) de la turbidité de deux rivières comparables. Avant le mois 12, les deux bassins versants sont boisés et les turbidités sont faibles et proches. Après le déboisement du second bassin (mois 12), la turbidité de sa rivière augmente fortement : le sol nu est érodé à chaque pluie et les particules sont charriées vers le cours d'eau.}
\end{popfigure}

\courssub{Les conséquences de la dégradation des sols}

\textbf{Baisse des rendements agricoles.}\par
La couche fertile perdue met des siècles à se reconstituer (on estime qu'il faut de 100 à 1000 ans pour former quelques centimètres de sol). Les rendements chutent, les paysans compensent en défrichant de nouvelles parcelles forestières, ce qui relance l'érosion : c'est un \emph{cercle vicieux} qui appauvrit à la fois le sol et les populations.

\textbf{Désertification dans le Grand Nord.}\par
Dans l'Extrême-Nord camerounais, la combinaison du surpâturage, de la coupe de bois, de la pression démographique et de l'aridité croissante du climat provoque l'avancée de paysages désertiques : c'est la \cle{désertification}, c'est-à-dire la dégradation irréversible des terres en zone sèche. Le lac Tchad, dont la superficie a été divisée par plus de dix depuis les années 1960, illustre dramatiquement la fragilité de cet équilibre entre eau, sols et activités humaines.

\textbf{Glissements de terrain.}\par
Sur les reliefs accidentés de l'Ouest et du Nord-Ouest, la construction sur des pentes déboisées supprime les racines qui « cousaient » le sol à la roche. Sous l'effet de pluies intenses, le sol saturé d'eau devient lourd et glisse brutalement.

\begin{exemplebox}[Les glissements de terrain de Bafoussam (octobre 2019)]
Dans la nuit du 28 au 29 octobre 2019, des glissements de terrain ont détruit plusieurs habitations du quartier Gouache à Bafoussam, causant la mort de plus de 40 personnes. Analysons la chaîne de causalité : \textbf{activité} (constructions sur des pentes raides, déboisées et drainées par des eaux usées) $\rightarrow$ \textbf{modification du milieu} (sol privé de racines, saturé en eau par les pluies exceptionnelles) $\rightarrow$ \textbf{conséquence sur la ressource} (rupture de la stabilité du versant, glissement de masse de terre) $\rightarrow$ \textbf{conséquence pour l'Homme} (plus de 40 morts, familles sinistrées, quartier détruit). La catastrophe est donc \emph{naturelle dans son déclenchement} (la pluie) mais \emph{humaine dans ses causes profondes} (occupation non maîtrisée du sol).
\end{exemplebox}

\courssub{Le lien entre pression démographique et surexploitation des terres}

Pourquoi ces pressions s'aggravent-elles ? Le facteur d'amplification central est la \cle{pression démographique}. La population camerounaise est passée d'environ 4 millions d'habitants en 1950 à plus de 28 millions aujourd'hui. À techniques agricoles constantes, nourrir, loger et chauffer une population qui double en une génération impose :
\begin{itemize}
\item de cultiver davantage de terres, donc de raccourcir les jachères et de défricher la forêt ;
\item de concentrer hommes et bêtes sur les mêmes parcours (surpâturage) ;
\item de construire sur tous les espaces disponibles, y compris les pentes instables et les zones inondables ;
\item de prélever toujours plus d'eau et de bois de chauffe.
\end{itemize}
La ressource, elle, ne s'étend pas : la surface des terres cultivables et le volume des nappes restent limités. Lorsque la demande excède durablement la capacité de renouvellement, on parle de \cle{surexploitation}. C'est exactement la situation que le développement durable — objet de la section suivante — cherche à corriger en adaptant les pratiques aux capacités réelles des milieux.

\begin{exoresolu}[Analyser l'impact d'une activité humaine]
\textbf{Énoncé.} Dans un village de la région du Centre, une coopérative installe une vaste plantation de maïs sur une colline précédemment couverte de forêt secondaire, en bordure de la rivière qui alimente le village en eau de boisson. Les plants reçoivent des engrais azotés et des herbicides. Deux ans plus tard, les habitants constatent que l'eau de la rivière est devenue verdâtre et trouble après chaque pluie, et que les poissons se font rares. En suivant la méthode de la chaîne de causalité, explique ces observations.
\tcblower
\textbf{Solution détaillée.} On applique la méthode en quatre étapes, en distinguant les deux types de pression.
\begin{enumerate}
\item \textbf{Activités} : déforestation de la colline, mise en culture sur pente, épandage d'engrais azotés et d'herbicides.
\item \textbf{Modifications du milieu} : (a) disparition du couvert végétal $\rightarrow$ le sol nu ruisselle à chaque pluie et l'érosion entraîne des particules de terre vers la rivière (d'où l'eau trouble) ; (b) les nitrates et phosphates des engrais, lessivés par la pluie, enrichissent la rivière en nutriments.
\item \textbf{Conséquences sur la ressource} : (a) augmentation de la turbidité de l'eau (pollution de la qualité) ; (b) eutrophisation : prolifération des algues (eau verdâtre), puis décomposition bactérienne consommant l'oxygène dissous, d'où l'asphyxie et la disparition des poissons.
\item \textbf{Conséquences pour l'Homme} : eau de boisson dégradée (risque de maladies hydriques, coût de traitement accru), effondrement de la pêche (perte de revenus et de protéines), et à terme baisse des rendements du maïs lui-même par érosion de la couche fertile de la colline.
\end{enumerate}
\textbf{Conclusion} : les observations des habitants s'expliquent entièrement par la combinaison de l'érosion (pression physique sur le sol) et de la pollution agricole (pression chimique sur l'eau). Des mesures de gestion durable — bandes végétalisées en bord de rivière, cultures en courbes de niveau, réduction raisonnée des engrais — permettraient de limiter ces impacts ; elles seront étudiées dans la suite de la leçon.
\end{exoresolu}

\begin{aretenir}[L'essentiel de la section]
\begin{itemize}
\item Deux grandes pressions sur l'eau : la \cle{surexploitation} (quantité : baisse des nappes phréatiques) et la \cle{pollution} (qualité : domestique, industrielle, agricole), sources de maladies hydriques (choléra, typhoïde) et d'eutrophisation.
\item Quatre grandes pressions sur les sols : déforestation, agriculture sur brûlis, surpâturage, urbanisation non maîtrisée.
\item Mécanisme clé : disparition du couvert végétal $\rightarrow$ ruissellement $\rightarrow$ érosion $\rightarrow$ perte de la couche fertile $\rightarrow$ baisse des rendements, désertification (Grand Nord) et glissements de terrain (Bafoussam, 2019, plus de 40 morts).
\item La pression démographique amplifie toutes ces pressions en augmentant la demande sur des ressources limitées.
\item Méthode d'analyse : activité $\rightarrow$ modification du milieu $\rightarrow$ conséquence sur la ressource $\rightarrow$ conséquence pour l'Homme.
\end{itemize}
\end{aretenir}

\coursec{Les pressions humaines sur les forêts et la biodiversité}

Après avoir analysé les pressions exercées sur l'eau et les sols, nous tournons notre regard vers les poumons verts de notre pays : les forêts. Le Cameroun, avec ses quelque \SI{22}{millions d'hectares} de forêts (près de \SI{46}{\%} du territoire national), abrite l'un des patrimoines forestiers et biologiques les plus riches du bassin du Congo. Pourtant, ce patrimoine recule chaque année. Pourquoi ? Comment l'activité humaine fragilise-t-elle ces écosystèmes et les espèces qu'ils abritent ? C'est ce que nous allons comprendre par une démarche d'investigation rigoureuse.

\courssub{Les causes de la déforestation au Cameroun : une analyse systémique}

\begin{experience}[Observation : le recul de la forêt camerounaise]
\textbf{Mise en situation.} Regardons deux images satellite de la région du Sud (ex. : autour de Lolodorf ou de Campo) prises à 20 ans d'intervalle. On observe une fragmentation nette du couvert forestier : des taches claires (champs, pistes, zones urbanisées) percent la masse verte continue.
\textbf{Question déclenchante.} Quelles activités humaines, directes ou indirectes, expliquent cette transformation du paysage ?
\textbf{Hypothèses.} (1) L'exploitation du bois pour l'export ou le bois-énergie local. (2) L'agriculture (itinérante vivrière ou agro-industrielle). (3) Les infrastructures (routes, barrages, mines). (4) La croissance démographique qui amplifie les besoins.
\textbf{Données d'enquête (sources : MINEPDED, FAO, Global Forest Watch).} Au Cameroun, la perte de couverture forestière primaire est estimée entre \num{80000} et \num{100000} hectares par an ces dernières décennies. Les causes identifiées se répartissent ainsi : agriculture itinérante sur brûlis (\approx \SI{60}{\%}), exploitation forestière (industrielle + artisanale, \approx \SI{20-25}{\%}), extension des plantations agro-industrielles (palmier à huile, hévéa, \approx \SI{10}{\%}), infrastructures et bois-énergie (le reste).
\textbf{Interprétation.} La déforestation n'a pas une cause unique mais résulte d'un faisceau de pressions directes (coupe, feu, labour) portées par des causes indirectes (demande de terres, pauvreté, gouvernance foncière, marchés internationaux).
\textbf{Conclusion.} Comprendre la déforestation exige de distinguer \cle{causes directes} (proximales) et \cle{causes indirectes} (sous-jacentes, structurelles).
\end{experience}

\begin{definition}[Biodiversité]
La \cle{biodiversité} (diversité biologique) est l'ensemble des espèces vivantes, de leurs gènes et des écosystèmes qu'elles constituent. Elle s'apprécie à trois niveaux imbriqués : la diversité spécifique (nombre d'espèces), la diversité génétique (variabilité au sein d'une espèce) et la diversité des écosystèmes (forêts, savanes, zones humides, etc.).
\end{definition}

\begin{popfigure}
\begin{tikzpicture}[
    mindmap,
    concept color=popgreen!70,
    level 1/.style={level distance=4.5cm, sibling angle=90},
    level 2/.style={level distance=3.5cm, sibling angle=45},
    level 3/.style={level distance=2.8cm},
    every node/.style={font=\small, text width=3.2cm, align=center},
    edge from parent/.style={draw=popdark, thick}
]
\node[concept, text width=3.8cm] {Déforestation\\ au Cameroun}
    child[concept color=poporange!70] {node[concept, align=center]{Causes directes\\ (Proximales)}
        child {node[concept, concept color=poporange!50, align=center]{Agriculture itinérante\\ sur brûlis (\SI{60}{\%})}}
        child {node[concept, concept color=poporange!50, align=center]{Exploitation forestière\\ industrielle \& artisanale\\ (bois d'œuvre, \SI{20}{\%})}}
        child {node[concept, concept color=poporange!50, align=center]{Agro-industrie\\ (palmier, hévéa,\\ cacao \approx \SI{10}{\%})}}
        child {node[concept, concept color=poporange!50, align=center]{Bois-énergie\\ (charbon, bois de feu)}}
        child {node[concept, concept color=poporange!50, align=center]{Infrastructures\\ (routes, barrages,\\ mines, urbanisation)}}
    }
    child[concept color=popblue!70] {node[concept, align=center]{Causes indirectes\\ (Sous-jacentes)}
        child {node[concept, concept color=popblue!50, align=center]{Croissance démographique\\ \& pression foncière}}
        child {node[concept, concept color=popblue!50, align=center]{Pauvreté rurale\\ \& absence d'alternatives}}
        child {node[concept, concept color=popblue!50, align=center]{Gouvernance foncière\\ faible, corruption}}
        child {node[concept, concept color=popblue!50, align=center]{Demande internationale\\ (bois, huile de palme,\\ cacao, minerais)}}
        child {node[concept, concept color=popblue!50, align=center]{Politiques sectorielles\\ non intégrées}}
    };
\end{tikzpicture}
\legende{Carte mentale des causes directes (orange) et indirectes (bleu) de la déforestation au Cameroun. Les pourcentages sont des ordres de grandeur issus des rapports FAO/MINEPDED.}
\end{popfigure}

\textbf{Analyse détaillée des causes directes.}\par
\begin{itemize}
    \item \textbf{Agriculture itinérante sur brûlis (\emph{slash-and-burn}).} Pratique séculaire adaptée à faible densité de population : coupe-abattis-brûlis $\rightarrow$ culture 1-2 ans $\rightarrow$ jachère longue (15-20 ans) pour régénération. \textbf{Problème actuel :} raccourcissement drastique des jachères (parfois < 5 ans) par pression foncière $\rightarrow$ appauvrissement des sols, invasion par \emph{Chromolaena odorata} (herbe à la fièvre), échec de la régénération forestière.
    \item \textbf{Exploitation forestière.} Distinction cruciale : \textbf{exploitation légale} (titres forestiers, plans d'aménagement, rotation 30 ans, prélèvement sélectif de quelques tiges/ha) vs \textbf{exploitation illégale/artisanale non contrôlée} (coupe à blanc, non-respect des diamètres minimaux, pistes non refermées). La seconde ouvre la forêt à l'agriculture et au braconnage.
    \item \textbf{Agro-industrie.} Extension rapide des plantations de palmier à huile (ex. : Sud, Sud-Ouest, Littoral) et d'hévéa. Souvent sur forêts secondaires ou dégradées, mais parfois sur forêts primaires de haute valeur de conservation (HVC). Conflits fonciers fréquents avec communautés locales.
    \item \textbf{Bois-énergie.} Le bois (charbon de bois, bois de feu) fournit \approx \SI{75}{\%} de l'énergie domestique au Cameroun. Prélèvement intense dans les périphéries urbaines (Yaoundé, Douala, Bafoussam) créant des \cle{halos de déforestation} de \SI{50-100}{km} de rayon.
    \item \textbf{Infrastructures.} Routes (ex. : route Kribi-Lolabé, autoroute Yaoundé-Douala), barrages (Mekin, Lom Pangar, Nachtigal), mines (fer de Mbalam, cobalt de Nkamouna). Ouverture de pistes = accès facilité pour chasseurs, agriculteurs, exploitants illégaux.
\end{itemize}

\begin{attention}[Piège fréquent : amalgame exploitation forestière = déforestation]
L'exploitation forestière \textbf{légale, planifiée et certifiée} (ex. : FSC, PAFC) n'est pas synonyme de déforestation : c'est une \cle{gestion durable} (prélèvement < accroissement, maintien des fonctions écologiques). C'est l'exploitation \textbf{non contrôlée, illégale ou non renouvelée} (coupe à blanc, absence de réensemblancement) qui détruit le capital forestier. Au Probatoire, précise toujours : « exploitation \emph{non durable} » ou « \emph{non maîtrisée} ».
\end{attention}

\courssub{Conséquences écologiques de la déforestation}

\begin{methode}[Démarche : relier cause $\rightarrow$ mécanisme $\rightarrow$ conséquence]
Pour chaque conséquence ci-dessous, identifie le mécanisme biophysique mis en jeu.
\begin{enumerate}
    \item \textbf{Perte d'habitats et fragmentation.} La forêt continue devient un archipel de fragments isolés. $\rightarrow$ Effet de bord (microclimat modifié, vents, lumière), isolement des populations (dérive génétique, consanguinité), disparition des espèces à grand domaine vital (éléphant, gorille, chimpanzé).
    \item \textbf{Érosion hydrique accélérée.} Suppression du couvert $\rightarrow$ énergie cinétique des gouttes de pluie non interceptée $\rightarrow$ battance $\rightarrow$ ruissellement de surface $\rightarrow$ ravinement. Perte des horizons humifères riches en nutriments. Exemple : bassins versants de la Sanaga, du Nyong, du Wouri $\rightarrow$ envasement des barrages (Edéa, Song Loulou, Lom Pangar), perte de capacité de stockage, coûts de curage énormes.
    \item \textbf{Perturbation du cycle de l'eau.} Forêt = pompe biotique (évapotranspiration massive $\approx$ \SI{1200-1500}{mm/an} recyclés localement). Déforestation $\rightarrow$ baisse de l'évapotranspiration $\rightarrow$ diminution des précipitations locales et régionales, assèchement des sources en saison sèche, crues plus violentes en saison des pluies (perte de l'effet éponge du sol forestier).
    \item \textbf{Émissions de gaz à effet de serre (GES).} Forêt = puits de carbone (stock dans la biomasse aérienne, racinaire, sol). Déforestation $\rightarrow$ combustion/décomposition $\rightarrow$ relargage de CO$_2$, CH$_4$, N$_2$O. Le secteur « Changement d'affectation des terres et foresterie » (UTCATF) est le 1er émetteur de GES au Cameroun (\approx \SI{60-70}{\%} des émissions nationales). \textbf{Lien Probatoire :} la forêt est un \cle{puits de carbone} ; sa destruction transforme un puits en source nette d'émissions, amplifiant le changement climatique.
\end{enumerate}
\end{methode}

\begin{popfigure}
\begin{tikzpicture}
\begin{axis}[
    ybar,
    bar width=16pt,
    width=\linewidth,
    height=9cm,
    enlargelimits=0.15,
    ylabel={Perte annuelle (milliers d'hectares)},
    xlabel={Période},
    symbolic x coords={1990-2000, 2000-2010, 2010-2020, 2020-2023*},
    xtick=data,
    nodes near coords,
    nodes near coords align={vertical},
    every node near coord/.append style={font=\small, rotate=90, anchor=west, yshift=2pt},
    ymin=0,
    ymax=130,
    ymajorgrids=true,
    grid style=dashed,
    tick label style={font=\small},
    label style={font=\normalsize},
    legend style={at={(0.5,-0.18)}, anchor=north, legend columns=-1, font=\small},
    bar cycle list={
        {fill=popgreen!70, draw=popgreen!90},
        {fill=poporange!70, draw=poporange!90},
        {fill=popblue!70, draw=popblue!90},
    }
]
\addplot coordinates {(1990-2000, 85) (2000-2010, 92) (2010-2020, 98) (2020-2023*, 105)};
\addplot coordinates {(1990-2000, 50) (2000-2010, 55) (2010-2020, 60) (2020-2023*, 65)};
\addplot coordinates {(1990-2000, 15) (2000-2010, 18) (2010-2020, 22) (2020-2023*, 25)};
\legend{Agriculture itinérante, Exploitation bois (toutes formes), Agro-industrie + Infrastructures}
\end{axis}
\end{tikzpicture}
\legende{Évolution estimée de la perte forestière annuelle au Cameroun par cause principale (données synthétiques FAO/FRA et Global Forest Watch, ordres de grandeur en milliers d'hectares/an). *Période partielle. L'agriculture itinérante reste le premier moteur direct.}
\end{popfigure}

\begin{exemplebox}[Cas concret : le bassin du lac Nyos et la déforestation des pentes]
Les pentes du Mont Oku (région du Nord-Ouest), zone de captage du lac Nyos, ont subi une déforestation massive pour l'agriculture (pommes de terre, maïs) et le pâturage. Conséquences observées : (1) érosion gullying spectaculaire, (2) turbidité accrue des affluents, (3) modification de l'hydrologie locale. La reforestation (projet PNRN, ONG) avec espèces locales (\emph{Prunus africana}, \emph{Podocarpus}) vise à restaurer la fonction éponge et limiter l'apport de sédiments dans le lac, réduisant le risque de remise en suspension du CO$_2$ dissous en profondeur. Cet exemple illustre la chaîne : déforestation $\rightarrow$ érosion $\rightarrow$ risque géologique accru.
\end{exemplebox}

\courssub{Pressions sur la biodiversité : braconnage, pêche abusive, destruction d'habitats}

\begin{definition}[Braconnage]
Le \cle{braconnage} est la chasse ou la pêche illégale, notamment d'espèces protégées ou hors des périodes et zones autorisées. Il vise soit la viande de brousse (consommation locale/urbaine), soit les produits à haute valeur marchande (ivoire, écailles, peaux, trophées) pour le trafic international.
\end{definition}

\begin{savaistu}[Le pangolin : mammifère le plus trafiqué au monde]
Le Cameroun abrite trois espèces de pangolins (pangolin géant \emph{Smutsia gigantea}, pangolin à longue queue \emph{Phataginus tetradactyla}, pangolin à ventre blanc \emph{Phataginus tricuspis}). Ils sont intensément braconnés pour leurs écailles (médecine traditionnelle asiatique) et leur viande. Entre 2010 et 2020, des saisies au port de Douala et à l'aéroport de Yaoundé-Nsimalen ont révélé des tonnes d'écailles en partance pour l'Asie, représentant des dizaines de milliers d'individus. Le pangolin est inscrit à l'Annexe I de la CITES (interdiction totale du commerce international) depuis 2016. Sa disparition locale déséquilibre les écosystèmes (régulation des termites et fourmis).
\end{savaistu}

\begin{itemize}
    \item \textbf{Braconnage d'espèces emblématiques.}
    \begin{itemize}
        \item \textbf{Éléphant de forêt (\emph{Loxodonta cyclotis}).} Pour l'ivoire. Dans le bassin du Congo, la population est passée de \approx 150 000 (années 80) à < 100 000 aujourd'hui ; au Cameroun, elle est estimée entre 10 000 et 20 000 individus. Disparition = perte d'un \cle{ingénieur écosystémique} (dispersal de graines, création de clairières, points d'eau).
        \item \textbf{Grands singes : Gorille de l'Ouest (\emph{Gorilla gorilla}), Chimpanzé (\emph{Pan troglodytes}).} Viande de brousse, trafic de bébés (animaux de compagnie), perte d'habitat. Statut UICN : \emph{En danger critique} (gorille) et \emph{En danger} (chimpanzé).
        \item \textbf{Autres :} Buffle nain, Céphalophes (antilopes naines), Perroquets gris (Gabon), Tortues.
    \end{itemize}
    \item \textbf{Pêche abusive.} Dans les cours d'eau (Sanaga, Nyong, Logone, Lac Tchad) et en zone côtière (Kribi, Limbé, Douala). Pratiques destructrices : filets mailles trop petites (capture juvéniles), empoisonnement (produits chimiques, plantes ichthyotoxiques comme \emph{Tephrosia vogelii}), dynamitage, pêche en période de frai. $\rightarrow$ Effondrement des stocks halieutiques, perte de revenus pour pêcheurs artisanaux, insécurité alimentaire.
    \item \textbf{Destruction d'habitats.} Déforestation (vue plus haut), drainage des zones humides (riziculture, urbanisation), pollution (pesticides cotonnières dans le Nord, rejets industriels à Douala/Bonabéri, déchets plastiques), espèces envahissantes (\emph{Chromolaena odorata}, \emph{Eichhornia crassipes} / jacinthe d'eau sur le Wouri et le Lac Ossa).
\end{itemize}

\begin{aretenir}[Espèces menacées et Liste rouge de l'UICN]
L'Union internationale pour la conservation de la nature (UICN) publie la \cle{Liste rouge des espèces menacées}, référence mondiale. Elle classe les espèces en catégories de risque : EX (Éteinte), EW (Éteinte à l'état sauvage), CR (Danger critique), EN (En danger), VU (Vulnérable), NT (Quasi menacée), LC (Préoccupation mineure), DD (Données insuffisantes). Au Cameroun, le Ministère des Forêts et de la Faune (MINFOF) établit la liste nationale des espèces de classe A (protection intégrale), B (protection partielle), C (chasse réglementée), en cohérence avec l'UICN et la CITES. \textbf{Au Probatoire :} savoir citer des exemples d'espèces de classe A (Éléphant, Gorille, Chimpanzé, Pangolin, Perroquet gris, Lamantin) et expliquer le rôle de la Liste rouge (outil d'alerte, base pour lois nationales, priorisation conservation).
\end{aretenir}

\courssub{Conséquences de la perte de biodiversité}

\begin{experience}[Investigation : que perd-on quand une espèce disparaît ?]
\textbf{Observation.} Dans une zone de forêt exploitée puis abandonnée près de Bertoua, on note : disparition des gros frugivores (éléphants, grands singes, gros becs), pullulement de rongeurs et de petites antilopes, difficulté de régénération des arbres à grosses graines (\emph{Irvingia}, \emph{Detarium}, \emph{Baillonella}).
\textbf{Hypothèse.} La perte des dispersers de graines (zoochorie) bloque le cycle de reproduction des arbres à grosses graines.
\textbf{Données (expérience de terrain / littérature).} Expériences d'exclusion (filets) vs parcelles ouvertes : germination et survie des plantules d'\emph{Irvingia gabonensis} (bush mango) 5 à 10 fois supérieures en présence d'éléphants/gorilles. Les graines non dispersées meurent sous l'arbre mère (prédation, compétition, pathogènes : hypothèse de Janzen-Connell).
\textbf{Interprétation.} Disparition d'une espèce clé (éléphant) $\rightarrow$ cascade trophique $\rightarrow$ modification de la composition floristique future $\rightarrow$ forêt moins diversifiée, moins stockante en carbone, moins résiliente.
\textbf{Conclusion.} La biodiversité n'est pas une liste d'espèces : c'est un réseau d'interactions. Perdre une espèce, c'est perdre des fonctions écosystémiques.
\end{experience}

\begin{propriete}[Conséquences majeures de l'érosion de la biodiversité]
\begin{enumerate}
    \item \textbf{Déséquilibres écosystémiques :} rupture des réseaux trophiques, perte de la pollinisation (cultures : cacao, café, fruits), de la dispersion des graines, de la régulation naturelle des ravageurs (oiseaux, chauves-souris, insectes auxiliaires), de la fertilité des sols (macrofaune, micro-organismes).
    \item \textbf{Perte de ressources alimentaires et médicinales :} \SI{70-80}{\%} de la population camerounaise recourt à la pharmacopée traditionnelle (écorces de \emph{Prunus africana} contre l'hypertrophie bénigne de la prostate, \emph{Ancistrocladus korupensis} anti-VIH, \emph{Garcinia kola}, etc.). Érosion génétique des plantes cultivées (variétés locales de maïs, sorgho, igname, manioc) $\rightarrow$ vulnérabilité aux maladies et au changement climatique.
    \item \textbf{Appauvrissement du patrimoine national et mondial :} valeur intrinsèque, valeur optionnelle (futurs médicaments, gènes d'intérêt agronomique), valeur culturelle et spirituelle (forêts sacrées, totems), valeur économique (écotourisme : parcs de Waza, Bouba Ndjida, Korup, Campo Ma'an, Dja ; chasse sportive encadrée).
    \item \textbf{Aggravation du changement climatique :} forêts appauvries en biodiversité stockent moins de carbone et sont plus vulnérables aux sécheresses, incendies, parasites.
\end{enumerate}
\end{propriete}

\begin{exoresolu}[Analyse d'un document : l'exploitation du bois d'œuvre dans la TRIDOM]
\textbf{Énoncé.} Un rapport d'ONG (2022) signale que dans la zone TRIDOM (Tri-national Dja-Odzala-Minkébé, sud-est Cameroun), l'ouverture de routes forestières pour l'exploitation sélective du bois d'œuvre (ayous, sapelli, azobé) a facilité l'accès aux braconniers. La densité d'éléphants a chuté de \SI{75}{\%} en 15 ans le long des axes routiers. L'entreprise forestière a un plan d'aménagement certifié FSC, mais le contrôle sur le terrain est faible.
\textbf{Question.} Analyse les interactions entre exploitation forestière légale, infrastructures et braconnage. Propose deux mesures concrètes pour concilier exploitation et conservation.
\textbf{Solution détaillée.}
\begin{enumerate}
    \item \textbf{Analyse des interactions.}
    \begin{itemize}
        \item \textit{Cause directe :} routes forestières (pistes principales, secondaires) = corridors d'accès pour braconniers (motos, 4x4, pirogues sur rivières).
        \item \textit{Effet de bord :} fragmentation de l'habitat éléphant, perturbation des couloirs de migration transfrontaliers (Cameroun-Gabon-Congo).
        \item \textit{Faiblesse de la gouvernance :} certification FSC impose des mesures anti-braconnage (éco-gardes, barrières, sensibilisation), mais l'application est lacunaire (manque de moyens, corruption, zone vaste). L'exploitation \emph{légale} devient ainsi le vecteur indirect de l'exploitation \emph{illégale} de la faune.
    \end{itemize}
    \item \textbf{Mesures de gestion durable proposées.}
    \begin{enumerate}
        \item \textbf{Fermeture systématique et réhabilitation des pistes secondaires} après exploitation (remblaiement, barrières physiques, reboisement naturel assisté) pour empêcher l'accès permanent. Seules les pistes principales de débardage restent ouvertes, avec postes de contrôle permanents (éco-gardes MINFOF + guards entreprise).
        \item \textbf{Mise en place d'un système de surveillance participative et technologique :} brigades mixtes (MINFOF + entreprise + communautés riveraines) équipées de GPS/SMART (Spatial Monitoring and Reporting Tool), pièges photographiques sur les pistes, drones pour survols ciblés. Partage des revenus (redevances forestières, redevances cynégétiques) avec les communautés pour inciter à la dénonciation du braconnage (alternative économique : écotourisme, cultures pérennes durables, PFNL).
    \end{enumerate}
\end{enumerate}
\end{exoresolu}

\begin{attention}[Erreur fréquente : confondre « espèce menacée » et « espèce rare »]
Une espèce peut être \textbf{rare} (faible abondance naturelle, aire restreinte) sans être \textbf{menacée} (risque d'extinction imminent si pressions maintiennent). Inversement, une espèce abondante peut devenir menacée si la pression est forte (ex. : pangolin, perroquet gris). Au Probatoire, utilise les catégories UICN (CR, EN, VU) pour qualifier le statut de menace, pas l'adjectif « rare ».
\end{attention}

\courssub{Synthèse : les leviers pour réduire les pressions sur forêts et biodiversité}

\begin{aretenir}[Mesures clés de gestion durable des ressources forestières et fauniques]
\begin{itemize}
    \item \textbf{Aménagement du territoire :} Plan national d'affectation des terres (PNAT) respecté : zones de production (UFA, plantations) vs zones de protection intégrale (parcs, réserves) vs zones communautaires (forêts communautaires, zones de chasse villageoises). Connectivité écologique (couloirs) maintenue.
    \item \textbf{Forêt :} Lutte contre l'illégalité (traçabilité bois, contrôle indépendant, certification FSC/PAFC généralisée), promotion de l'exploitation à faible impact (EFI), reboisement/restauration (objectif Bonn Challenge / AFR100 : \SI{12}{millions ha} restaurés à l'horizon 2030), développement filières bois locale (transformation 2e/3e degré = valeur ajoutée + emplois + moins de pression sur bois brut).
    \item \textbf{Agriculture :} Intensification écologique (agroforesterie : cacao/café sous couvert, systèmes d'intercalaires, engrais verts) pour produire plus sur moins de surface $\rightarrow$ arrêt de l'extension des champs dans la forêt. Sécurisation foncière des paysans (certificats d'occupation, titres) pour inciter à l'investissement long terme.
    \item \textbf{Faune :} Application stricte loi 94/01 (régime forêts/faune/pêche) et CITES. Démantèlement réseaux trafic (coopération régionale COMIFAC, OCFSA). Développement élevages alternatifs (cane-rat/aulacode, poulet local, poisson) pour réduire demande viande de brousse. Valorisation économique faune vivante (écotourisme, chasse sportive villageoise encadrée, redevances reversées aux COVAREF).
    \item \textbf{Bois-énergie :} Promotion foyers améliorés (foyers « Mayo » ou « Kené » à haut rendement), charbonnages améliorés (four à chaudière, rendement \SI{30-35}{\%} vs \SI{10-15}{\%} traditionnel), plantations énergétiques à rotation courte (\emph{Eucalyptus}, \emph{Acacia}, \emph{Gliricidia}) en zones dégradées hors forêt.
    \item \textbf{Éducation et gouvernance :} Éducation environnementale (programmes scolaires, radios communautaires), participation effective des populations locales (CLIP, COVAREF, comités de gestion forêts communautaires), transparence revenus (ITIE - Initiative pour la transparence dans les industries extractives, adaptée au secteur forestier).
\end{itemize}
\end{aretenir}

\begin{exercice}[Entraînement Probatoire : Étude de cas - La plantation de palmier à huile à Campo]
\textbf{Document.} Une société agro-industrielle obtient une concession de \SI{60000}{ha} dans l'arrondissement de Campo (Sud), en bordure du Parc national de Campo Ma'an. La zone est un mélange de forêts secondaires, de jachères et de forêts primaires de haute valeur de conservation (couloirs éléphants, habitat gorilles/chimpanzés). L'étude d'impact environnemental et social (EIES) a été réalisée mais contestée par des ONG (sous-estimation impacts cumulés, consultation communautés insuffisante). L'entreprise s'engage à réserver \SI{30}{\%} de la concession en zones de conservation (HVC/HCV) et à développer un programme d'appui aux planteurs villageois (outgrowers).
\textbf{Questions.}
\begin{enumerate}
    \item Identifie deux pressions directes et deux pressions indirectes que ce projet risque d'exacerber sur la biodiversité du parc adjacent.
    \item Explique pourquoi la simple mise en réserve de \SI{30}{\%} de la concession ne garantit pas la connectivité écologique avec le parc.
    \item Propose trois mesures concrètes, applicables par l'entreprise et l'État, pour minimiser les impacts négatifs et favoriser la coexistence homme-faune.
\end{enumerate}
\end{exercice}

\begin{corrige}[Exercice : La plantation de palmier à huile à Campo]
\begin{enumerate}
    \item \textbf{Pressions directes :} (1) Destruction/fragmentation d'habitats critiques (couloirs de migration éléphants, zones de nidification grands singes) par le défrichage. (2) Augmentation massive du braconnage (main-d'œuvre nombreuse, pistes d'accès, marchés de viande de brousse). \textbf{Pressions indirectes :} (1) Afflux migratoire (chercheurs d'emploi, commerçants) $\rightarrow$ pression foncière accrue en périphérie, déforestation anarchique. (2) Faible gouvernance locale / corruption $\rightarrow$ non-respect des clauses EIES, trafic d'influence.
    \item \textbf{Limites des \SI{30}{\%} réservés :} (1) Morcellement : si les \SI{30}{\%} sont éparpillés en petites taches isolées, ils ne forment pas un bloc fonctionnel. (2) Effet de bord : les fragments subissent microclimat altéré, invasion espèces pionnières, braconnage facilité. (3) Absence de couloirs fonctionnels : les éléphants et grands singes ont besoin de vastes territoires continus pour se nourrir, se reproduire, accéder à l'eau/minéraux (baïs). Une réserve en « îlot » dans la plantation ne suffit pas ; il faut des \textbf{couloirs écologiques} larges (> 1-2 km), connectés au parc, exempts de cultures attractives (palmiers) et de perturbations humaines.
    \item \textbf{Mesures concrètes :}
    \begin{enumerate}
        \item \textbf{Aménagement spatial rigoureux :} Regrouper les \SI{30}{\%} HCV/HCS (High Carbon Stock) en \textbf{un ou deux grands blocs contigus} le long de la limite du parc, formant une zone tampon fonctionnelle. Créer des \textbf{couloirs de migration} physiques (largeur min. 1 km) reliant ces blocs au parc, interdits à toute culture/plantation, gérés par une brigade mixte (Entreprise + MINFOF + Communautés).
        \item \textbf{Stratégie « Zéro braconnage / Zéro déforestation » opérationnelle :} Recrutement/formation/équipement d'éco-gardes privés (certifiés) sous supervision MINFOF, patrouilles quotidiennes (SMART), barrières contrôlées aux entrées pistes, sanctions internes immédiates (licenciement) pour tout personnel impliqué. Programme alternatif protéines : élevage aulacodes/volailles/poissons subventionné pour ouvriers et villages riverains.
        \item \textbf{Gouvernance partagée et partage des bénéfices :} Mise en place d'un \textbf{Comité de suivi environnemental et social (CSES)} paritaire (Entreprise, Administration, ONG locales, Représentants villages, Chefs traditionnels) se réunissant trimestriellement, avec pouvoir de suspension d'activités en cas de non-conformité. Fonds de développement local (1-2 \% CA ou redevance hectare) géré transparemment pour projets eau, santé, écoles, agroforesterie villageoise (cacao/plantain sous couvert) $\rightarrow$ appropriation locale de la conservation.
    \end{enumerate}
\end{enumerate}
\end{corrige}

\coursec{Le développement durable et la gestion raisonnée des ressources}

Après avoir examiné les pressions que les activités humaines exercent sur les forêts et la biodiversité (section précédente), nous comprenons qu'il ne suffit pas de constater la dégradation : il faut proposer des modes d'exploitation qui préservent le capital naturel pour les générations futures. C'est précisément l'objet du développement durable et de la gestion raisonnée des ressources.

\courssub{Définition et origine du concept de développement durable}

\begin{experience}[Observation d'une situation concrète : l'extraction de sable dans le Sanaga]
Dans le bassin du fleuve Sanaga, l'extraction intensive de sable pour le bâtiment a entraîné un abaissement du lit du fleuve, une érosion des berges et une perte d'habitats pour les poissons. Les riverains constatent une baisse des captures et une augmentation des inondations. On se demande : comment continuer à prélever du sable sans détruire la ressource ni nuire aux populations ?
\tcblower
\textbf{Hypothèse} : Il existe un seuil de prélèvement au‑delà duquel la ressource ne se renouvelle plus assez vite, et en deçà duquel l'exploitation peut se poursuivre indéfiniment.
\end{experience}

La communauté internationale a formalisé cette idée en 1987 dans le \textbf{Rapport Brundtland} (Commission mondiale sur l'environnement et le développement, présidée par Gro Harlem Brundtland). Ce rapport définit le développement durable comme suit :

\begin{definition}[Développement durable (Rapport Brundtland, 1987)]
Le développement durable est un développement qui répond aux besoins des générations présentes sans compromettre la capacité des générations futures à répondre aux leurs.
\end{definition}

\courssub{Les trois piliers du développement durable}

Pour qu'un projet soit « durable », il doit satisfaire simultanément trois dimensions indissociables :

\begin{itemize}
\item \textbf{Économique} : l'activité doit être viable, générer des revenus et des emplois durables.
\item \textbf{Social} : elle doit garantir l'équité, le bien‑être, la santé et la participation des populations concernées.
\item \textbf{Environnemental} : elle doit protéger les écosystèmes, la biodiversité et les ressources naturelles (eau, sols, air, forêts).
\end{itemize}

Ces trois piliers sont souvent représentés par trois cercles entrelacés : l'intersection centrale correspond aux solutions vraiment durables.

\begin{popfigure}
\begin{tikzpicture}[scale=1.2, every node/.style={font=\small}]
% Cercles avec couleurs autorisées (popblue, popgreen, poporange)
\draw[fill=popblue,opacity=0.3] (-1.2,0) circle (1.5);
\draw[fill=popgreen,opacity=0.3] (1.2,0) circle (1.5);
\draw[fill=poporange,opacity=0.3] (0,1.04) circle (1.5);
% Étiquettes
\node[popblue,font=\bfseries] at (-1.2,-1.7) {Économie};
\node[popgreen,font=\bfseries] at (1.2,-1.7) {Social};
\node[poporange,font=\bfseries] at (0,2.7) {Environnement};
% Intersection centrale
\node[font=\bfseries\small,align=center] at (0,0) {Développement\\ durable};
% Flèches d'interaction
\draw[->,>=Stealth,thick,popblue] (-2.7,0) -- (-1.5,0) node[midway,above] {investissements};
\draw[->,>=Stealth,thick,popgreen] (2.7,0) -- (1.5,0) node[midway,above] {emplois};
\draw[->,>=Stealth,thick,poporange] (0,2.54) -- (0,1.8) node[midway,right] {protection};
\end{tikzpicture}
\legende{Les trois piliers du développement durable : économie, social, environnement. Leur intersection centrale définit la zone des choix durables.}
\end{popfigure}

\begin{propriete}[Ressource renouvelable et rythme d'exploitation (admis)]
Une ressource renouvelable exploitée à un rythme inférieur à son rythme de régénération naturelle peut être utilisée indéfiniment ; si le prélèvement dépasse la régénération, le stock diminue jusqu'à l'épuisement.
\end{propriete}

\begin{methode}[Proposer une mesure de gestion durable]
Pour vérifier qu'une mesure est durable, procéder en trois étapes :
\begin{enumerate}
\item \textbf{Test économique} : la mesure est‑elle rentable ou au moins viable à long terme pour les acteurs locaux ?
\item \textbf{Test social} : améliore‑t‑elle l'équité, la santé, la sécurité alimentaire des communautés ?
\item \textbf{Test environnemental} : réduit‑elle la pression sur la ressource, favorise‑t‑elle la régénération et la biodiversité ?
\end{enumerate}
Seule une mesure validée par les trois tests peut être considérée comme durable.
\end{methode}

\courssub{Principes de la gestion raisonnée des ressources naturelles}

La gestion raisonnée traduit les trois piliers en règles opérationnelles :

\begin{itemize}
\item \textbf{Fixer des quotas de prélèvement} basés sur des études d'état des stocks (ex. : quota annuel de bois d'œuvre par hectare).
\item \textbf{Respecter les périodes de reproduction} (interdiction de pêche pendant la frayère, de chasse pendant la mise bas).
\item \textbf{Reboiser après coupe} : pour chaque arbre prélevé, planter au moins un jeune plant d'essence locale.
\item \textbf{Recycler et valoriser les déchets} : transformer les résidus de scierie en panneaux, composter les déchets agricoles, traiter les eaux usées avant rejet.
\end{itemize}

Au Cameroun, plusieurs concessions forestières (UFA -- Unités Forestières d'Aménagement) appliquent la \textbf{coupe sélective} : seuls les arbres d'un diamètre minimal (ex. 60 cm pour l'ayous \emph{Triplochiton scleroxylon}) sont abattus, les autres sont laissés pour assurer la régénération naturelle. Après l'exploitation, un plan de reboisement est obligatoire.

\begin{popfigure}
\begin{tikzpicture}[scale=1.1, every node/.style={font=\small}]
% Styles
\tikzset{
  box/.style={rectangle,rounded corners,draw=popdark,fill=popcream,thick,minimum width=2.8cm,minimum height=1.2cm,align=center},
  arrow/.style={->,>=Stealth,thick,popdark}
}
% Nœuds
\node[box, align=center] (coupe){Coupe sélective\\ (arbres matures)};
\node[box, right=3.5cm of coupe, align=center] (rebois){Reboisement\\ (essences locales)};
\node[box, below=2.5cm of rebois, align=center] (regen){Régénération naturelle\\ (croissance, faune)};
\node[box, left=3.5cm of regen, align=center] (nouvelle){Nouvelle coupe\\ (après 25‑30 ans)};
% Flèches
\draw[arrow] (coupe.east) -- (rebois.west) node[midway,above] {Plan de gestion};
\draw[arrow] (rebois.south) -- (regen.north) node[midway,right] {Suivi écologique};
\draw[arrow] (regen.west) -- (nouvelle.east) node[midway,below] {Inventaire};
\draw[arrow] (nouvelle.north) -- (coupe.south) node[midway,left] {Cycle renouvelé};
% Légende centrale
\node[font=\bfseries\small,align=center] at (0,0) {Cycle de gestion\\ raisonnée d'une forêt};
\end{tikzpicture}
\legende{Cycle de gestion raisonnée d'une forêt tropicale : coupe sélective → reboisement → régénération → nouvelle coupe. Chaque étape est encadrée par un plan de gestion approuvé par l'État.}
\end{popfigure}

\begin{exemplebox}[Certification FSC au Cameroun]
La certification \textbf{FSC} (Forest Stewardship Council) garantit qu'un bois provient d'une forêt gérée selon des critères stricts : respect des quotas, protection des zones à haute valeur de conservation, droits des travailleurs et des communautés locales. Plusieurs UFA dans la région de l'Est ont obtenu cette certification, ce qui leur ouvre l'accès aux marchés européens exigeants et valorise une exploitation responsable.
\end{exemplebox}

\courssub{L'empreinte écologique comme indicateur de pression}

L'\textbf{empreinte écologique} mesure la surface de territoire productif (forêts, terres agricoles, zones de pêche, surfaces nécessaires pour absorber le CO₂) requise pour soutenir le mode de vie d'une population. Elle s'exprime en hectares globaux par personne (gha/hab). Si l'empreinte dépasse la biocapacité disponible, le territoire est en situation de \textbf{déficit écologique}. Au Probatoire, la définition est admise ; le calcul n'est pas exigible.

\begin{aretenir}[Empreinte écologique]
Indicateur synthétique de la pression humaine sur la biosphère. Une empreinte supérieure à la biocapacité signifie que l'humanité puise dans le capital naturel (déforestation, surpêche, émissions de CO₂ non absorbées).
\end{aretenir}

\courssub{Responsabilités partagées : État, collectivités, entreprises, citoyens}

La durabilité ne repose pas sur un seul acteur :

\begin{center}
\begin{tabularx}{\linewidth}{|l|X|}\hline
\textbf{Acteur} & \textbf{Rôle principal} \\ \hline
État & Légiférer (code forestier, loi sur l'eau), créer des aires protégées, contrôler les quotas, délivrer les certificats. \\ \hline
Collectivités locales & Gérer les ressources de proximité (forêts communales, points d'eau), sensibiliser, organiser la collecte sélective des déchets. \\ \hline
Entreprises & Adopter des chartes RSE (Responsabilité Sociétale des Entreprises), obtenir des certifications (FSC, ISO 14001), investir dans l'économie circulaire. \\ \hline
Citoyens & Réduire la consommation d'eau, trier les déchets, privilégier les produits locaux certifiés, participer aux reboisements communautaires. \\ \hline
\end{tabularx}
\end{center}

\begin{savaistu}[Histoire du concept]
Le terme « développement durable » a été popularisé en 1987 par la Commission mondiale sur l'environnement et le développement, présidée par \textbf{Gro Harlem Brundtland} (ancienne Première ministre de Norvège). Le rapport \emph{Notre avenir à tous} (Brundtland Report) a posé les bases des conférences de Rio (1992), Johannesburg (2002) et des Objectifs de développement durable (ODD) de l'ONU en 2015.
\end{savaistu}

\courssub{Éducation à l'environnement : levier de changement des comportements}

L'éducation à l'environnement (ÉE) vise à faire comprendre les interdépendances entre activités humaines et écosystèmes, pour transformer les savoirs en \textbf{savoir‑être} (respect, responsabilité) et en \textbf{savoir‑faire} (tri, économie d'eau, agroforesterie). Au Cameroun, les clubs « Jeunes pour l'environnement » dans les lycées, les campagnes « Un arbre, un élève » et les modules d'ÉE intégrés au programme de SVT sont des exemples concrets. L'ÉE est le levier le plus puissant pour ancrer la gestion raisonnée dans la culture quotidienne.

\courssub{Synthèse : analyser une situation et proposer des solutions durables}

\begin{exoresolu}[Exploitation de sable dans le fleuve Sanaga]
\textbf{Énoncé} : Une entreprise prélève 150 000 m³ de sable par an dans le lit mineur du Sanaga, près de la ville d'Édéa. Les riverains signalent une baisse de la pêche, une érosion des berges et une turbidité accrue de l'eau. L'entreprise emploie 80 personnes et fournit du sable pour la construction de routes. Proposer un plan de gestion durable.

\tcblower
\textbf{Solution détaillée} :

\begin{enumerate}
\item \textbf{Diagnostic (observation + données)} : Mesurer le débit solide charrié estimé sur le tronçon concerné (≈ 200 000 m³/an, valeur indicative) et comparer au prélèvement (150 000 m³/an = 75 \% du flux). Le prélèvement dépasse le seuil de renouvellement naturel (souvent fixé à 30‑40 \% du flux sédimentaire charrié).
\item \textbf{Analyse selon les trois piliers} :
\begin{itemize}
\item \textit{Économique} : L'entreprise génère des revenus et des emplois → maintenir l'activité mais réduire le volume.
\item \textit{Social} : Les pêcheurs perdent leurs revenus, les berges s'effondrent → protéger les communautés.
\item \textit{Environnemental} : Érosion, perte d'habitats, turbidité → laisser le lit se reconstituer.
\end{itemize}
\item \textbf{Mesures proposées} (validées par la méthode des trois tests) :
\begin{itemize}
\item Fixer un \textbf{quota annuel} à 60 000 m³ (≈ 30 \% du flux sédimentaire charrié), basé sur une étude hydrologique.
\item Instaurer une \textbf{période de repos} de 3 mois par an (saison des crues) pour permettre le rechargement naturel.
\item Obliger le \textbf{reboisement des berges} (espèces ripicoles : \emph{Mitragyna inermis}, \emph{Alchornea cordifolia}) sur 5 km linéaires.
\item Créer un \textbf{fonds de compensation} alimenté par l'entreprise pour soutenir les pêcheurs (formation à l'aquaculture, matériel).
\item Mettre en place un \textbf{suivi participatif} : comité local (administration, entreprise, pêcheurs, ONG) qui contrôle les volumes prélevés et l'état des berges.
\end{itemize}
\item \textbf{Conclusion} : Ce plan respecte les trois piliers : l'exploitation reste viable (quota réduit mais stable), les populations sont protégées et compensées, l'écosystème fluvial peut se régénérer. C'est un exemple de gestion raisonnée applicable à d'autres ressources (bois, eau, sols).
\end{enumerate}
\end{exoresolu}

\begin{attention}[Pièges fréquents à éviter]
\begin{itemize}
\item Confondre « développement durable » et « protection de l'environnement » : le développement durable inclut obligatoirement les dimensions économique et sociale.
\item Oublier le test social lors de la proposition d'une mesure (ex. : interdire toute coupe sans alternative de revenus pour les populations).
\item Croire qu'une ressource renouvelable est inépuisable : sans respect du rythme de régénération, elle s'épuise (cf. propriété admise).
\item Négliger l'échelle locale : une mesure durable au niveau national peut être inappropriée pour une communauté donnée.
\end{itemize}
\end{attention}

\coursec{Exemples camerounais d'exploitation et de conservation}

Après avoir défini les principes du \cle{développement durable} et de la \cle{gestion raisonnée}, il est indispensable d'ancrer ces concepts dans la réalité du terrain camerounais. Le Cameroun, souvent décrit comme une « Afrique en miniature » du fait de sa diversité écologique (zone sahélienne au nord, forêt dense humide au sud, montagnes à l'ouest, littoral atlantique), concentre l'essentiel des enjeux de conservation du continent. Observons comment le pays tente de concilier exploitation économique de ses richesses naturelles et préservation de son capital écologique pour les générations futures.

\courssub{Les aires protégées : piliers de la conservation \emph{in situ}}

\begin{savaistu}[L'origine des aires protégées au Cameroun]
La première réserve de faune, celle du Dja, a été créée en 1950 sous l'administration coloniale. Après l'indépendance, la loi n°81-13 du 27 novembre 1981 relative à la forêt, à la faune et à la pêche pose le cadre juridique moderne. Elle définit les catégories d'aires protégées et institue la notion de domaine forestier permanent (forêts destinées à rester forestières) et de domaine forestier non permanent. Aujourd'hui, le réseau d'aires protégées couvre environ 19 \% du territoire national, soit plus de 9 millions d'hectares, avec l'objectif international d'atteindre 30 \% (objectif d'Aichi, puis Cadre mondial de Kunming-Montréal).
\end{savaistu}

\begin{definition}[Aire protégée]
Un espace géographique délimité, reconnu, dédié et géré, par des moyens légaux ou effectifs, en vue d'assurer la conservation à long terme de la nature, des services écosystémiques et des valeurs culturelles associées. Au Cameroun, les principales catégories sont les \textbf{parcs nationaux} (protection intégrale, interdiction de toute chasse, exploitation forestière, agriculture), les \textbf{réserves de faune} (protection de la faune, droits d'usage limités pour les populations riveraines) et les \textbf{réserves de biosphère} (zones cœur, tampon, transition ; réconciliation conservation/développement).
\end{definition}

\begin{popfigure}
\begin{tikzpicture}[scale=0.75, every node/.style={transform shape}]
    % Carte simplifiée du Cameroun (contour approximatif)
    \draw[popdark, thick, fill=popcream] 
        (0,0) -- (1,2) -- (3,4) -- (5,5) -- (7,4.5) -- (9,4) -- (10,2) -- (9.5,0) -- (8,-1) -- (6,-1.5) -- (4,-1) -- (2,-0.5) -- cycle;
    
    % Lac Tchad (extrême nord)
    \draw[popblue, fill=popblue!20] (1.5, -0.5) ellipse (0.8 and 0.4);
    \node[popblue, font=\tiny] at (1.5, -0.5) {Lac Tchad};

    % Fleuves
    \draw[popblue, thick] (4, -1) .. controls (4.5, 1) and (5, 3) .. (5.5, 5); % Bénoué -> Sanaga
    \draw[popblue, thick] (6, 4.5) .. controls (7, 3) .. (7.5, 1); % Sanaga aval
    \node[popblue, font=\tiny, rotate=30] at (4.5, 1.5) {Bénoué};
    \node[popblue, font=\tiny, rotate=-20] at (6.5, 3) {Sanaga};
    % Nyong : coule vers l'Ouest (x décroissant)
    \draw[popblue, thick] (4, 4) .. controls (3, 3) .. (2, 2); 
    \node[popblue, font=\tiny, rotate=-20] at (3, 3) {Nyong};

    % Mont Cameroun
    \node[poporange, font=\small\bfseries, fill=popcream, rounded corners, inner sep=1pt] at (9.5, 2.5) {Mont Cameroun (4095 m)};
    \draw[poporange, thick] (9.5, 2.5) circle (0.3);

    % Villes repères
    \node[popdark, fill=popcream, circle, inner sep=1.5pt, font=\tiny] (Ya) at (5.5, 3.2) {};
    \node[below, font=\tiny] at (Ya) {Yaoundé};
    \node[popdark, fill=popcream, circle, inner sep=1.5pt, font=\tiny] (Do) at (8.5, 1.5) {};
    \node[below, font=\tiny] at (Do) {Douala};
    \node[popdark, fill=popcream, circle, inner sep=1.5pt, font=\tiny] (Ga) at (3.5, 0.5) {};
    \node[below, font=\tiny] at (Ga) {Garoua};
    \node[popdark, fill=popcream, circle, inner sep=1.5pt, font=\tiny] (Ma) at (2.5, -0.8) {};
    \node[below, font=\tiny] at (Ma) {Maroua};

    % Aires protégées (Parcs Nationaux = Carré vert, Réserves = Losange bleu)
    % 1. Waza (Nord)
    \node[popgreen, fill=popgreen!20, rectangle, rounded corners, inner sep=2pt, font=\tiny\bfseries] (Waza) at (2.5, 0.2) {Waza};
    \node[below, font=\tiny, popmuted] at (Waza.south) {PN};

    % 2. Bénoué (Nord)
    \node[popgreen, fill=popgreen!20, rectangle, rounded corners, inner sep=2pt, font=\tiny\bfseries] (Ben) at (4.2, 1.0) {Bénoué};
    \node[below, font=\tiny, popmuted] at (Ben.south) {PN};

    % 3. Faro (Nord-Ouest)
    \node[popgreen, fill=popgreen!20, rectangle, rounded corners, inner sep=2pt, font=\tiny\bfseries] (Faro) at (3.5, 2.0) {Faro};
    \node[below, font=\tiny, popmuted] at (Faro.south) {PN};

    % 4. Bouba Ndjida (Nord)
    \node[popgreen, fill=popgreen!20, rectangle, rounded corners, inner sep=2pt, font=\tiny\bfseries] (Bouba) at (3.0, 1.5) {Bouba Ndjida};
    \node[below, font=\tiny, popmuted] at (Bouba.south) {PN};

    % 5. Mbam et Djerem (Centre-Est)
    \node[popgreen, fill=popgreen!20, rectangle, rounded corners, inner sep=2pt, font=\tiny\bfseries] (Mbam) at (5.5, 2.5) {Mbam-Djerem};
    \node[below, font=\tiny, popmuted] at (Mbam.south) {PN};

    % 6. Deng Deng (Est)
    \node[popgreen, fill=popgreen!20, rectangle, rounded corners, inner sep=2pt, font=\tiny\bfseries] (Deng) at (6.5, 3.0) {Deng Deng};
    \node[below, font=\tiny, popmuted] at (Deng.south) {PN};

    % 7. Lobéké (Sud-Est)
    \node[popgreen, fill=popgreen!20, rectangle, rounded corners, inner sep=2pt, font=\tiny\bfseries] (Lob) at (7.0, 3.8) {Lobéké};
    \node[below, font=\tiny, popmuted] at (Lob.south) {PN};

    % 8. Boumba Bek (Sud-Est)
    \node[popgreen, fill=popgreen!20, rectangle, rounded corners, inner sep=2pt, font=\tiny\bfseries] (Boumba) at (6.0, 4.0) {Boumba Bek};
    \node[below, font=\tiny, popmuted] at (Boumba.south) {PN};

    % 9. Nki (Sud-Est)
    \node[popgreen, fill=popgreen!20, rectangle, rounded corners, inner sep=2pt, font=\tiny\bfseries] (Nki) at (5.5, 4.2) {Nki};
    \node[below, font=\tiny, popmuted] at (Nki.south) {PN};

    % 10. Campo-Ma'an (Sud)
    \node[popgreen, fill=popgreen!20, rectangle, rounded corners, inner sep=2pt, font=\tiny\bfseries] (Campo) at (7.5, 2.0) {Campo-Ma'an};
    \node[below, font=\tiny, popmuted] at (Campo.south) {PN};

    % 11. Korup (Sud-Ouest)
    \node[popgreen, fill=popgreen!20, rectangle, rounded corners, inner sep=2pt, font=\tiny\bfseries] (Korup) at (8.5, 2.8) {Korup};
    \node[below, font=\tiny, popmuted] at (Korup.south) {PN};

    % 12. Mont Cameroun (Sud-Ouest) - Parc National récent
    \node[popgreen, fill=popgreen!20, rectangle, rounded corners, inner sep=2pt, font=\tiny\bfseries] (MtnCam) at (9.2, 2.0) {Mont Cameroun};
    \node[below, font=\tiny, popmuted] at (MtnCam.south) {PN};

    % Réserves de Faune / Biosphère (Losanges)
    % Dja (Sud-Est) - UNESCO
    \node[popblue, fill=popblue!15, diamond, inner sep=3pt, font=\tiny\bfseries] (Dja) at (6.0, 3.2) {Dja};
    \node[above, font=\tiny, popmuted] at (Dja.north) {RF/UNESCO};

    % Réserve de faune du Dja (contour approximatif)
    \draw[popblue, dashed, thick] (5.2, 2.8) ellipse (1.0 and 0.6);

    % Santchou (Ouest)
    \node[popblue, fill=popblue!15, diamond, inner sep=3pt, font=\tiny\bfseries] (Sant) at (7.0, 3.0) {Santchou};
    \node[above, font=\tiny, popmuted] at (Sant.north) {RF};

    % Légende
    \begin{scope}[shift={(10.5, 4.5)}]
        \node[popgreen, fill=popgreen!20, rectangle, rounded corners, inner sep=2pt, font=\tiny\bfseries] (leg1) at (0,0) {PN};
        \node[right, font=\small] at (leg1.east) {Parc National};
        \node[popblue, fill=popblue!15, diamond, inner sep=3pt, font=\tiny\bfseries] (leg2) at (0,-0.6) {RF};
        \node[right, font=\small] at (leg2.east) {Réserve Faune / Biosphère};
    \end{scope}

    % Échelle
    \draw[popdark, thick] (0.5, -1.8) -- (2.5, -1.8);
    \node[below, font=\tiny] at (1.5, -1.9) {~200 km};
\end{tikzpicture}
\legende{Carte simplifiée du Cameroun localisant les principaux parcs nationaux (PN, vert) et réserves de faune/biosphère (RF, bleu). Le réseau couvre les trois grandes zones écologiques : soudano-sahélienne (Waza, Bénoué), guinéo-congolaise (Korup, Campo-Ma'an, Dja, Lobéké) et montagnarde (Mont Cameroun).}
\end{popfigure}

\begin{exemplebox}[La Réserve de Faune du Dja : joyau de la biodiversité et patrimoine mondial]
La réserve du Dja (région du Sud, Est) illustre parfaitement l'intérêt des aires protégées. Classée \textbf{Réserve de Biosphère} en 1981 et inscrite au \textbf{Patrimoine Mondial de l'UNESCO} en 1987 (critères ix et x), elle s'étend sur \textbf{526 000 hectares}. Elle abrite une forêt dense humide de type guinéo-congolais quasi intacte.
\begin{itemize}
    \item \textbf{Biodiversité :} Plus de 1 500 espèces de plantes, 107 mammifères (dont gorilles de plaine, chimpanzés, éléphants de forêt, bongos, mandrills), 320 oiseaux.
    \item \textbf{Populations locales :} Les Baka (pygmées) y pratiquent chasse, cueillette, pêche traditionnelles (droits d'usage reconnus en zone périphérique).
    \item \textbf{Gestion :} Plan de gestion participatif (MINFOF, ONG comme WWF, ECOFAC, populations). Surveillance par écogardes, bio-monitoring (pièges photographiques, transects).
    \item \textbf{Menaces persistantes :} Braconnage commercial (ivoire, viande de brousse), orpaillage illégal, pression agricole en périphérie.
\end{itemize}
\end{exemplebox}

\begin{attention}[Aire protégée « papier » vs aire protégée effective]
La création légale (décret) ne suffit pas. Une aire protégée sans \textbf{moyens de surveillance} (écogardes formés, équipés, payés), sans \textbf{plan de gestion} actualisé et sans \textbf{implication réelle des populations riveraines} (partage des retombées, concertation) devient une « aire protégée de papier » où le braconnage et l'exploitation illégale (bois, or, agriculture) prospèrent. C'est le défi majeur de nombreux parcs du Septentrion (Waza, Bénoué) confrontés à l'insécurité transfrontalière et au manque de personnel.
\end{attention}

\courssub{Forêts communautaires et communales : la gestion participative en action}

La loi forestière de 1994 (loi n°94/01 du 20 janvier 1994) a introduit une innovation majeure : la possibilité pour les communautés villageoises de devenir gestionnaires légaux d'une portion de forêt du domaine permanent. C'est l'application concrète du principe de \textbf{gouvernance partagée} (subsidiarité).

\begin{definition}[Forêt communautaire]
Portion du domaine forestier permanent, d'une superficie maximale de 5 000 hectares, attribuée par l'État à une communauté villageoise (groupement d'habitants d'un ou plusieurs villages) pour une durée de 25 ans renouvelable. La communauté en assure la gestion durable selon un \textbf{plan d'aménagement simple} validé par l'administration (MINFOF), en tire des revenus (vente de bois, produits forestiers non ligneux - PFNL, écotourisme) et a l'obligation de la conserver.
\end{definition}

\begin{popfigure}
\begin{tikzpicture}[node distance=1.2cm and 1.5cm, >=Stealth, font=\small]
    % Styles
    \tikzset{
        actor/.style={ellipse, draw=popdark, thick, fill=popcream, minimum width=2.5cm, minimum height=1.2cm, align=center},
        process/.style={rectangle, rounded corners, draw=popblue, thick, fill=popblue!10, minimum width=3cm, minimum height=1.2cm, align=center},
        benefit/.style={diamond, draw=popgreen, thick, fill=popgreen!10, aspect=2, minimum width=2.5cm, align=center},
        arrow/.style={->, thick, popdark},
        label/.style={font=\tiny, popmuted, midway, fill=popcream, inner sep=1pt}
    }

    % Nœuds
    \node[actor, align=center] (pop){Population locale\\ (Communauté villageoise)};
    \node[process, right=of pop, align=center] (gestion){Gestion durable\\ (Plan d'aménagement simple,\\ Inventaires, Coupes sélectives,\\ Suivi biodiversité)};
    \node[benefit, below=of gestion, align=center] (revenus){Revenus\\ (Vente bois certifié,\\ PFNL, Écotourisme,\\ Reversement taxes)};
    \node[actor, right=of gestion, align=center] (etat){Administration\\ (MINFOF, Commune,\\ Services déconcentrés)};
    \node[benefit, above=of gestion, align=center] (protection){Protection de la forêt\\ (Stock carbone, Biodiversité,\\ Services écosystémiques,\\ Lutte érosion)};

    % Flèches principales
    \draw[arrow] (pop) -- node[label, above, align=center]{Demande attribution,\\ Mise en œuvre} (gestion);
    \draw[arrow] (gestion) -- node[label, above] {Rapports, Reddition comptes} (etat);
    \draw[arrow] (etat) -- node[label, above, align=center]{Appui technique,\\ Contrôle, Sécurisation foncière} (gestion);
    \draw[arrow] (gestion) -- node[label, right] {Génère} (revenus);
    \draw[arrow] (revenus) -- node[label, left, align=center]{Réinvestit (écoles, santé,\\ micro-projets)} (pop);
    \draw[arrow] (gestion) -- node[label, left] {Assure} (protection);
    \draw[arrow] (protection) -- node[label, left] {Bénéficie à} (pop);
    
    % Boucle de rétroaction
    \draw[arrow, poporange, dashed] (revenus.east) .. controls +(1,0) and +(1,0) .. (etat.south) node[label, right] {Taxes / Redevances};
    \draw[arrow, poporange, dashed] (protection.west) .. controls +(-1,0) and +(-1,0) .. (pop.north) node[label, left] {Ressources durables};

    % Cadre légal
    \node[below=0.5cm of pop, font=\tiny\bfseries, poppurple] {Cadre légal : Loi 1994, Décret 95/531/PM, Manuel de procédures};
\end{tikzpicture}
\legende{Fonctionnement d'une forêt communautaire : boucle vertueuse entre population gestionnaire, administration tutélaire, génération de revenus locaux et conservation de l'écosystème forestier.}
\end{popfigure}

\begin{experience}[Analyse d'un document : Convention de forêt communautaire]
\textbf{Contexte :} Un groupement villageois de la région de l'Est signe une convention avec l'État pour une forêt communautaire de 3 000 ha.
\textbf{Document fourni :} Extrait du plan d'aménagement simple (inventaire : 15 m³/ha exploitable ; rotation 25 ans ; assemblage 1 : essence principale \emph{ayous}, 2 : essences secondaires).
\textbf{Travail demandé :}
\begin{enumerate}
    \item Identifier la ressource et l'activité.
    \item Calculer le volume annuel de coupe (possibilité annuelle).
    \item Expliquer pourquoi la rotation de 25 ans est une mesure de gestion durable.
    \item Citer deux avantages socio-économiques pour le village.
\end{enumerate}
\begin{center}
\textbf{Solution guidée}
\end{center}
\begin{enumerate}
    \item \textbf{Ressource :} Forêt dense humide (bois d'œuvre, PFNL). \textbf{Activité :} Exploitation forestière artisanale / communautaire.
    \item \textbf{Possibilité Annuelle de Coupe (PAC) =} Superficie / Rotation = 3 000 ha / 25 ans = 120 ha/an. Volume = 120 ha $\times$ 15 m³/ha = \textbf{1 800 m³/an}. C'est le volume maximal prélevable sans épuiser le capital.
    \item La rotation (25 ans) correspond au temps nécessaire pour que les \textbf{arbres rémanents} et la régénération naturelle reconstituent le stock exploitable. Elle garantit la \textbf{pérennité de la ressource} (principe de rendement soutenu).
    \item Revenus financiers (vente grumes/bois transformé) $\rightarrow$ fonds de développement communautaire (école, forage, santé). Emplois locaux (abattage, débardage, sciage). Valorisation PFNL (gnetum, Irvingia, champignons).
\end{enumerate}
\end{experience}

\begin{aretenir}[Forêt communale vs Forêt communautaire]
\begin{itemize}
    \item \textbf{Forêt communautaire :} Gérée par une \emph{communauté villageoise} (groupement d'habitants), initiative « par le bas », objectif prioritaire : amélioration conditions de vie locales + conservation.
    \item \textbf{Forêt communale :} Gérée par la \emph{commune} (collectivité territoriale décentralisée, maire/conseil municipal), initiative institutionnelle, objectif : développement local communal (infrastructures, budget) + conservation. Superficie non plafonnée à 5 000 ha.
\end{itemize}
Les deux modèles visent à détourner les populations de l'exploitation illicite en leur donnant un \textbf{intérêt économique direct} à la survie de la forêt.
\end{aretenir}

\courssub{Aménagement forestier durable dans les concessions industrielles}

Le domaine forestier permanent est en grande partie alloué sous forme de \textbf{Unités de Gestion Forestière (UGF)} à des entreprises privées (concessions). La loi de 1994 impose l'élaboration d'un \textbf{Plan d'Aménagement (PA)} sur 30 ans, révisable tous les 5 ans, et un \textbf{Cahier des Charges} strict.

\begin{propriete}[Principes de l'aménagement forestier durable (AFD) en concession]
\begin{enumerate}
    \item \textbf{Découpage en Assises Annuelles de Coupe (AAC) :} La concession est divisée en 30 blocs (un par an). On n'exploite qu'une AAC par an $\rightarrow$ 30 ans de repos pour la forêt avant retour.
    \item \textbf{Inventaire pré-exploitation (100\% ou échantillonnage) :} Localisation, identification, mesure (\cle{DHP} $\ge$ 60 cm généralement) de chaque arbre exploitable. Cartographie GPS.
    \item \textbf{Choix des arbres :} Seuls les arbres \textbf{mûrs} (\cle{DHP} $\ge$ \cle{DME} - Diamètre Minimum d'Exploitation, variable selon essence) et \textbf{sains} sont marqués pour la coupe (marqueau au marteau, peinture). \textbf{Arbres semenciers} (gros, bien formés) et \textbf{arbres d'avenir} (jeunes, vigoureux) sont \textbf{protégés} (marque de conservation).
    \item \textbf{Coupe sélective directionalisée :} Abattage dirigé pour minimiser les dégâts sur la rémanence (arbres restants) et la régénération. Débardage sur pistes planifiées (skidders) $\rightarrow$ limitation de la compaction des sols.
    \item \textbf{Suivi post-exploitation :} Contrôle de la régénération naturelle, enrichissement si nécessaire (plantations d'essences locales à croissance rapide : \emph{Terminalia superba} - Fraké, \emph{Triplochiton scleroxylon} - Ayous).
    \item \textbf{Traçabilité et certification :} Système informatisé (SIGIF - Système d'Information de Gestion de l'Information Forestière) du bois de la souche au port (Douala/Kribi). Certification FSC (Forest Stewardship Council) ou PAFC (Pan African Forest Certification) comme preuve de marché.
\end{enumerate}
\end{propriete}

\begin{savaistu}[ONADEF / ANAFOR : l'effort national de reboisement]
L'\textbf{Office National de Développement des Forêts (ONADEF)}, devenu \textbf{ANAFOR} (Agence Nationale d'Appui au Développement Forestier) en 2016, pilote les plantations forestières. Depuis les années 70, plus de \textbf{30 000 ha} de plantations industrielles (principalement \emph{Pinus caribaea}, \emph{Eucalyptus} spp. dans l'Adamaoua et l'Ouest ; \emph{Terminalia}, \emph{Triplochiton} au Sud) ont été établis. Objectif : réduire la pression sur la forêt naturelle, fournir du bois d'œuvre, pâte à papier, bois-énergie, et séquestrer du carbone. Défis : entretien (feu, broussailles), rentabilité économique, acceptation sociale (accaparement terres).
\end{savaistu}

\courssub{Lutte contre la désertification : l'héritage de l'Opération Sahel Vert}

Dans les régions du Nord (Nord, Extrême-Nord, Adamaoua nord), l'avancée du désert (sable éolien, dégradation sols, perte couverture végétale) menace l'agriculture, l'élevage et l'habitat. La réponse historique est l'\textbf{Opération Sahel Vert} (lancée 1976, relancée 2008 dans le cadre de la \cle{Grande Muraille Verte} panafricaine).

\begin{methode}[Techniques de conservation des sols et de l'eau (CES/DRS) au Sahel camerounais]
\textbf{Objectif :} Ralentir le ruissellement, augmenter l'infiltration, retenir les sédiments, restaurer la fertilité.
\begin{enumerate}
    \item \textbf{Diagnostic terrain :} Identifier les zones de ruissellement concentré, pentes, type de sol (sableux, latéritique), cultures.
    \item \textbf{Choix de la technique adaptée :}
        \begin{itemize}
            \item \textbf{Cordons pierreux (digues filtrantes) :} Rangées de pierres le long des courbes de niveau. Freinent l'eau, laissent passer le fin, retiennent le gros. Faciles, matériaux locaux, main-d'œuvre villageoise (\cle{HIMO} - Haute Intensité de Main d'Œuvre).
            \item \textbf{Demi-lunes / Zai :} Fosses demi-circulaires (demi-lunes) ou trous (Zai) en amont des plants. Concentrent eau + matière organique (fumier) au pied de la plante. Très efficaces pour réhabilitation sols crustés (glacis).
            \item \textbf{Agroforesterie / Parc à \emph{Faidherbia albida} (Kad) :} Arbre légumineux à phénologie inversée (feuilles en saison sèche, nu en saison des pluies). Fixe l'azote, fertilise (litière), ombre le bétail, ne concurrence pas les cultures (mil, sorgho, maïs). \textbf{Espèce clé du Sahel}.
            \item \textbf{Bandes enherbées / Haies vives :} \emph{Andropogon gayanus}, \emph{Vetiveria zizanioides} (vetiver) en courbes de niveau.
        \end{itemize}
    \item \textbf{Mise en œuvre participative :} Groupements villageois, ONG (IRAD, CIRAD, GIZ, FAO), services déconcentrés (MINEPDED, MINADER, MINEPIA). Formation, outillage, suivi.
    \item \textbf{Évaluation :} Taux de survie plants, rendements cultures, niveau nappe phréatique, couverture végétale (NDVI satellite).
\end{enumerate}
\end{methode}

\begin{exemplebox}[Résultats observés dans le Mayo-Danay (Extrême-Nord)]
Sur le site de \textbf{Kalmalo}, la mise en place de cordons pierreux + demi-lunes + plantation de \emph{Faidherbia albida} et \emph{Acacia senegal} (gommier) sur 500 ha a permis :
\begin{itemize}
    \item Augmentation rendement mil/sorgho de \textbf{+ 40 à 60 \%} (passage de 400-500 kg/ha à 800-1000 kg/ha).
    \item Remontée nappe phréatique (puits moins profonds).
    \item Revenu complémentaire : gomme arabique (\emph{Acacia senegal}), fourrage (feuilles \emph{Faidherbia}), bois de chauffe.
    \item Fixation des dunes mobiles menaçant le village.
\end{itemize}
\textbf{Limite :} La maintenance annuelle des ouvrages (réfection cordons après grosses pluies) nécessite une organisation villageoise pérenne, souvent fragile sans appui externe continu.
\end{exemplebox}

\courssub{Gestion de l'eau : de la ressource au service}

L'eau douce est une ressource critique, inégalement répartie (abondante au Sud, rare et saisonnière au Nord) et menacée par la pollution (urbaine, industrielle, agricole) et l'érosion des bassins versants.

\begin{propriete}[Mesures clés de gestion durable de l'eau au Cameroun]
\begin{itemize}
    \item \textbf{Protection des bassins versants :} Reboisement des têtes de bassin (ex. bassin de la Sanaga, du Wouri, de la Bénoué), lutte contre l'érosion (CES), régulation des cultures sur berges (bandes tampons de 10-15 m minimum, Code de l'Eau 1998).
    \item \textbf{Traitement des eaux usées :} Stations d'épuration (Yaoundé : station de Mfoundi, Douala : station de Bassa - capacités encore insuffisantes), lagunage naturel (stabilisation par bassins) pour villes moyennes. \textbf{Enjeu :} Réduction charge organique (DBO5, DCO), azote, phosphore, pathogènes avant rejet en milieu récepteur.
    \item \textbf{Adduction d'eau potable (AEP) :} Barrages (Lom Pangar, Memve'ele, Nachtigal - hydroélectricité + réserve eau), forages équipés pompes solaires/humaines (zones rurales), réseaux urbains (CAMWATER). Objectif ODD 6 : accès universel 2030.
    \item \textbf{Gestion intégrée des ressources en eau (GIRE) :} Comités de bassin (Comité de bassin de la Sanaga, du Lac Tchad, du Wouri) réunissant usagers (agriculteurs, éleveurs, industriels, CAMWATER, ONG, administration) pour planifier le partage de l'eau (irrigation, eau potable, énergie, écosystèmes).
\end{itemize}
\end{propriete}

\begin{attention}[Conflits d'usage de l'eau : le cas du bassin de la Bénoué / Lac Tchad]
Le bassin de la Bénoué (affluent majeur du Niger, via le Lac Tchad) concentre les tensions : \textbf{barrages hydroélectriques} (Lagdo, projet de Garoua) modifient le régime de crue $\rightarrow$ impact sur \textbf{récession agricole} (riz, sorgho) et \textbf{pêche} en aval (plaine d'inondation de Waza-Logone) $\rightarrow$ conflit avec \textbf{élevage transhumant} (point d'eau, pâturage) et \textbf{conservation} (parc de Waza, zone RAMSAR). La \cle{GIRE} est l'unique cadre de résolution pacifique, mais sa mise en œuvre effective reste lente.
\end{attention}

\courssub{L'écotourisme : valorisation économique de la conservation}

L'\cle{écotourisme} (tourisme responsable dans les zones naturelles, conservant l'environnement et améliorant le bien-être des populations locales) est promu comme alternative à l'exploitation extractive (bois, minerais, braconnage).

\begin{exemplebox}[Deux modèles d'écotourisme au Cameroun]
\begin{enumerate}
    \item \textbf{Parc National de Campo-Ma'an (Sud) :} Programme \textbf{« Gorilles de Campo »} (WWF, MINFOF, communauté). Habituation de groupes de gorilles de plaine de l'Ouest (\emph{Gorilla gorilla gorilla}) à la présence humaine. Touristes (nationaux/expats) paient permis (\textasciitilde 50 000 - 100 000 FCFA) + guide + hébergement communautaire. \textbf{Retombées :} Salaires guides/pisteurs (jeunes locaux), fonds de développement villageois (écoles, santé), sensibilisation $\rightarrow$ réduction braconnage local. \textbf{Défi :} Accessibilité (pistes dégradées), insécurité régionale, faible notoriété internationale.
    \item \textbf{Parc National de Waza (Extrême-Nord) :} Classique « safari » (lions, éléphants, kob, girafes, avifaune exceptionnelle zone humide Waza-Logone). Infrastructure (campements) vieillissante. \textbf{Potentiel :} Couplage avec tourisme culturel (royaumes Kotoko, Mousgoum), ornithologie (zone RAMSAR). \textbf{Freins majeurs :} Insécurité (Boko Haram zone frontière), baisse effectifs faune (braconnage, sècheresse), manque de promotion.
\end{enumerate}
\end{exemplebox}

\begin{aretenir}[Conditions de succès de l'écotourisme communautaire]
\begin{itemize}
    \item \textbf{Droits clairs :} Contrat de concession touristique ou accord de cogestion signé (MINFOF, Commune, Communauté).
    \item \textbf{Partage équitable des revenus :} Mécanisme transparent (ex. 40\% communauté, 30\% gestion parc, 30\% Etat/Commune).
    \item \textbf{Capacités locales :} Formation guides, hôtellerie, gestion comptable, anglais.
    \item \textbf{Produit attractif :} Faune visible, paysage unique, accessibilité raisonnable, sécurité.
    \item \textbf{Marketing :} Partenariat agences de voyage, plateformes numériques, foires (FITUR, ITB, Salon du Tourisme Yaoundé).
\end{itemize}
Sans ces conditions, l'écotourisme reste marginal et ne génère pas l'alternative économique espérée face au braconnage ou à l'agriculture itinérante sur brûlis.
\end{aretenir}

\courssub{Limites structurelles et défis persistants}

Malgré un arsenal juridique complet (Loi 94, Loi 96 sur l'environnement, Code de l'Eau, Plan National de Développement Durable, engagements internationaux : CBD, CCD, UNFCCC, RAMSAR, CITES), l'efficacité de la gestion durable bute sur des obstacles systémiques.

\begin{attention}[Les 5 freins majeurs à la gestion durable des ressources au Cameroun]
\begin{enumerate}
    \item \textbf{Braconnage et trafic d'espèces (faune/flore) :} Réseaux criminels organisés (ivoire, écailles de pangolin, bois de rose \emph{Pterocarpus erinaceus}, viande de brousse urbaine). Moyens de surveillance (écogardes) \textbf{insuffisants} (ratio 1 écogarde / 100-200 km² vs norme UICN 1/10-20 km²), armement inadapté, corruption.
    \item \textbf{Conflits d'usage et fonciers :} Superposition titres (concession forestière $\cap$ permis minier $\cap$ agriculture $\cap$ aire protégée $\cap$ terres coutumières). Absence de cadastre rural abouti. Conflits agriculteurs/éleveurs (couloirs de transhumance bouchés par cultures).
    \item \textbf{Gouvernance et reddition de comptes :} Traçabilité bois perfectible (bois « gris » / illégal dans filière formelle). Gestion des redevances forestières (RFA - Redevance Forestière Annuelle, taxes d'abattage) : redistribution aux communes et communautés souvent retardée, détournée ou mal utilisée (manque de capacitation gestion budgétaire locale).
    \item \textbf{Pression démographique et pauvreté :} Croissance démographique forte (\textasciitilde 2,6\%/an). Pauvreté rurale $\rightarrow$ dépendance immédiate aux ressources « gratuites » (bois énergie 75\% énergie nationale, PFNL, chasse, agriculture sur brûlis). Court terme prime sur long terme.
    \item \textbf{Changement climatique :} Sécheresses récurrentes (Nord), inondations (Logone, Wouri), montée des eaux (cri de l'estuaire), modification aires de répartition espèces. Exige adaptation des plans d'aménagement (essences résilientes, couloirs de migration).
\end{enumerate}
\end{attention}

\begin{methode}[Analyser un document sur l'exploitation d'une ressource naturelle au Cameroun (Compétence évaluée au Probatoire)]
Suivez rigoureusement ces 4 étapes :
\begin{enumerate}
    \item \textbf{Identifier la ressource concernée :} Eau (superficielle/souterraine), Sol (terres arables, pâturages), Forêt (bois d'œuvre, PFNL, carbone), Biodiversité (espèces, habitats), Minerais (sable, or, bauxite, fer).
    \item \textbf{Identifier l'activité exercée :} Exploitation forestière (industrielle/artisanale/communautaire), Agriculture (itinérante, intensive, plantations), Élevage (nomade, sédentaire, ranch), Pêche (artisanale, industrielle), Minage (artisanal, industriel), Urbanisme/Infrastructures, Tourisme.
    \item \textbf{Lister les impacts décrits (négatifs / positifs) :} \textit{Négatifs :} Déforestation, érosion, pollution (chimique, sédimentaire, thermique), perte biodiversité, fragmentation habitats, conflits sociaux, émissions GES. \textit{Positifs :} Emplois, revenus (devises, budget local), infrastructures (routes, écoles, santé), transfert technologie, valorisation déchets.
    \item \textbf{Analyser les mesures de gestion proposées ou mises en œuvre :} Cadre légal (loi, décret, arrêté), Plan d'aménagement / \cle{EIES} (Étude d'Impact Environnemental et Social), Certification (FSC, RSPO), Mesures techniques (coupes sélectives, CES, traitement effluents, corridors écologiques), Mesures institutionnelles (comités de gestion, cogestion, écogardes), Mesures économiques (\cle{PSE} - Paiements pour Services Environnementaux, taxes, éco-certification), Implication populations (CLIP - Comités Locaux d'Information et de Participation, forêts communautaires).
\end{enumerate}
\textbf{Conclusion attendue :} Évaluer si les mesures sont \textbf{adéquates, suffisantes, applicables et suivies}. Proposer des améliorations concrètes (ex. « Renforcer la surveillance par drones », « Instaurer un PSE pour les agriculteurs en amont du barrage »).
\end{methode}

\begin{exoresolu}[Analyse d'un cas réel : Exploitation de la bauxite à Ngaoundal (Adamaoua)]
\textbf{Énoncé :} Une entreprise minière obtient un permis d'exploitation de bauxite sur 1 200 km² dans la région de l'Adamaoua, zone de savane boisée, tête de bassin de la Sanaga et de la Bénoué. Le projet prévoit une mine à ciel ouvert, une usine de traitement (Bayer), un chemin de fer vers Kribi. Une EIES a été réalisée. Analysez la situation.
\textbf{Solution détaillée :}
\begin{enumerate}
    \item \textbf{Ressource :} Bauxite (minerai d'aluminium, ressource non renouvelable). Ressources associées impactées : Eau (rivières Mbéré, Vina, Sanaga), Sols (ferralitiques, fragiles), Forêt/Savane (habitat antilopes, primates), Air (poussières, SO2, CO2), Paysage.
    \item \textbf{Activité :} Minage industriel à ciel ouvert (décapage stérile, extraction minerais, traitement chimique), Transport (rail/route), Infrastructure (ville minière, énergie).
    \item \textbf{Impacts (EIES type) :}
        \begin{itemize}
            \item \textit{Négatifs :} Destruction irréversible sol/sols sur emprise mine ($\sim$ 20-30 km² actifs). Risque pollution eaux (boues rouges \emph{red mud} : soude, métaux lourds, pH extrême) si digues de résidus défaillantes $\rightarrow$ catastrophe aval (Sanaga $\rightarrow$ Edéa, Douala). Érosion massive décapage. Fragmentation habitat. Poussières (silice, alumine) santé riverains/ouvriers. Afflux population $\rightarrow$ pression bois-énergie, braconnage, déchets urbains. Conflits fonciers (expropriation pasteurs/agriculteurs).
            \item \textit{Positifs :} Emplois directs/indirects (milliers). Revenus Etat (impôts, redevances minières). Infrastructures (rail, route, électricité, hôpital, école). Développement régional (Ngaoundal, Tibati).
        \end{itemize}
    \item \textbf{Mesures de gestion (à vérifier dans le Plan de Gestion Environnementale et Sociale - PGES) :}
        \begin{itemize}
            \item \textbf{Eau :} Digues de résidus aux normes internationales (géomembrane, surveillance piézométrique, traitement eaux de process recyclage). Bassins de décantation eaux de ruissellement mine. Monitoring qualité eau aval continu (pH, $Al$, $Fe$, $As$, turbidité).
            \item \textbf{Sols/Paysage :} Remise en état progressive (réhabilitation) : remise en place sol végétal, reboisement essences locales (acacias, eucalyptus, essences fourragères), création zones humides compensatoires.
            \item \textbf{Biodiversité :} Couloirs de migration maintenus. Plan anti-braconnage (écogardes financés par mine). Sauvegarde espèces menacées (ex. \emph{Hippotragus equinus koba} - cob de Buffon).
            \item \textbf{Social :} Plan de réinstallation (RAP - Resettlement Action Plan) conforme standards Banque Mondiale (IFC PS5) : compensation juste, terres de remplacement, maintien moyens d'existence. Plan de développement communautaire (PDC) financé par fonds minier local (1-2\% CA). Clause de main-d'œuvre locale.
            \item \textbf{Gouvernance :} Comité de suivi tripartite (État, Entreprise, Société Civile/Communautés). Publication contrats (ITIE - Initiative pour la Transparence dans les Industries Extractives). Audit environnemental annuel indépendant.
        \end{itemize}
    \item \textbf{Évaluation critique :} L'exploitation est \textbf{acceptable conditionnellement} si : (1) Garantie technique \textbf{absolue} sur la stabilité des digues de boues rouges (risque majeur systémique pour la Sanaga), (2) Réhabilitation \textbf{concomitante} (pas à la fin), (3) Partage des bénéfices \textbf{équitable et transparent} pour prévenir conflits sociaux, (4) Surveillance indépendante \textbf{perpétuelle} (post-fermeture). Sans cela, le coût environnemental/social dépasse le bénéfice économique.
\end{enumerate}
\end{exoresolu}

\begin{savaistu}[Le Cameroun et la « Grande Muraille Verte » (GMV)]
Lancée en 2007 par l'Union Africaine, la \cle{Grande Muraille Verte} vise à restaurer 100 millions d'hectares de terres dégradées au Sahel d'ici 2030, séquestrer 250 Mt de \ce{CO2}, créer 10 millions d'emplois verts. Le Cameroun (régions Nord, Extrême-Nord, Adamaoua-Nord) est un pays pilote. Les actions concrètes : reboisement \emph{Acacia senegal} (gomme arabique), \emph{Balanites aegyptiaca} (dattier du désert), \emph{Faidherbia albida}, fixation dunes (palmiers doum \emph{Hyphaene thebaica}, \emph{Calotropis procera}), agroforesterie, gestion pastorale (aménagement parcours, points d'eau). C'est l'application à grande échelle des techniques CES vues plus haut, couplée à une approche \textbf{paysage} (intégration agriculture-élevage-forêt-eau).
\end{savaistu}

En synthèse, le Cameroun possède un cadre légal et institutionnel solide et des exemples de réussite (forêts communautaires de l'Est, reboisement Sahel Vert, aires protégées du Dja/Campo-Ma'an). Cependant, l'effectivité de la gestion durable repose sur la résolution de l'équation : \textbf{Moyens de surveillance + Gouvernance transparente + Partage équitable des bénéfices + Adaptation au changement climatique}. C'est à cette condition que les ressources naturelles deviendront le moteur d'un développement \textbf{vraiment} durable pour la jeunesse camerounaise.

\coursec{Synthèse : analyser une situation et proposer des solutions durables}

Après avoir exploré, dans la section précédente, des exemples concrets d'exploitation et de conservation des ressources au Cameroun — du barrage de Lom Pangar aux aires protégées du Dja et de la Bénoué, en passant par la gestion communautaire des forêts du Sud — nous disposons maintenant d'une « boîte à outils » conceptuelle et factuelle. Cette dernière section a pour but de t'apprendre à \textbf{mobiliser ces connaissances} pour analyser une situation environnementale nouvelle, en identifier les causes profondes, et proposer des solutions \emph{réalistes, argumentées et durables}. C'est exactement la compétence évaluée au Probatoire : face à un document (texte, carte, graphique, photo), tu dois produire un raisonnement scientifique structuré.

\courssub{La démarche d'analyse en quatre étapes : du diagnostic à l'évaluation}

Toute analyse rigoureuse suit une démarche logique, calquée sur la méthode scientifique. Ne saute jamais d'étape : une solution proposée sans diagnostic précis est une opinion, pas une réponse scientifique.

\begin{popfigure}
\begin{tikzpicture}[
    node distance=1.6cm and 1.6cm,
    box/.style={rectangle, draw=popblue, thick, fill=popblueL, rounded corners, minimum width=3.4cm, minimum height=1.6cm, align=center, font=\small},
    arrow/.style={-Stealth, thick, popdark},
    labelbox/.style={rectangle, draw=poporange, fill=poporangeL, rounded corners, align=center, font=\scriptsize, text=popdark, inner sep=2pt}
]
% Nœuds principaux
\node[box, align=center] (diag){\textbf{1. DIAGNOSTIC}\\\emph{Quoi ? Où ? Qui ?}\\Ressource concernée\\Pressions identifiées\\Acteurs impliqués};
\node[box, right=of diag, align=center] (imp){\textbf{2. IMPACTS}\\\emph{Comment ? Pourquoi ?}\\Mécanismes physiques\\et biologiques\\Court / long terme};
\node[box, right=of imp, align=center] (sol){\textbf{3. SOLUTIONS}\\\emph{Quelles actions ?}\\Protection\\Restauration\\Alternatives économiques};
\node[box, right=of sol, align=center] (eval){\textbf{4. ÉVALUATION}\\\emph{Est-ce viable ?}\\Grille de critères\\Coût / Acceptabilité\\Efficacité / Faisabilité};

% Flèches principales
\draw[arrow] (diag.east) -- (imp.west) node[midway, above, labelbox, align=center]{Données\\de terrain};
\draw[arrow] (imp.east) -- (sol.west) node[midway, above, labelbox, align=center]{Lien cause-\\effet};
\draw[arrow] (sol.east) -- (eval.west) node[midway, above, labelbox, align=center]{Critères\\du DD};

% Boucle de rétroaction (amélioration continue)
\draw[arrow, popgreen] (eval.south) .. controls +(0,-1.6) and +(0,-1.6) .. (diag.south) node[midway, below, labelbox, fill=popgreenL, draw=popgreen] {Réajustement : suivi-évaluation (gestion adaptative)};
\end{tikzpicture}
\legende{Organigramme de la démarche d'analyse d'une situation environnementale. La flèche verte en bas symbolise la gestion adaptative : les résultats de l'évaluation nourrissent un nouveau diagnostic.}
\end{popfigure}

\textbf{Étape 1 — Diagnostic : « Quel est le patient ? »}\par
Tu dois extraire du document (ou de l'observation de terrain) trois informations non négociables :
\begin{itemize}
    \item \textbf{La ressource naturelle concernée} : eau (rivière, nappe souterraine), sol (terres agricoles, pâturages), forêt (bois d'œuvre, produits forestiers non ligneux ou PFNL), biodiversité (espèces endémiques, corridors écologiques).
    \item \textbf{Les pressions humaines directes} : déboisement (exploitation forestière, agriculture sur brûlis), pollution (rejets industriels, pesticides, plastiques), surexploitation (pêche, chasse, prélèvements d'eau), artificialisation des sols (urbanisation, barrages).
    \item \textbf{Les acteurs et leurs intérêts} : populations locales (besoins vitaux), entreprises (profit), État (régulation, recettes fiscales), ONG (conservation), administrations techniques (forêts, eau, mines). Les \emph{conflits d'usage} sont fréquents et doivent être repérés.
\end{itemize}

\textbf{Étape 2 — Analyse des impacts : « Quels sont les symptômes ? »}\par
Relie chaque pression à ses conséquences \textbf{en utilisant le vocabulaire précis du cours} :
\begin{itemize}
    \item \textbf{Sur l'eau} : érosion hydrique $\rightarrow$ turbidité, comblement du lit (colmatage), destruction des frayères, eutrophisation (apports d'azote N et de phosphore P), tarissement (perte d'infiltration), pollution chimique ou bactérienne.
    \item \textbf{Sur le sol} : perte de l'horizon humifère, tassement (compaction), salinisation, acidification, glissements de terrain sur les pentes dénudées.
    \item \textbf{Sur la forêt et la biodiversité} : fragmentation des habitats, perte de connectivité entre les populations, disparition d'espèces clés (par exemple les grands singes et éléphants, disperseurs de graines), invasion d'espèces exotiques, perturbation des cycles biogéochimiques (carbone, azote, eau).
\end{itemize}
Distingue toujours les \textbf{impacts directs} (immédiats, locaux) des \textbf{impacts indirects ou en cascade} (différés, à distance : par exemple déboisement en amont $\rightarrow$ crue dévastatrice en aval).

\textbf{Étape 3 — Proposition de solutions : « Quel traitement ? »}\par
Une solution durable ne se contente pas d'interdire ; elle \textbf{réconcilie} l'humain et son milieu. Classe tes propositions en trois niveaux complémentaires :
\begin{enumerate}
    \item \textbf{Protection / Conservation stricte} : aires protégées, zones tampons, corridors écologiques, interdiction des pratiques destructrices (feux de brousse tardifs, produits chimiques, engins de pêche prohibés).
    \item \textbf{Restauration / Réhabilitation} : reboisement avec des essences locales (jamais de monoculture d'eucalyptus !), agroforesterie, restauration des berges par végétalisation, désimperméabilisation des sols urbains.
    \item \textbf{Gestion raisonnée / Alternatives économiques} : plans d'aménagement forestier, agriculture de conservation (semis direct sous couvert végétal), pêche réglementée (tailles minimales de capture, périodes de repos biologique), écotourisme, paiements pour services environnementaux (PSE).
\end{enumerate}

\textbf{Étape 4 — Évaluation et argumentaire : « Le traitement est-il supportable ? »}\par
C'est l'étape où tu gagnes des points au Probatoire. Utilise la grille ci-dessous pour \textbf{justifier} chaque mesure proposée.

\begin{definition}[Grille d'évaluation d'une mesure de gestion durable (critères du développement durable)]
Une solution est dite « durable » si elle satisfait \emph{simultanément} les quatre critères suivants :
\begin{itemize}
    \item \textbf{Efficacité environnementale} : résout-elle le problème écologique identifié ? (Exemple : la haie vive réduit-elle l'érosion mesurée ?)
    \item \textbf{Viabilité économique / Coût} : est-elle finançable par les acteurs locaux ou l'État ? Le rapport coût/bénéfice est-il favorable à long terme ? (Exemple : l'agroforesterie rapporte du cacao \emph{et} du bois.)
    \item \textbf{Acceptabilité sociale / Équité} : les populations l'acceptent-elles ? Respecte-t-elle les droits coutumiers et les besoins vitaux (alimentation, santé) ? Il ne faut pas de « conservation forteresse » qui chasserait les populations autochtones de leurs terres.
    \item \textbf{Faisabilité technique et institutionnelle} : les compétences, la main-d'œuvre, la législation et le suivi existent-ils localement ?
\end{itemize}
Une mesure qui échoue sur un seul critère (par exemple interdire toute culture sur pente raide sans alternative de revenu = échec social) n'est \textbf{pas} une solution durable.
\end{definition}

\begin{methode}[Rédiger une réponse de type Probatoire en trois temps]
Au Probatoire, l'exercice « Analyse de document » ou « Étude de cas » suit ce canevas rigoureux :
\begin{enumerate}
    \item \textbf{Je décris la situation (diagnostic)} : « Le document montre une rivière dont le débit diminue et dont l'eau est trouble. La ressource concernée est l'eau douce et le sol du bassin versant. La pression principale est le déboisement des rives lié à l'agriculture sur brûlis pratiquée par les populations locales. Les acteurs sont... »
    \item \textbf{J'explique les impacts (mécanismes)} : « Le déboisement supprime la couverture végétale qui protégeait le sol de l'énergie des gouttes de pluie (érosion par impact ou \emph{splash}) et qui ralentissait le ruissellement en favorisant l'infiltration. Conséquence : augmentation du ruissellement de surface $\rightarrow$ érosion hydrique accrue $\rightarrow$ comblement du lit de la rivière par les sédiments $\rightarrow$ perte de la capacité de rétention $\rightarrow$ tarissement en saison sèche et crues violentes en saison des pluies, avec perte de la biodiversité de la ripisylve (forêt galerie). » \emph{Utilise les termes : infiltration, ruissellement, érosion, colmatage, frayère, corridor, service écosystémique.}
    \item \textbf{Je propose des solutions justifiées (grille DD)} : « Mesure 1 : reboisement des berges (bande de 10 à 20 m) avec des essences locales (ex. \emph{Alchornea cordifolia}, raphias). \textit{Justification} : efficace (stabilise les berges, ombrage, filtre les apports), faible coût (pépinières villageoises), acceptée (bois de feu, PFNL), faisable (technique connue). Mesure 2 : agroforesterie cacao sous arbres d'ombrage en amont. \textit{Justification} : maintient le revenu paysan tout en assurant une couverture permanente du sol... » \textbf{Termine toujours par une phrase de conclusion globale.}
\end{enumerate}
\end{methode}

\courssub{Étude de cas guidée : tarissement de la rivière Mfoundi après déboisement de son bassin versant (Yaoundé)}

\begin{experience}[Analyse documentée : le cas de la Mfoundi à Yaoundé]
\textbf{Contexte réel.} La Mfoundi, rivière emblématique traversant Yaoundé, subit un tarissement marqué en saison sèche et des crues rapides destructrices en saison des pluies. Son bassin versant a perdu l'essentiel de sa couverture forestière en quelques décennies (urbanisation rapide et mal planifiée, agriculture sur brûlis en périphérie, coupe de bois de feu).

\textbf{Données d'observation (ordres de grandeur d'après rapports publics et études universitaires) :}
\begin{itemize}
    \item Débit moyen annuel : en nette diminution depuis les années 1990, avec un étiage (débit minimal de saison sèche) de plus en plus sévère.
    \item Turbidité en saison des pluies : très supérieure aux normes de potabilité (norme OMS pour l'eau potable : turbidité inférieure à 5 NTU).
    \item Taux de couverture végétale du bassin : très inférieur au seuil d'alerte érosion (environ 30 \%).
    \item Présence de déchets plastiques et de rejets d'eaux usées non traitées en zone urbaine.
\end{itemize}

\textbf{Travail demandé :} applique la démarche en 4 étapes pour proposer un plan d'action priorisé sur 5 ans.

\medskip
\textbf{Solution guidée (modèle de rédaction Probatoire)}

\textbf{1. Diagnostic}\par
\textbf{Ressource :} l'eau de la rivière Mfoundi (ressource en eau potentielle, milieu de vie aquatique, drainage urbain).
\textbf{Pressions :} (a) déforestation massive du bassin versant (perte de la couverture végétale) $\rightarrow$ perte de la fonction « éponge » du sol ; (b) artificialisation des sols (béton, toitures) $\rightarrow$ ruissellement direct, infiltration quasi nulle ; (c) pollution ponctuelle (ordures, eaux usées) $\rightarrow$ dégradation de la qualité de l'eau.
\textbf{Acteurs :} communauté urbaine et mairies (urbanisme, déchets), MINEE (eau), MINFOF (forêts), riverains (agriculteurs, habitants), société de distribution d'eau, ONG environnementales.

\textbf{2. Analyse des impacts (mécanismes biophysiques)}\par
\begin{itemize}
    \item \textbf{Hydrologie :} sans couvert forestier, l'interception de la pluie par la canopée chute. L'énergie des gouttes détache les particules de sol (\emph{splash}) $\rightarrow$ érosion en nappe puis en ravines. Le ruissellement de surface augmente fortement (le coefficient de ruissellement passe d'environ 0,1 sous forêt à 0,8 en zone urbanisée) et concentre l'eau vers le lit mineur $\rightarrow$ \textbf{crues rapides} (temps de concentration du bassin raccourci).
    \item \textbf{Stockage :} moins d'infiltration $\rightarrow$ recharge de la nappe phréatique réduite $\rightarrow$ \textbf{tarissement de l'écoulement de base} en saison sèche : la rivière devient presque intermittente.
    \item \textbf{Morphologie et qualité :} apports massifs de sédiments (matières en suspension, MES) $\rightarrow$ \textbf{colmatage} du lit, destruction des frayères, élargissement du lit mineur. La pollution organique et azotée $\rightarrow$ consommation du dioxygène dissous par les micro-organismes décomposeurs $\rightarrow$ asphyxie de la faune aquatique.
\end{itemize}

\textbf{3. Solutions hiérarchisées (plan d'action 5 ans)}\par
\begin{center}
\begin{tabularx}{\linewidth}{|l|X|X|X|X|}\hline
\rowcolor{popblueL}\textbf{Mesure} & \textbf{Efficacité env.} & \textbf{Coût / Éco.} & \textbf{Accept. sociale} & \textbf{Faisabilité} \\ \hline
\textbf{Urgent (années 1-2) :} suppression des décharges sauvages et ramassage des ordures sur les berges & Forte (qualité de l'eau immédiate) & Faible (main-d'œuvre locale, ONG) & Forte (visibilité, santé) & Forte (mairies, travaux HIMO) \\ \hline
\textbf{Court terme (années 2-3) :} bandes tampons végétalisées de 10 à 20 m (essences ripicoles : raphias, \emph{Alchornea}, bambous) & Forte (stabilisation des berges, filtration, ombrage) & Moyen (pépinières, plantations HIMO) & Moyenne (perte de cultures en berge $\rightarrow$ compensation / PSE) & Moyenne (foncier, entretien) \\ \hline
\textbf{Moyen terme (années 3-5) :} agroforesterie et reboisement des zones amont (pentes fortes) : cacaoyers, arbres fixateurs d'azote (\emph{Leucaena}, \emph{Calliandra}) associés aux cultures vivrières & Très forte (restaure la fonction « éponge », stocke du carbone, protège le sol) & Investissement initial fort, rentable à 3-5 ans (cacao, bois, fruits) & Forte si accompagnement technique et sécurisation foncière & Moyenne (vulgarisation, pépinières nécessaires) \\ \hline
\textbf{Structurel (continu) :} assainissement (réseaux d'eaux usées, stations d'épuration) et plan d'urbanisme (zones inconstructibles = lit majeur) & Indispensable à long terme & Très fort (État et partenaires) & Complexe (relogements) & Faible à court terme, indispensable à long terme \\ \hline
\end{tabularx}
\end{center}

\textbf{4. Argumentaire de synthèse (conclusion type Probatoire)}\par
« La dégradation de la Mfoundi résulte de la rupture du cycle de l'eau par la déforestation et l'imperméabilisation du bassin versant. La solution unique n'existe pas : il faut une \textbf{approche intégrée à l'échelle du bassin versant}. La priorité immédiate est l'assainissement des berges (qualité de l'eau, santé) couplé à la restauration de la ripisylve (fonction hydrologique, biodiversité). À moyen terme, seule l'agroforesterie en amont, source de revenus pour les paysans, peut restaurer durablement la fonction « éponge » du sol et garantir le débit d'étiage. Cela nécessite une gouvernance concertée (comité de bassin) et des financements innovants (PSE, redevances sur l'eau). Sans sécurisation foncière paysanne, tout reboisement est voué à l'échec. »
\end{experience}

\begin{attention}[Piège Probatoire : les solutions « clé en main » irréalistes]
\emph{« Il faut interdire toute agriculture sur les pentes. »} $\rightarrow$ \textbf{Zéro point.} Pourquoi ? Cette mesure ignore le besoin alimentaire, le droit coutumier et l'absence d'alternative de revenu. \textbf{Sanction du correcteur :} « hors sujet / non faisable ».

\emph{« Il faut reboiser avec de l'eucalyptus, ça pousse vite. »} $\rightarrow$ \textbf{Piège écologique.} L'eucalyptus transpire abondamment et peut abaisser le niveau des nappes, appauvrit le sol et n'offre pas d'habitat à la faune locale. \textbf{Bon réflexe :} choisir des \textbf{essences locales, adaptées et multi-usages}.

\emph{« L'État doit tout payer. »} $\rightarrow$ \textbf{Non durable.} La gestion communautaire, les PSE et la responsabilité élargie du producteur sont les leviers modernes.

\textbf{Règle d'or :} une solution sans les \textbf{acteurs locaux} comme \textbf{acteurs principaux} (et pas seulement bénéficiaires) est une solution qui échouera.
\end{attention}

\begin{exemplebox}[Exemple chiffré : efficacité des haies vives et bandes boisées anti-érosives]
Des essais agronomiques menés dans l'Ouest camerounais (hautes terres de l'Ouest, pentes de 15 à 25 \%, sols ferrallitiques) comparent une parcelle témoin (labour conventionnel, sol nu) et une parcelle avec haies vives de \emph{Calliandra calothyrsus} (espacement 4 m) associées à des bandes de résidus de culture. Ordres de grandeur typiques de ce type de dispositif :
\begin{center}
\begin{tabular}{|l|c|c|c|}\hline
\rowcolor{popblueL}\textbf{Paramètre} & \textbf{Témoin (sol nu)} & \textbf{Haies vives + résidus} & \textbf{Variation} \\ \hline
Perte en sol (t/ha/an) & 42,5 & 14,8 & $-65\ \%$ \\ \hline
Ruissellement (\% de la pluie) & 38 \% & 12 \% & $-68\ \%$ \\ \hline
Rendement maïs (t/ha) & 1,2 & 2,8 & $+133\ \%$ \\ \hline
Revenu net (FCFA/ha) & 150 000 & 420 000 & $+180\ \%$ \\ \hline
\end{tabular}
\end{center}
\textbf{Interprétation :} la haie vive freine le ruissellement (effet barrière physique), ses racines structurent le sol (porosité, infiltration) et son feuillage apporte de la matière organique (fertilité). \textbf{Gagnant-gagnant :} environnement (sol, eau) + économie (rendement, revenu). C'est l'archétype de la \cle{gestion raisonnée}.
\end{exemplebox}

\begin{savaistu}[L'agroforesterie : réconcilier production et conservation au Cameroun]
Le Cameroun figure parmi les premiers producteurs mondiaux de cacao. Le modèle dominant (cacao de plein soleil, intrants chimiques) épuise les sols en 15 à 20 ans et pousse au défrichement de nouvelles parcelles forestières. L'\textbf{agroforesterie cacaoyère} associe le cacaoyer à des arbres d'ombrage (bois d'œuvre : \emph{Terminalia superba}, \emph{Triplochiton scleroxylon} ; fruitiers : safoutier, avocatier ; légumineuses fixatrices d'azote : \emph{Albizia}, \emph{Gliricidia}).
\begin{itemize}
    \item \textbf{Sols :} couverture permanente $\rightarrow$ érosion quasi nulle, matière organique et activité biologique accrues.
    \item \textbf{Eau :} infiltration améliorée, régulation du microclimat (température, hygrométrie) $\rightarrow$ meilleure résistance aux sécheresses.
    \item \textbf{Biodiversité :} corridors, habitats pour la faune (oiseaux, pollinisateurs, auxiliaires des cultures).
    \item \textbf{Climat :} stockage de carbone dans la biomasse et le sol (atténuation), ombrage (adaptation).
    \item \textbf{Socio-économie :} diversification des revenus (bois, fruits, PFNL, plantes médicinales), sécurité alimentaire, moindre dépendance aux engrais et pesticides.
\end{itemize}
Des programmes de développement accompagnent des milliers de planteurs du Sud (Ebolowa, Sangmélima) et du Centre pour rénover les vergers vieillissants. C'est la \textbf{preuve par l'exemple} que le développement durable n'est pas un slogan, mais une technique agricole rentable.
\end{savaistu}

\courssub{Bilan synthétique de la leçon : carte mentale}

Cette carte mentale résume l'architecture complète de la leçon. Utilise-la pour réviser : chaque branche principale correspond à une section du cours, les sous-branches aux notions clés à maîtriser pour le Probatoire.

\begin{popfigure}
\begin{tikzpicture}[
    mindmap,
    concept color=popblue,
    level 1 concept/.append style={level distance=4.6cm, sibling angle=72, font=\small\bfseries},
    level 2 concept/.append style={level distance=3.4cm, sibling angle=45, font=\scriptsize},
    every node/.style={concept}
]
\node[concept] {NCT11 : Réduire les conséquences néfastes des activités humaines sur les ressources naturelles}
    child[concept color=popgreen] { node[concept] {1. Ressources naturelles}
        child { node[concept] {Eau (de surface, souterraine)} }
        child { node[concept] {Sols (agricoles, forestiers)} }
        child { node[concept] {Forêts (bois, PFNL, carbone)} }
        child { node[concept] {Biodiversité (espèces, écosystèmes)} }
    }
    child[concept color=poporange] { node[concept] {2. Pressions humaines}
        child { node[concept] {Eau / Sols : déforestation $\rightarrow$ érosion ; pollution ; artificialisation} }
        child { node[concept] {Forêts / Biodiversité : exploitation illégale, brûlis, braconnage, fragmentation} }
    }
    child[concept color=poppurple] { node[concept] {3. Développement durable et gestion}
        child { node[concept] {3 piliers : environnement, social, économie} }
        child { node[concept] {Gestion raisonnée : protection, restauration, alternatives (agroforesterie, PSE, écotourisme)} }
        child { node[concept] {Outils : plans d'aménagement, études d'impact, zonage, législation} }
    }
    child[concept color=poppink] { node[concept] {4. Exemples camerounais}
        child { node[concept] {Barrages : Lom Pangar, Nachtigal, Memve'ele} }
        child { node[concept] {Aires protégées : Dja, Bénoué, Waza, Korup} }
        child { node[concept] {Forêts communautaires du Sud} }
        child { node[concept] {Lac Nyos : risque naturel et dégazage} }
        child { node[concept] {Mfoundi / Yaoundé : bassin versant urbain} }
    }
    child[concept color=popgold] { node[concept] {5. Démarche d'analyse (Probatoire)}
        child { node[concept] {Diagnostic : ressource, pressions, acteurs} }
        child { node[concept] {Impacts : mécanismes, court / long terme} }
        child { node[concept] {Solutions : protection, restauration, gestion raisonnée} }
        child { node[concept] {Évaluation : grille des 4 critères, argumentaire} }
    };
\end{tikzpicture}
\legende{Carte mentale de synthèse de la leçon NCT11. Chaque couleur correspond à une section du plan. Les concepts de niveau 2 sont les « savoirs » précis à mobiliser dans une argumentation.}
\end{popfigure}

\courssub{Entraînement Probatoire : exercice résolu type}

\begin{exoresolu}[Analyse d'un document : projet de plantation industrielle d'hévéas dans la région de l'Est]
\textbf{Document :} une multinationale obtient une concession de 45 000 ha présentée comme des « terres nues » en zone de forêt dense humide (région de l'Est) pour créer une hévéaculture industrielle. L'étude d'impact environnemental (EIE) annonce : création de 3 000 emplois, construction d'une usine de transformation, de routes, d'écoles et d'un centre de santé. Les populations locales (Baka et Bantous) craignent la perte de leurs zones de chasse, de cueillette et de leurs lieux sacrés. L'administration signale la présence d'éléphants, de gorilles et de chimpanzés dans la zone.

\textbf{Question :} en t'appuyant sur le document et sur tes connaissances, analyse les impacts prévisibles de ce projet sur les ressources naturelles et propose un plan de gestion durable pour la concession.

\tcblower
\textbf{Solution détaillée (modèle de copie)}

\textbf{1. Diagnostic de la situation}\par
\textbf{Ressources concernées :} forêt dense humide (stock de carbone, biodiversité, bois, PFNL), sols ferrallitiques (fragiles, lessivés), grands mammifères (éléphant, grands singes — espèces intégralement protégées), ressources en eau (bassins versants des affluents de la Sangha), droits coutumiers des populations autochtones (Baka).
\textbf{Pression annoncée :} conversion de forêt naturelle en monoculture industrielle d'hévéa (\emph{Hevea brasiliensis}) sur 45 000 ha (déforestation massive).
\textbf{Acteurs :} multinationale (profit, emplois), État (recettes, infrastructures), populations locales et autochtones (droits, subsistance, culture), ONG de conservation (biodiversité), administrations (MINFOF, MINEPDED, MINADER).

\textbf{2. Analyse des impacts prévisibles (mécanismes)}\par
\begin{itemize}
    \item \textbf{Biodiversité :} perte d'un habitat critique (45 000 ha) $\rightarrow$ fragmentation du paysage $\rightarrow$ isolement des populations animales $\rightarrow$ conflits homme-faune (éléphants dans les cultures) et braconnage facilité par les routes $\rightarrow$ \textbf{disparition locale d'espèces protégées}. Perte des corridors écologiques vers les aires protégées voisines (Boumba Bek, Nki).
    \item \textbf{Sols et eau :} déboisement et préparation mécanique du sol $\rightarrow$ érosion hydrique massive (sols nus en pente) $\rightarrow$ colmatage des rivières $\rightarrow$ perturbation de l'hydrologie (crues et étiages accentués). L'hévéa transpire abondamment $\rightarrow$ risque de baisse des nappes et des débits en saison sèche. Perte de matière organique et tassement du sol par les engins lourds.
    \item \textbf{Climat :} destruction du stock de carbone d'une forêt primaire (biomasse aérienne et souterraine) $\rightarrow$ émissions massives de gaz à effet de serre $\rightarrow$ contradiction avec les engagements climatiques du Cameroun (Accord de Paris, contribution nationale).
    \item \textbf{Social et droits :} perte d'accès aux ressources (chasse, cueillette, pharmacopée, sites sacrés) pour les Baka et les Bantous $\rightarrow$ insécurité alimentaire et perte d'identité culturelle. Emplois souvent précaires et peu qualifiés ; arrivée de main-d'œuvre migrante $\rightarrow$ tensions sociales et pression accrue sur les ressources résiduelles.
\end{itemize}

\textbf{3. Plan de gestion durable proposé (mesures hiérarchisées et justifiées)}\par
\begin{center}
\begin{tabularx}{\linewidth}{|l|X|X|X|X|}\hline
\rowcolor{popblueL}\textbf{Mesure} & \textbf{Efficacité env.} & \textbf{Coût / Éco.} & \textbf{Accept. sociale} & \textbf{Faisabilité} \\ \hline
\textbf{1. Zonage strict :} exclusion des zones de haute valeur de conservation (HCV) : corridors d'éléphants, sites sacrés Baka, berges (50 m), pentes fortes. Surface nette plantée inférieure à 30 000 ha. & Critique (biodiversité, eau, carbone, droits) & Perte de revenu pour l'entreprise, compensée par la certification & Forte (respect du droit coutumier, consentement des populations) & Forte (exigence légale de l'EIE, standards de certification) \\ \hline
\textbf{2. Agroforesterie et corridors :} bandes de forêt naturelle conservées tous les 100 à 200 m ; couverture végétale des inter-rangs (légumineuses : \emph{Pueraria}, \emph{Centrosema}). & Forte (connectivité, sol, eau, auxiliaires) & Faible (semences, main-d'œuvre) ; gain (moins d'engrais et d'herbicides) & Moyenne (changement de pratiques) & Forte (technique maîtrisée par la recherche agronomique) \\ \hline
\textbf{3. Partage des bénéfices / PSE :} fonds communautaire géré par un comité local (Baka, Bantous, entreprise, État) finançant écoles, santé et alternatives protéiques (élevage de volailles et de porcs, pisciculture) $\rightarrow$ réduction du braconnage. & Indirecte mais forte (réduit la pression résiduelle) & Coût supportable pour l'entreprise / investissement social & Très forte (appropriation locale) & Moyenne (gouvernance transparente requise) \\ \hline
\textbf{4. Zéro déforestation nette :} engagement « zéro déforestation, zéro tourbe, zéro exploitation » ; reboisement compensatoire hors concession (zones dégradées) avec essences locales. & Forte (carbone, image, accès aux marchés) & Coût de certification et de reboisement & Forte (exigence des marchés européen et américain) & Forte (standards internationaux adaptés à l'hévéa) \\ \hline
\textbf{5. Gestion de l'eau et des sols :} bassins de rétention, seuils en pierre, bandes enherbées ; suivi de la qualité de l'eau en aval. & Forte (ressource en eau) & Investissement en infrastructures & Forte (santé des populations) & Forte \\ \hline
\end{tabularx}
\end{center}

\textbf{4. Conclusion argumentée}\par
« Ce projet, tel que présenté initialement (monoculture industrielle en forêt primaire), est \textbf{incompatible avec le développement durable} : il sacrifie le capital naturel (biodiversité, carbone, sols, eau) et les droits humains pour un profit privé à court terme. L'acceptabilité sociale est conditionnée par le \textbf{consentement libre, éclairé et préalable} des peuples autochtones Baka. Seule une refonte profonde intégrant des \textbf{zones de conservation strictes (HCV)}, l'\textbf{agroforesterie dans les blocs de culture}, des \textbf{mécanismes équitables de partage des bénéfices} et une \textbf{certification zéro déforestation} peut le rendre viable. L'État doit conditionner le bail à ces exigences via le cahier des charges et un suivi indépendant. Sans cela, il s'agit d'un projet « brun » maquillé en vert (greenwashing), et non d'un projet durable. »
\end{exoresolu}

\begin{aretenir}[Check-list finale pour le Probatoire — Leçon NCT11]
\begin{itemize}
    \item \textbf{Vocabulaire imposé :} ressource renouvelable / non renouvelable, service écosystémique (d'approvisionnement, de régulation, de support, culturel), pression anthropique, impact direct / indirect / cumulatif, érosion hydrique / éolienne, colmatage, eutrophisation, fragmentation d'habitat, corridor écologique, espèce clé / parapluie / indicatrice, développement durable (3 piliers), gestion raisonnée / intégrée, plan d'aménagement, aire protégée (catégories de l'UICN), agroforesterie, paiement pour services environnementaux (PSE), consentement libre, éclairé et préalable, haute valeur de conservation (HCV).
    \item \textbf{Exemples camerounais à citer :} Lom Pangar / Nachtigal (barrages et mesures de compensation), Dja / Bouba Ndjida (aires protégées, braconnage), forêts communautaires du Sud (modèle de gestion locale), lac Nyos (risque naturel, dégazage), Mfoundi / Yaoundé (bassin versant urbain), agroforesterie cacaoyère du Sud et du Centre, Mont Cameroun (biodiversité, érosion, agriculture de pente).
    \item \textbf{Structure de réponse type :} \textbf{Diagnostic} (1 phrase ressource + 1 phrase pression + acteurs) $\rightarrow$ \textbf{Explication des impacts} (mécanismes biophysiques enchaînés, vocabulaire précis) $\rightarrow$ \textbf{Solutions} (au moins 3 mesures : protection, restauration, alternative) \textbf{chacune justifiée par la grille des 4 critères} $\rightarrow$ \textbf{Conclusion synthétique}.
    \item \textbf{Interdits :} « interdire tout », « l'État doit payer », « planter des eucalyptus », solutions « hors sol » (hors contexte local), oublier les acteurs locaux.
\end{itemize}
\end{aretenir}

\courssub{Exercices d'entraînement}

\begin{exercice}[Analyse de situation : le lac de retenue de Bamendjin (Ouest)]
Le lac de retenue du barrage de Bamendjin (mis en eau en 1974, environ 300 km² ; hydroélectricité, irrigation de la riziculture, pêche) subit une sédimentation rapide (perte d'environ 0,5 \% de capacité par an), une prolifération de jacinthes d'eau (\emph{Eichhornia crassipes}) et une baisse des captures de poissons. Les rives sont cultivées (maïs, pommes de terre) sans bande tampon. Les éleveurs font boire le bétail dans le lac (piétinement des berges, déjections). Les riverains se plaignent du paludisme et des odeurs.
\begin{enumerate}
    \item Réalise le \textbf{diagnostic complet} (ressource, pressions, acteurs).
    \item Explique les \textbf{mécanismes biophysiques} liant l'agriculture de rive et l'envahissement par la jacinthe d'eau.
    \item Propose \textbf{3 mesures de gestion durable} hiérarchisées et justifie chacune avec la \textbf{grille des 4 critères du développement durable}.
    \item Rédige la \textbf{conclusion argumentée} (5 lignes maximum).
\end{enumerate}
\end{exercice}

\begin{exercice}[Cas pratique : extraction de sable marin à Kribi (région du Sud)]
Pour alimenter le chantier du port en eau profonde de Kribi et la construction de routes, des dragues aspirent le sable du fond marin dans la zone estuarienne du fleuve Lobé, près des chutes de la Lobé (site touristique et zone de ponte des tortues marines). Le sable est vendu 15 000 FCFA/m³. Les pêcheurs artisanaux signalent la disparition des zones de frayères et la destruction de leurs filets par les dragues. L'État perçoit une redevance.
\begin{enumerate}
    \item Identifie le \textbf{conflit d'usage} principal et les \textbf{services écosystémiques} en jeu.
    \item Analyse l'\textbf{impact de l'extraction de sable} sur la dynamique côtière (érosion des plages, intrusion saline) et sur la biodiversité marine (tortues, poissons).
    \item Propose une \textbf{alternative durable d'approvisionnement en granulats} et une \textbf{mesure de restauration} du site, en les justifiant.
\end{enumerate}
\end{exercice}

\begin{corrige}[Exercice 1 : lac de Bamendjin]
\textbf{1. Diagnostic}\par
\textbf{Ressource :} eau stockée (énergie, irrigation, pêche, eau potable), biodiversité lacustre, sols des rives.
\textbf{Pressions :} (a) agriculture de rive sans bande tampon $\rightarrow$ érosion $\rightarrow$ sédimentation ; (b) apports d'azote et de phosphore (engrais, déjections du bétail) $\rightarrow$ eutrophisation $\rightarrow$ jacinthe d'eau ; (c) piétinement des berges par le bétail $\rightarrow$ érosion et pollution fécale ; (d) pêche non réglementée (mailles trop fines, absence de période de repos).
\textbf{Acteurs :} société de production d'électricité (hydroélectricité), sociétés de riziculture irriguée, MINEPIA (pêche et élevage), MINEPDED (environnement), mairies, riverains (paysans, éleveurs, pêcheurs), ONG.

\textbf{2. Mécanisme agriculture $\rightarrow$ jacinthe}\par
Sol nu + labour en pente $\rightarrow$ ruissellement fort $\rightarrow$ exportation des particules fines (argiles) et des éléments nutritifs (N, P, K) fixés sur ces particules $\rightarrow$ arrivée dans le lac $\rightarrow$ dépôt (colmatage) et enrichissement de l'eau en N et P $\rightarrow$ \textbf{eutrophisation} $\rightarrow$ prolifération de la jacinthe d'eau (plante flottante à croissance très rapide : sa biomasse peut doubler en 10 à 15 jours si la température dépasse 25 °C et si N et P sont disponibles). Conséquences : anoxie de l'eau (décomposition de la biomasse morte par les micro-organismes, qui consomment le dioxygène), interception de la lumière (mort des plantes submergées), obstacle à la navigation et à la pêche, gîtes larvaires pour les moustiques (paludisme).

\textbf{3. Mesures hiérarchisées}\par
\begin{enumerate}
    \item \textbf{Urgent : bandes tampons végétalisées de 20 à 50 m (interdiction de cultiver et de pâturer).} Espèces : vétiver, bambous, arbres fruitiers et fixateurs d'azote. \textit{Justification :} efficacité forte (filtre N et P, stabilise les berges), coût faible (pépinières, travaux HIMO), acceptabilité moyenne (perte de surface cultivée $\rightarrow$ compensation par PSE ou cultures pérennes en bande), faisabilité forte.
    \item \textbf{Moyen terme : gestion intégrée de la jacinthe :} enlèvement manuel et mécanique (emploi des jeunes en HIMO) + valorisation (compost, biogaz, artisanat, alimentation animale après séchage) ; la lutte biologique (charançons du genre \emph{Neochetina}) n'est envisageable qu'\textbf{après une étude de risque}. \textit{Justification :} efficacité démontrée (lac Victoria), coût modéré avec revenus de valorisation, acceptabilité forte (emplois, propreté), faisabilité bonne (technique connue).
    \item \textbf{Structurel : plan d'aménagement du bassin versant, concerté :} zonage (cultures / pâturages / forêts / zone tampon), agroforesterie en amont (réduit les apports de sédiments et de N/P à la source), assainissement des villages (latrines, fosses septiques), réglementation de la pêche (repos biologique, mailles réglementaires). \textit{Justification :} seule solution à long terme ; coût fort (État et partenaires techniques et financiers) ; acceptabilité forte si concertation ; faisabilité institutionnelle (comité de bassin).
\end{enumerate}

\textbf{4. Conclusion}\par
Le lac de Bamendjin s'envase et s'asphyxie (eutrophisation) sous l'effet combiné de l'amont (érosion, intrants agricoles) et de l'aval (surpêche, pollution). La bande tampon est le « pansement » urgent ; le plan d'aménagement concerté du bassin versant, accompagné de PSE pour les riverains, est le traitement curatif. Sans sécurisation foncière ni alternatives de revenus (agroforesterie, valorisation de la jacinthe), le lac perdra une fraction majeure de sa capacité dans les prochaines décennies, menaçant la production d'électricité et la sécurité alimentaire (riziculture) de toute la région.
\end{corrige}

\begin{corrige}[Exercice 2 : sable marin à Kribi]
\textbf{1. Conflit d'usage et services écosystémiques}\par
\textbf{Conflit :} exploitation minière (sable) \textbf{contre} pêche artisanale, tourisme (chutes de la Lobé), conservation (tortues marines, biodiversité) et protection naturelle du trait de côte.
\textbf{Services écosystémiques impactés :} \textit{approvisionnement :} poisson (pêche), sable (construction) ; \textit{régulation :} protection de la côte contre l'érosion (stock de sable), qualité de l'eau ; \textit{support :} habitats de frayères, plages de ponte des tortues marines (tortue luth, tortue verte, tortue olivâtre — espèces protégées), nurseries de juvéniles ; \textit{culturel :} tourisme des chutes de la Lobé, patrimoine des communautés côtières.

\textbf{2. Impacts de l'extraction de sable}\par
\begin{itemize}
    \item \textbf{Dynamique côtière :} le prélèvement de sable sur le fond crée un déficit sédimentaire $\rightarrow$ érosion des plages situées en aval (Kribi, Grand Batanga) $\rightarrow$ recul du trait de côte, menace sur les infrastructures, intrusion d'eau salée dans les nappes et les sols agricoles côtiers.
    \item \textbf{Biodiversité :} destruction physique des fonds (organismes benthiques, herbiers) $\rightarrow$ perte des frayères et des nurseries. Bruit et turbidité des dragues $\rightarrow$ dérangement des tortues (ponte, émergence des nouveau-nés), mortalité des juvéniles. Modification de la bathymétrie $\rightarrow$ modification des courants et de la salinité de l'estuaire.
\end{itemize}

\textbf{3. Alternative durable et restauration}\par
\begin{itemize}
    \item \textbf{Alternative :} \textbf{carrières terrestres de roches concassées (granite, gneiss) dans l'arrière-pays} + \textbf{recyclage des déblais du chantier portuaire et du béton de démolition}. \textit{Justification :} ressource abondante, impact côtier nul, création d'emplois (carrières, concassage), prix compétitif si le transport est optimisé. \textbf{Interdiction stricte du dragage dans la zone estuarienne et les zones de frayères.}
    \item \textbf{Restauration :} \textbf{réensablement artificiel des plages érodées} avec du sable propre (de carrière) + \textbf{revégétalisation des dunes} (plantes fixatrices : ipomée, filaos) pour piéger le sable. \textbf{Sanctuarisation de la zone de ponte des tortues} (interdiction nocturne, éclairage dirigé vers la terre, surveillance communautaire). \textit{Justification :} restaure la protection de la côte et l'habitat des tortues ; coût porté par l'exploitant (principe pollueur-payeur) et le port ; acceptabilité forte (tourisme, pêche) ; faisabilité bonne (techniques maîtrisées).
\end{itemize}
\end{corrige}

\newpage

\coursec{Activité d'intégration}

\coursec{NCT11 : Réduction des conséquences néfastes des activités humaines sur les ressources naturelles}

\courssub{Objectifs de l'activité}
Cette activité d'intégration vise à mobiliser l'ensemble des savoirs et savoir-faire de la leçon NCT11 dans une situation concrète et contextualisée. À l'issue de ce travail, vous devez être capables de :
\begin{itemize}
    \item \textbf{Analyser} un jeu de données environnementales (évolution temporelle de la qualité de l'eau, de la couverture forestière et des rendements agricoles) pour en extraire des tendances.
    \item \textbf{Identifier} les activités anthropiques responsables de la dégradation des ressources naturelles (eau, sol, forêt) et \textbf{construire la chaîne de causalité} reliant ces activités aux impacts observés.
    \item \textbf{Proposer} un plan de gestion durable structuré, hiérarchisé et justifié au regard des trois piliers du \cle{développement durable} (environnemental, social, économique).
    \item \textbf{Communiquer} oralement une argumentation scientifique rigoureuse en simulant une concertation locale (conseil municipal).
\end{itemize}

\courssub{Situation : Le village de Nkolbisson face à la dégradation de son environnement}
\emph{Contexte géographique et socio-économique.} Nkolbisson est un village périurbain situé à une quinzaine de kilomètres au sud de Yaoundé, dans la région du Centre. La population, estimée à environ 3\,500 habitants, vit principalement d'une agriculture de subsistance (manioc, maïs, plantain, légumes) et de la cueillette de produits forestiers non ligneux (\emph{Gnetum africanum} « eru », champignons, fruits sauvages). Le village s'approvisionne en eau potable grâce à une source captée sur le flanc de la colline de Nkolbisson, alimentée par un petit bassin versant boisé de 120 ha.

\emph{Évolution récente.} Depuis une dizaine d'années, l'extension de la ville de Yaoundé a engendré une forte demande en matériaux de construction. Trois carrières de sable (exploitation artisanale et semi-mécanisée) se sont installées sur les flancs de la colline, empiétant sur la forêt. Parallèlement, la pression foncière a poussé les agriculteurs à défricher les dernières zones boisées en pente pour étendre les champs de manioc.

\emph{Constats des villageois.} Depuis 5 ans, les plaintes se multiplient : « L'eau de la source devient trouble à la moindre pluie, elle a mauvais goût et les enfants ont souvent des diarrhées » (Mme Atangana, mère de famille). « Nos champs de manioc donnent de moins en moins ; le sol est lessivé, on voit les racines à nu » (M. Essomba, chef cultivateur). « Les carrières font du bruit, de la poussière, et les camions abîment notre unique piste » (Jeunes du village).

\emph{Dossier technique fourni.} Le conseil municipal a commandité une étude simplifiée. Voici les données synthétisées sur la période 2014--2023.

\begin{center}
\begin{tabularx}{\linewidth}{|c|X|X|X|}\hline
\rowcolor{popblueL}
\textbf{Année} & \textbf{Turbidité de l'eau à la source (NTU)} & \textbf{Superficie boisée résiduelle du bassin versant (ha)} & \textbf{Rendement moyen du manioc (t/ha)} \\ \hline
2014 & 4,2 & 118 & 18,5 \\ \hline
2015 & 5,1 & 112 & 17,8 \\ \hline
2016 & 7,8 & 101 & 15,2 \\ \hline
2017 & 12,5 & 89 & 12,9 \\ \hline
2018 & 18,3 & 76 & 10,4 \\ \hline
2019 & 25,6 & 65 & 8,7 \\ \hline
2020 & 32,1 & 58 & 7,2 \\ \hline
2021 & 38,9 & 52 & 6,5 \\ \hline
2022 & 45,4 & 47 & 5,8 \\ \hline
2023 & 52,7 & 42 & 5,1 \\ \hline
\end{tabularx}
\end{center}
\legende{Évolution de trois indicateurs environnementaux et agricoles à Nkolbisson (2014--2023). Norme OMS turbidité eau de boisson : $< 5$ NTU.}

\emph{Témoignage complémentaire.} « Avant, la forêt retenait l'eau. Quand il pleuvait, l'eau filtrait doucement et arrivait claire à la source. Maintenant, avec les carrières et les champs nus en pente, la boue descend tout droit dans la source. Le sable des carrières est lavé par la pluie et bouche le captage. Nous avons tout perdu : l'eau, la terre, la forêt. » — Extrait du procès-verbal de la concertation villageoise du 12 mars 2024.

\courssub{Consignes}
Travaillant en groupes de 4 à 5 élèves, vous réaliserez les tâches suivantes :

\begin{enumerate}
    \item \textbf{Représentation graphique.} À l'aide d'un logiciel tableur (ou sur papier millimétré), tracez l'évolution des trois indicateurs (turbidité, surface boisée, rendement manioc) en fonction du temps (2014--2023). Utilisez un graphique à double échelle (axe gauche pour surface boisée et rendement, axe droit pour turbidité). Choisissez des échelles adaptées, légendez axes, unités, courbes et titre.
    \item \textbf{Analyse des tendances et seuils.} Décrivez l'évolution de chaque indicateur. Repérez l'année où la turbidité franchit le seuil de potabilité OMS (5 NTU). Calculez le taux de perte de surface boisée et le taux de baisse du rendement sur la période.
    \item \textbf{Chaîne de causalité.} Identifiez les activités humaines directes (pressions) et construisez une chaîne de causalité complète (minimum 5 maillons) reliant ces activités aux trois impacts majeurs observés (pollution eau, érosion sol, baisse rendement). Utilisez des flèches $\rightarrow$ et du vocabulaire précis (ruissellement, érosion hydrique, colmatage, lessivage, matière organique).
    \item \textbf{Plan de gestion durable (5 mesures).} Proposez 5 mesures concrètes, applicables à Nkolbisson, hiérarchisées de la plus urgente/structurante à la plus complémentaire. Pour \emph{chaque} mesure, précisez :
    \begin{itemize}
        \item L'action technique ou organisationnelle.
        \item Le(s) pilier(s) du \cle{développement durable} concerné(s) : \textbf{Environnement}, \textbf{Social}, \textbf{Économique}.
        \item L'acteur principal responsable de la mise en œuvre (mairie, villageois, État, ONG, exploitants).
        \item Un indicateur de suivi pour évaluer l'efficacité de la mesure dans 3 ans.
    \end{itemize}
    \item \textbf{Restitution orale.} Désignez un porte-parole. Vous avez 3 minutes pour présenter votre plan au « Conseil municipal » (la classe). La présentation doit suivre l'ordre : Diagnostic $\rightarrow$ Causalité $\rightarrow$ Plan d'action $\rightarrow$ Appel à décision.
\end{enumerate}

\courssub{Ressources à mobiliser}
\begin{itemize}
    \item \textbf{Notion de développement durable :} « Développement qui répond aux besoins du présent sans compromettre la capacité des générations futures de répondre aux leurs » (Rapport Brundtland, 1987). Trois piliers indissociables : \textbf{Environnement} (préservation ressources, biodiversité), \textbf{Social} (équité, santé, culture), \textbf{Économique} (viabilité activités, revenus).
    \item \textbf{Chaîne de causalité environnementale :} Pression (activité) $\rightarrow$ Facteur physique/chimique modifié $\rightarrow$ Impact sur le milieu $\rightarrow$ Conséquence sur les \cle{services écosystémiques} $\rightarrow$ Impact sur les sociétés humaines.
    \item \textbf{Érosion hydrique et ruissellement :} Déforestation $\rightarrow$ perte couverture végétale $\rightarrow$ augmentation ruissellement (coefficient de ruissellement $\uparrow$) $\rightarrow$ détachement particules sol (érosion) $\rightarrow$ transport sédiments $\rightarrow$ colmatage cours d'eau/captages.
    \item \textbf{Turbidité :} Mesure de la trouble de l'eau (unités néphélométriques NTU). Liée aux matières en suspension (argiles, limons, matière organique, micro-organismes). Seuil OMS eau de boisson : \SI{5}{NTU}.
    \item \textbf{Gestion raisonnée des carrières :} Réhabilitation progressive (remblai, reboisement), bassins de décantation, zones tampons riveraines, plan d'exploitation validé par étude d'impact.
\end{itemize}

\courssub{Démarche et corrigé commenté}
\begin{exoresolu}{Corrigé type et commentaires pédagogiques}
\emph{\textbf{1. Représentation graphique (Consigne 1)}}
Le graphique ci-dessous utilise un double axe des ordonnées pour respecter les unités distinctes (ha et t/ha à gauche, NTU à droite). Il met en évidence des tendances opposées : la surface boisée et le rendement décroissent de façon quasi-linéaire, tandis que la turbidité croît de façon exponentielle (accélération marquée après 2018).

\begin{popfigure}
\begin{tikzpicture}
\begin{axis}[
    width=\linewidth, height=10cm,
    xlabel={Année},
    ylabel={Surface boisée (ha) / Rendement (t/ha)},
    ylabel style={text=popblue},
    xmin=2013.5, xmax=2023.5,
    ymin=0, ymax=130,
    xtick={2014,2015,2016,2017,2018,2019,2020,2021,2022,2023},
    xticklabel style={rotate=45, anchor=east, font=\small},
    legend pos=north west,
    legend cell align=left,
    grid=major,
    grid style={popline},
    axis line style={popdark},
    tick style={popdark},
    label style={font=\small},
    title={Évolution des indicateurs à Nkolbisson (2014--2023)},
    title style={font=\bfseries\small, yshift=1.1cm},
    axis y line*=left
]
% Surface boisée (ha)
\addplot[popblue, very thick, mark=square*, mark size=2.5] coordinates {
    (2014,118) (2015,112) (2016,101) (2017,89) (2018,76) (2019,65) (2020,58) (2021,52) (2022,47) (2023,42)
};
\addlegendentry{Surface boisée (ha)}

% Rendement manioc (t/ha)
\addplot[popgreen, very thick, mark=triangle*, mark size=2.5] coordinates {
    (2014,18.5) (2015,17.8) (2016,15.2) (2017,12.9) (2018,10.4) (2019,8.7) (2020,7.2) (2021,6.5) (2022,5.8) (2023,5.1)
};
\addlegendentry{Rendement manioc (t/ha)}
\end{axis}

\begin{axis}[
    width=\linewidth, height=10cm,
    ylabel={Turbidité (NTU)},
    ylabel style={text=poporange},
    xmin=2013.5, xmax=2023.5,
    ymin=0, ymax=60,
    xtick={2014,2015,2016,2017,2018,2019,2020,2021,2022,2023},
    xticklabels={},
    axis y line*=right,
    axis x line=none,
    legend pos=north east,
    legend cell align=left,
    label style={font=\small},
]
% Turbidité (NTU)
\addplot[poporange, very thick, mark=*, mark size=2.5] coordinates {
    (2014,4.2) (2015,5.1) (2016,7.8) (2017,12.5) (2018,18.3) (2019,25.6) (2020,32.1) (2021,38.9) (2022,45.4) (2023,52.7)
};
\addlegendentry{Turbidité (NTU)}

% Ligne seuil OMS
\draw[popred, dashed, thick] (2013.5,5) -- (2023.5,5) node[right, font=\tiny, popred] {Seuil OMS (5 NTU)};
\end{axis}
\end{tikzpicture}
\legende{Graphique à double échelle : évolution opposée de la couverture forestière et des rendements (axe gauche, décroissance) vs turbidité (axe droit, croissance accélérée). Franchissement du seuil de potabilité OMS (5 NTU) dès 2015.}
\end{popfigure}

\emph{\textbf{2. Analyse des tendances et seuils (Consigne 2)}}
\begin{itemize}
    \item \textbf{Turbidité :} Croissance lente (2014--2017) puis exponentielle (2018--2023). Franchissement du seuil OMS (\SI{5}{NTU}) \textbf{dès 2015} (5,1 NTU). En 2023, elle est \textbf{10 fois supérieure} à la norme (52,7 NTU).
    \item \textbf{Surface boisée :} Perte quasi-linéaire de \SI{76}{ha} en 10 ans (118 $\rightarrow$ 42 ha). Taux de déforestation moyen : \SI{7,6}{ha/an}. \textbf{Perte relative : 64\%} de la couverture initiale.
    \item \textbf{Rendement manioc :} Baisse continue de \SI{13,4}{t/ha} (18,5 $\rightarrow$ 5,1 t/ha). \textbf{Baisse relative : 72\%}. Corrélation visuelle forte avec la perte de surface boisée (coefficient de corrélation $r$ proche de -1 si calculé).
\end{itemize}

\emph{\textbf{3. Chaîne de causalité (Consigne 3)}}
Voici une chaîne de causalité structurée (les maillons peuvent être détaillés différemment tant que la logique est respectée) :

\begin{center}
\begin{tabularx}{\linewidth}{|X|}\hline
\rowcolor{popblueL} \textbf{Chaîne de causalité principale à Nkolbisson} \\ \hline
\textbf{Pression 1 : Déforestation (champs + carrières)} sur pentes $\rightarrow$ \textbf{Perte couverture végétale} (racines, litière) $\rightarrow$ \textbf{Augmentation ruissellement surfacique} (vitesse, volume) $\rightarrow$ \textbf{Érosion hydrique accélérée} (détachement/transport particules sol) $\rightarrow$ \textbf{Colmatage source / Augmentation turbidité} (Matières En Suspension) $\rightarrow$ \textbf{Eau impropre à la consommation} (risque sanitaire) \textbf{ET} \textbf{Appauvrissement sol} (perte horizon argilo-humique, lessivage éléments nutritifs) $\rightarrow$ \textbf{Baisse rendements agricoles} (manioc) $\rightarrow$ \textbf{Insécurité alimentaire / Perte revenus} villageois. \\ \hline
\end{tabularx}
\end{center}

\emph{\textbf{4. Plan de gestion durable : 5 mesures hiérarchisées (Consigne 4)}}

\begin{center}
\begin{tabularx}{\linewidth}{|c|X|X|X|X|}\hline
\rowcolor{popblueL}
\textbf{N°} & \textbf{Mesure (Action technique/organisationnelle)} & \textbf{Piliers DD} & \textbf{Acteur principal} & \textbf{Indicateur de suivi (3 ans)} \\ \hline
1 & \textbf{Arrêt immédiat de l'extension des carrières} et mise en place d'un \textbf{Plan de Réhabilitation Environnementale} (remblaiement, reprofilage pentes, bassins de décantation obligatoires, reboisement espèces pionnières locales \emph{Acacia, Leucaena}). & \textbf{Env} (stop érosion, restauration sol) \textbf{Eco} (maintien activité encadrée) & Mairie / MINEPDED / Exploitants (sous contrôle) & \% surface carrières réhabilitée ; Turbidité source $< 20$ NTU. \\ \hline
2 & \textbf{Reconstitution de la couverture forestière du bassin versant} : Reboisement communautaire des 76 ha perdus (agroforesterie \emph{Gliricidia/manioc}, boisements purs pentes $> 30\%$), création d'une \textbf{Aire de Protection de la Source} (rayon 200 m, interdiction culture/pâturage). & \textbf{Env} (régulation hydrique, biodiversité) \textbf{Soc} (eau potable, cadre vie) \textbf{Eco} (bois-énergie, fruits à terme) & Villageois / ONG / MINFOF (appui plants) & Surface reboisée (ha) ; Taux de survie plants $> 80\%$ ; Turbidité $< 30$ NTU (tendance baissière). \\ \hline
3 & \textbf{Adoption de pratiques agroécologiques} : Culture sur billons/ridelles suivant courbes de niveau, paillage résidus manioc, associations manioc/légumineuses (haricot, soja), haies vives \emph{Calliandra} en bordure champs. & \textbf{Env} (lutte érosion, fertilité sol) \textbf{Eco} (rendement $\uparrow$, intrants $\downarrow$) \textbf{Soc} (sécurité alimentaire, charge travail femmes $\downarrow$) & Cultivateurs / MINADER (vulgarisation) & Rendement manioc $> 12$ t/ha ; Taux matière organique sol $\uparrow$. \\ \hline
4 & \textbf{Mise en place d'une gestion communautaire de l'eau} : Protection physique captage (clôture, regard étanche), installation filtre lent sable / chloration domestique promotion, comité de gestion eau villageois (cotisation maintenance). & \textbf{Soc} (santé, accès eau) \textbf{Env} (protection ressource) \textbf{Eco} (coût santé $\downarrow$) & Comité villageois / Mairie / ONG (appui technique) & Taux maladies diarrhéiques $-50\%$ ; Fonctionnement comité (réunions, caisse). \\ \hline
5 & \textbf{Diversification économique et valorisation} : Appui filières produits forestiers non ligneux (PFNL : eru, miel), écotourisme léger (randonnée colline reboisée), formation jeunes maintenance filtres/agroforesterie. Réduction dépendance carrières. & \textbf{Eco} (revenus alternatives) \textbf{Soc} (emploi jeunes, genre) \textbf{Env} (valorisation forêt debout) & Mairie / MINEPAT / Partenaires (PNDP, ONG) & Nb ménages revenus diversifiés ; Taux chômage jeunes $\downarrow$. \\ \hline
\end{tabularx}
\end{center}

\emph{\textbf{5. Restitution orale (Consigne 5)}}
Le porte-parole structure son intervention (3 min) :
\begin{enumerate}
    \item \textbf{Diagnostic (45 s) :} « À Nkolbisson, 10 ans d'exploitation non régulée (carrières, déforestation) ont détruit 64\% de la forêt du bassin versant. Conséquence : l'eau est 10 fois trop trouble (53 NTU vs 5 NTU norme), les sols sont lessivés, le manioc a perdu 72\% de rendement. »
    \item \textbf{Causalité (45 s) :} « La chaîne est claire : sol nu $\rightarrow$ ruissellement $\rightarrow$ érosion $\rightarrow$ boue dans source + terre partie des champs. »
    \item \textbf{Plan d'action (90 s) :} « 5 mesures : 1. Stopper l'hémorragie (carrières + bassins décantation). 2. Recoudre le tissu végétal (reboisement + protection source). 3. Cultiver autrement (agroécologie). 4. Boire sain (gestion eau communautaire). 5. Vivre mieux sans détruire (alternatives économiques). »
    \item \textbf{Appel (30 s) :} « Monsieur le Maire, ce plan coûte moins cher que l'inaction. Nous demandons l'arrêté municipal d'urgence pour la mesure 1 et le budget communal pour les plants mesure 2. Le village s'engage sur les mesures 3, 4, 5. »
\end{enumerate}
\end{exoresolu}

\courssub{Bilan de l'activité}
Cette activité a permis de mettre en évidence plusieurs points fondamentaux du programme :
\begin{itemize}
    \item \textbf{L'interdépendance des ressources :} La dégradation de la forêt (ressource bois/biodiversité) entraîne directement la dégradation de l'eau (ressource vitale) et du sol (ressource productive). Agir sur l'une sans les autres est inefficace.
    \item \textbf{La non-linéarité des réponses écologiques :} La turbidité n'augmente pas proportionnellement à la déforestation ; elle explose après un seuil de dégradation du bassin versant (perte de la capacité de tampon/filtration). Cela justifie l'urgence d'agir \emph{avant} le point de basculement.
    \item \textbf{L'approche systémique du développement durable :} Aucune mesure isolée ne suffit. La mesure technique (bassins de décantation) échouera sans la mesure foncière (reboisement) et sans l'appropriation sociale (comité eau, agroécologie). La hiérarchisation (Urgence $\rightarrow$ Structurel $\rightarrow$ Accompagnement) est la marque d'une gestion \emph{raisonnée}.
    \item \textbf{L'ancrage local :} Les solutions (espèces pionnières locales, billons courbes de niveau, comité villageois) sont adaptées au contexte pédoclimatique (ferralsols, climat équatorial) et socio-culturel (agriculture familiale, chefferie, mairie) de Nkolbisson / Yaoundé.
    \item \textbf{Compétences mobilisées :} Lecture critique de données (tableau, graphique), modélisation (chaîne causalité), argumentation multidimensionnelle (3 piliers DD), communication orale structurée.
\end{itemize}

\begin{attention}[Pièges fréquents à éviter dans votre copie/oral]
\begin{itemize}
\item Ne pas confondre \emph{turbidité} (trouble, MES) et \emph{pollution chimique} (pesticides, nitrates) : ici, c'est une pollution physique/sédimentaire.
\item Oublier de quantifier : « la forêt a diminué » $\rightarrow$ « la forêt a perdu 64\% (76 ha) ».
\item Proposer des mesures « clé en main » venues d'ailleurs (ex. « construire une usine de traitement d'eau ») inadaptées au village (coût, maintenance, énergie).
\item Négliger le pilier \textbf{Économique} : si les exploitants de carrières et les paysans n'ont pas d'alternatives de revenus, ils contourneront les interdictions.
\item Présenter des mesures en « liste de courses » sans hiérarchisation ni articulation (quelle mesure conditionne les autres ?).
\end{itemize}
\end{attention}

\begin{savaistu}[Le contexte réel de Nkolbisson et la gestion des carrières au Cameroun]
Le village de Nkolbisson existe réellement (arrondissement de Yaoundé VI). Il illustre parfaitement les conflits d'usage périurbains décrits dans la \textbf{Loi n°98/005 du 14 avril 1998 relative au régime de l'eau} (pollution, police de l'eau) et la \textbf{Loi n°94/01 du 20 janvier 1994 sur les forêts} (domaine forestier permanent/non permanent). L'exploitation artisanale des carrières de sable est encadrée par le \textbf{Code minier (Loi n°2016/017)} qui impose une \textbf{Étude d'Impact Environnemental et Social (EIES)} et un \textbf{Plan de Réhabilitation} pour toute carrière, mais l'application reste faible face à la pression urbaine. Le \textbf{Plan National de Développement (PND 2020-2030)} et la \textbf{Stratégie Nationale de Développement Durable (SNDD30)} prônent l'approche « Paysage » (intégration forêt-eau-sol-agriculture) que vous venez de simuler.
\end{savaistu}

\newpage

\coursec{Méthodes et exercices résolus}

\courssub{Méthodes et exercices résolus}

\courssub{Savoir-faire 1 : Identifier les impacts d'une activité humaine donnée sur une ressource naturelle}

\begin{methode}[Analyse d'impacts anthropiques sur une ressource naturelle]
    \textbf{Objectif :} Passer d'une description d'activité à une liste structurée de pressions (prélèvement, pollution, fragmentation, introduction d'espèces) et d'impacts (qualitatifs et quantitatifs) sur la ressource visée.
    \begin{enumerate}
        \item \textbf{Identifier l'activité et la ressource cible} : Nommer précisément l'activité (ex. : irrigation par aspersion, exploitation forestière industrielle, urbanisation non planifiée) et la ressource concernée (nappe phréatique, sol ferrallitique, forêt dense humide, population de \emph{Prunus africana}).
        \item \textbf{Lister les pressions directes} (mécanismes physiques/chimiques immédiats) :
        \begin{itemize}
            \item \textbf{Prélèvement / Surexploitation} : débit pompé $>$ recharge ; volume de bois prélevé $>$ accroissement naturel.
            \item \textbf{Pollution} : apports d'intrants (nitrates, phosphates, pesticides), rejets thermiques, sédimentation excessive, métaux lourds.
            \item \textbf{Modification physique de l'habitat} : déforestation, drainage, imperméabilisation des sols, barrages, pistes forestières (fragmentation).
            \item \textbf{Introduction d'espèces envahissantes} : volontaires (aquaculture, reboisement monospécifique) ou accidentelles (eau de ballast, matériel agricole).
        \end{itemize}
        \item \textbf{Qualifier les impacts} (conséquences sur l'état de la ressource) :
        \begin{itemize}
            \item \textbf{Réversibilité} : réversible à court/moyen/long terme ou irréversible (extinction locale, perte de fertilité définitive, salinisation).
            \item \textbf{Échelle spatiale} : locale (champ, parcelle), bassin versant, régionale, globale (climat).
            \item \textbf{Échelle temporelle} : aiguë (crue de pollution), chronique (eutrophisation lente), retardée (bioaccumulation).
        \end{itemize}
        \item \textbf{Utiliser une matrice Pression–État–Réponse (PER) ou DPSIR} (Drivers, Pressions, État, Impacts, Réponses) pour structurer la réponse écrite.
        \item \textbf{Rédiger la conclusion} : « L'activité X exerce une pression de type Y sur la ressource Z, entraînant un impact [réversible/irréversible] caractérisé par [indicateurs mesurables : baisse du niveau piézométrique de 2 m/an, perte de 15 t/ha/an de sol, disparition de 3 espèces endémiques]. »
    \end{enumerate}
    \textbf{Reflexe Probatoire :} Toujours citer des \textbf{indicateurs chiffrés} (débits, concentrations, surfaces, taux d'érosion, richesse spécifique) si le document les fournit, ou les estimer à partir de données de référence (ex. : taux d'érosion tolérable $T \approx 1$ à $5$ t/ha/an en zone tropicale humide).
\end{methode}

\begin{exoresolu}[Impacts de l'irrigation intensive dans la plaine de la Logone (Extrême-Nord)]
    \textbf{Énoncé :} Dans la plaine de la Logone, l'irrigation gravitaire du riz (deux campagnes/an) prélève $120 \times 10^6$ m$^3$/an d'eau du fleuve. Les sols, initialement ferrallitiques peu profonds sur matériel alluvial, reçoivent $250$ kg/ha/an d'urée (46\% N) et $100$ kg/ha/an de NPK (15-15-15). Le drainage est insuffisant. Le fleuve Logone alimente aussi le parc national de Waza en aval. À partir de ces données, identifiez les pressions et leurs impacts sur l'eau, le sol et la biodiversité.
    \tcblower
    \textbf{Solution détaillée :}
    \begin{enumerate}
        \item \textbf{Analyse de l'activité :} Riziculture irriguée intensive, double campagne, intrants chimiques élevés, drainage défaillant, zone soudano-sahélienne (évapotranspiration potentielle $ETP \approx 2000$ mm/an).
        \item \textbf{Matrice Pression–Impact :}
        \begin{center}
            \begin{tabularx}{\linewidth}{|l|X|X|}\hline
            \rowcolor{popblueL}
            \textbf{Compartiment} & \textbf{Pressions identifiées} & \textbf{Impacts qualifiés (avec indicateurs)} \\ \hline
            \textbf{Eau (Logone \& nappe)} & \begin{itemize}\item Prélèvement $120 \times 10^6$ m$^3$/an (environ 30\% du module sec du fleuve).\item Retour d'eau drainage chargé en \ce{NO3-}, \ce{PO4^3-}, pesticides.\item Évaporation nette forte (perte d'eau pure).\end{itemize} & \begin{itemize}\item \textbf{Baisse du niveau d'eau} en aval : assèchement partiel des mares de Waza (impact sur la faune).\item \textbf{Eutrophisation} : $[\ce{NO3-}] > 50$ mg/L (norme OMS eau potable) dans les drains ; prolifération de \emph{Typha australis} (envahissante).\item \textbf{Salinisation} de la nappe par concentration évaporative (conductivité $> 2000$ $\mu$S/cm).\end{itemize} \\ \hline
            \textbf{Sol} & \begin{itemize}\item Apport azoté total : $250 \times 0,46 + 100 \times 0,15 = 130$ kg N/ha/an.\item Apport $\ce{P2O5}$ : $15$ kg/ha/an ; $\ce{K2O}$ : $15$ kg/ha/an.\item Absence de drainage $\rightarrow$ montée de la nappe, capillarité.\end{itemize} & \begin{itemize}\item \textbf{Acidification} : nitrification $\ce{NH4+ -> NO3- + 2H+}$ $\rightarrow$ lessivage des bases (Ca, Mg), pH chute $< 5,5$.\item \textbf{Salinisation / Sodification} : accumulation de sels solubles (ECe $> 4$ dS/m) et $\ce{Na+}$ échangeable (ESP $> 15\%$) $\rightarrow$ dégradation structurelle (dispersion argile, croûtes).\item \textbf{Pollution diffuse} : lessivage \ce{NO3-} vers nappe ($> 50$ mg/L).\end{itemize} \\ \hline
            \textbf{Biodiversité (Milieu aquatique \& Waza)} & \begin{itemize}\item Réduction débit aval (connectivité fleuve-plaine inondable rompue).\item Pollution azotée/phosphatée.\item Monoculture riz (perte habitats naturels).\end{itemize} & \begin{itemize}\item \textbf{Disparition frayères} : poissons migrateurs (ex. \emph{Alestes baremoze}) bloqués.\item \textbf{Perte richesse spécifique} : remplacement herbiers aquatiques divers par \emph{Typha} monospécifique.\item \textbf{Impact faune grande} : éléphants, kob de Buffon, girafes de Waza privés de points d'eau secs $\rightarrow$ conflits homme-faune.\end{itemize} \\ \hline
        \end{tabularx}
    \end{center}
    \item \textbf{Synthèse rédigée type Probatoire :}
    « L'irrigation rizicole intensive de la plaine de la Logone exerce trois pressions majeures : (1) un prélèvement hydrique excessif ($120$ Mm$^3$/an) réduisant le débit d'étiage du Logone de $\sim 30\%$, asséchant les mares du parc de Waza ; (2) une pollution diffuse azotée ($130$ kg N/ha/an apportés, moins de 50\% exportés par la récolte) entraînant eutrophisation (\emph{Typha}) et contamination de la nappe ($[\ce{NO3-}] > 50$ mg/L) ; (3) une dégradation pédologique par salinisation/sodification (montée nappe + évaporation) et acidification (nitrification), menaçant la durabilité des sols alluviaux. Ces impacts sont \textbf{réversibles à long terme} seulement si le drainage est maîtrisé et les intrants rationalisés ; la perte d'espèces piscicoles locales peut être \textbf{irréversible}. »
    \end{enumerate}
    \textbf{Conclusion :} L'activité compromet la multifonctionnalité du bassin (eau potable, agriculture, biodiversité, tourisme) par une surexploitation de l'eau et une minéralisation non maîtrisée des sols.
\end{exoresolu}

\courssub{Savoir-faire 2 : Proposer des mesures de gestion durable adaptées à une situation locale}

\begin{methode}[Élaboration de mesures de gestion durable (locale, concertée, mesurable)]
    \textbf{Objectif :} Passer du diagnostic d'impacts à un plan d'actions hiérarchisé (Éviter – Réduire – Compenser – Restaurer) avec indicateurs de suivi, adapté au contexte socio-économique local (Cameroun).
    \begin{enumerate}
        \item \textbf{Hiérarchiser les mesures selon la séquence \cle{Éviter – Réduire – Compenser – Restaurer} (séquence ERC) :}
        \begin{itemize}
            \item \textbf{Éviter} : Changer de localisation, de technologie, de calendrier (ex. : ne pas drainer une zone humide, interdire l'exploitation dans une aire protégée).
            \item \textbf{Réduire} : Optimiser l'existant (efficacité irrigation, agriculture de conservation, traitements phytosanitaires ciblés, coupes sélectives à faible impact - RIL).
            \item \textbf{Compenser} : Créer un gain écologique équivalent ailleurs (réhabilitation zone humide, corridor biologique, reboisement espèces locales) \emph{uniquement si résidus inévitables}.
            \item \textbf{Restaurer} : Réhabiliter les milieux déjà dégradés (reboisement, amendements calcaires, revégétalisation berges).
        \end{itemize}
        \item \textbf{Intégrer les piliers du développement durable} pour chaque mesure :
        \begin{itemize}
            \item \textbf{Environnement} : efficacité écologique (biodiversité, qualité eau/sol, carbone).
            \item \textbf{Économie} : coût d'investissement / fonctionnement, rentabilité pour l'exploitant, création d'emplois verts.
            \item \textbf{Social} : acceptabilité populations locales, équité d'accès à la ressource, savoirs locaux, genre.
        \end{itemize}
        \item \textbf{Définir des indicateurs SMART} (Spécifiques, Mesurables, Atteignables, Réalistes, Temporels) pour le suivi :
        \begin{itemize}
            \item Ex. : « Taux de couverture du sol $> 30\%$ en permanence d'ici 2026 » ; « Concentration \ce{NO3-} à la sortie du bassin $< 25$ mg/L d'ici 3 ans » ; « Surface en agroforesterie $+ 500$ ha/an ».
        \end{itemize}
        \item \textbf{Identifier les acteurs et le cadre juridique camerounais} :
        \begin{itemize}
            \item \textbf{État} : MINEPDED, MINFOF, MINADER, MINEA (Police de l'eau, Police des forêts).
            \item \textbf{Collectivités} : Communes (Plans Communaux de Développement - PCD), Comités de Gestion des Ressources Naturelles (COGERN).
            \item \textbf{Textes} : Loi n°96/12 (Cadre environnement), Loi n°94/01 (Forêts), Loi n°98/005 (Régime de l'eau), Décret n°2013/0171 (Étude d'impact), Plan National de Développement (PND 2020-2030).
        \end{itemize}
        \item \textbf{Rédiger la proposition} sous forme de tableau « Mesure – Acteurs – Indicateur – Échéance – Coût estimé ».
    \end{enumerate}
    \textbf{Reflexe Probatoire :} Ne proposez jamais une mesure « idéale » mais irréaliste (ex. « arrêter toute agriculture »). Privilégiez l'\textbf{agroécologie}, l'\textbf{agroforesterie}, la \cle{gestion intégrée des ressources en eau (GIRE)} et la \cle{foresterie communautaire}. Citez explicitement un texte de loi ou une institution camerounaise.
\end{methode}

\begin{exoresolu}[Plan de gestion durable pour la plaine de la Logone (suite ex. précédent)]
    \textbf{Énoncé :} À partir des impacts identifiés (salinisation sols, pollution azotée, baisse débit aval, perte biodiversité Waza), proposez un plan de gestion durable concerté pour les 5 prochaines années. Précisez pour chaque mesure : le type (Éviter/Réduire/Compenser/Restaurer), les acteurs principaux, un indicateur SMART et le texte de référence camerounais.
    \tcblower
    \textbf{Solution détaillée :}
    \begin{center}
        \begin{tabularx}{\linewidth}{|l|X|l|X|l|}\hline
        \rowcolor{popblueL}
        \textbf{Problème cible} & \textbf{Mesure proposée (Type ERC)} & \textbf{Acteurs clés} & \textbf{Indicateur SMART (5 ans)} & \textbf{Cadre juridique / Réf. Cameroun} \\ \hline
        \textbf{Prélèvement eau excessif} & \textbf{Réduire} : Passage irrigation gravitaire $\rightarrow$ goutte-à-goutte / pivot basse pression. Calendrier cultural décalé (éviter campagne sèche critique). & MINADER, MINEA, ONAD (Office National de Développement de la Riziculture), Groupements riziculteurs. & \begin{itemize}\item Efficacité irrigation $> 70\%$ (vs 40\% gravitaire).\item Prélèvement réduit $< 80 \times 10^6$ m$^3$/an.\item Débit réservé Waza respecté $> 5$ m$^3$/s en étiage.\end{itemize} & Loi n°98/005 Art. 14 (Débit réservé) ; Décret 2001/165 (Police de l'eau) ; PND 2020-2030 Axe 3. \\ \hline
        \textbf{Pollution azotée / Eutrophisation} & \textbf{Réduire} : Fertilisation raisonnée (analyse sol, fractionnement apports, engrais à libération contrôlée). Bandes enherbées 10 m bord drains/fleuve (tampon). & MINADER (vulgarisation), IRAD (recherche variétés efficientes N), Communes (PCD). & \begin{itemize}\item Dose N $\downarrow$ 130 $\rightarrow$ 80 kg/ha/an.\item Rendement maintenu $> 5$ t/ha.\item $[\ce{NO3-}]$ drains $< 25$ mg/L.\item Largeur bandes tampon $100\%$ drains principaux.\end{itemize} & Loi n°96/12 Art. 17 (Pollution) ; Arrêté 007/MINEPDED (Normes rejet) ; Code de bonnes pratiques agricoles (MINADER). \\ \hline
        \textbf{Salinisation / Acidification sols} & \textbf{Restaurer + Réduire} : Drainage souterrain collecteur + lessivage contrôlé. Amendement calcaire (chaux agricole) selon analyse. Rotation Riz – Légumineuse (niébé, arachide) / Jachère améliorée (\emph{Mucuna}, \emph{Stylosanthes}). Agroforesterie (haies \emph{Acacia nilotica}, \emph{Faidherbia albida}). & IRAD (cartes fertilité), MINADER, ONG (ex. CIFOR, WWF), Groupements. & \begin{itemize}\item pH sol $> 5,5$ ; ESP $< 10\%$ ; ECe $< 4$ dS/m sur 80\% périmètre.\item Taux matière organique $\uparrow$ 0,5\%/an.\item Surface agroforesterie $> 20\%$ périmètre.\end{itemize} & Loi n°94/01 Art. 26 (Gestion durable forêts/terroirs) ; Stratégie Nationale de Lutte contre la Désertification (SNLD). \\ \hline
        \textbf{Perte biodiversité / Conflits Homme-Faune (Waza)} & \textbf{Compenser + Restaurer} : Réhabilitation corridor écologique Logone-Waza (revégétalisation rives). Points d'eau artificiels hors parc pour faune. Aires de pâturage gérées (zonage) pour pasteurs. Écotourisme communautaire (revenus alternatifs). & MINFOF (Conservateur Waza), MINEPDED, Communes, ONG (WCS, IUCN), Populations riveraines (COGERN). & \begin{itemize}\item Corridor fonctionnel (largeur $> 500$ m) restauré 5 km.\item 5 points d'eau hors parc opérationnels.\item Conflits H-F déclarés $\downarrow 50\%$.\item Revenus écotourisme $> 5$ M FCFA/an/village.\end{itemize} & Loi n°94/01 (Aires protégées) ; Décret 95/466 (Forêts communautaires) ; Plan de Gestion Parc Waza. \\ \hline
        \textbf{Gouvernance / Concertation} & \textbf{Éviter (conflits)} : Création/mise en œuvre effective d'un \textbf{Comité de Bassin Logone} (GIRE). Plan d'Allocation de l'Eau (PAE). Contrats d'objectifs riziculteurs / éleveurs / conservation. & MINEA (Chef Service Eau), Gouverneur Région, Communes, Représentants usagers (riziculteurs, éleveurs, pêcheurs, conservation). & \begin{itemize}\item Comité réuni 2 fois/an.\item PAE validé an 1.\item 100\% conflits résolus par médiation.\end{itemize} & Loi n°98/005 Art. 38 (Comités de bassin) ; Décret 2001/165 Art. 44 (Police de l'eau). \\ \hline
        \end{tabularx}
    \end{center}
    \textbf{Conclusion :} Ce plan intègre la \cle{GIRE} et l'\cle{agroécologie}. Il est \textbf{réaliste} car il maintient l'activité rizicole (sécurité alimentaire, revenus) tout en restaurant les fonctions écologiques (eau, sol, biodiversité) via des mesures techniques accessibles (goutte-à-goutte, rotation, haies) et une gouvernance locale renforcée (Comité de bassin, COGERN), conformément aux lois 98/005 et 94/01.
\end{exoresolu}

\courssub{Savoir-faire 3 : Analyser un document sur l'exploitation d'une ressource naturelle au Cameroun}

\begin{methode}[Analyse critique de document (texte, carte, graphique, tableau) sur une ressource camerounaise]
    \textbf{Objectif :} Extraire, structurer et critiquer l'information pour répondre à une question de type Probatoire (souvent : « Montrez que... », « Expliquez... », « Proposez... »).
    \begin{enumerate}
        \item \textbf{Identifier la nature et la source du document} : Type (rapport MINEPDED, article presse, carte MINFOF, graphique FAO, photo satellite, témoignage ONG), Auteur, Date, Contexte (projet, conflit, politique publique). \emph{Vérifier la fiabilité / biais potentiel.}
        \item \textbf{Extraire les données brutes} (chiffres, dates, localisations, noms d'acteurs, espèces, lois citées) : les surligner / lister. \textbf{Unités obligatoires.}
        \item \textbf{Structurer l'information selon le modèle DPSIR} (Drivers – Pressions – État – Impacts – Réponses) ou \textbf{Analyse SWOT} (Forces, Faiblesses, Opportunités, Menaces) si c'est une évaluation de projet.
        \item \textbf{Croiser avec les connaissances du cours} (notions : surexploitation, services écosystémiques, valeur d'option, valeur d'existence, principe pollueur-payeur, principe de précaution, droits coutumiers vs droit moderne).
        \item \textbf{Rédiger la réponse argumentée} :
        \begin{itemize}
            \item Phrase d'introduction : « Le document X (source, date) présente la situation de la ressource Y dans la zone Z... »
            \item Argumentation paragraphée : 1 idée = 1 paragraphe = 1 donnée du doc + 1 connaissance cours.
            \item Utiliser le vocabulaire précis : \cle{déforestation}, \cle{dégradation}, \cle{fragmentation}, \cle{braconnage}, \cle{exploitation illicite}, \cle{certification FSC/PEFC}, \cle{foresterie communautaire}, \cle{reboisement}, \cle{restauration}.
            \item Conclure par une évaluation critique : durabilité ? équité ? efficacité des réponses ?
        \end{itemize}
    \end{enumerate}
    \textbf{Reflexe Probatoire :} \textbf{Citez explicitement le document} (« Selon le document... », « La figure 1 montre que... »). Ne récitez pas le cours seul. Si c'est un graphique : lisez les axes, les unités, les tendances, les ruptures. Si c'est une carte : légende, échelle, localisation précise (région, département, aire protégée).
\end{methode}

\begin{exoresolu}[Analyse de l'exploitation du bois d'œuvre dans le massif forestier du Dja (Sud-Est)]
    \textbf{Énoncé :} Document : Extrait du rapport « État des Forêts du Bassin du Congo 2021 » (OFAC/CBF) + Carte MINFOF 2020.
    \begin{itemize}
        \item Le massif du Dja (Réserve de Faune du Dja, Patrimoine UNESCO + UFA périphériques) couvre 1,6 million ha.
        \item Production bois d'œuvre 2020 : $1,2 \times 10^6$ m$^3$ (dont 60\% Ayous, 15\% Sapelli, 10\% Azobé).
        \item Taux de déforestation nette 2010-2020 : $0,15\%$/an dans les UFA certifiées FSC ; $0,45\%$/an dans les UFA non certifiées ; $0,02\%$/an dans la Réserve (cœur).
        \item Braconnage : 120 indices/km$^2$ (pièges, cartouches) en périphérie UFA vs 5 indices/km$^2$ en cœur de Réserve.
        \item Populations locales (Baka, Bantu) : 45 villages, 12 000 hab. Dépendance produits forestiers non ligneux (PFNL) : \emph{Irvingia} (bush mango), \emph{Gnetum} (okok), chenilles, miel. Revenus PFNL $\approx 30\%$ revenu monétaire.
        \item Foresterie communautaire : 2 forêts communautaires créées (2018, 2020), 5 000 ha chacune, plans de gestion simples validés.
    \end{itemize}
    \textbf{Question :} En vous appuyant sur le document et vos connaissances, analysez la durabilité de l'exploitation forestière dans ce massif et proposez deux mesures pour améliorer la prise en compte des populations locales.
    \tcblower
    \textbf{Solution détaillée :}
    \begin{enumerate}
        \item \textbf{Analyse DPSIR du document :}
        \begin{itemize}
            \item \textbf{Drivers} : Demande internationale bois tropicaux (Chine, UE), demande urbaine bois énergie/PFNL, pression démographique, pauvreté rurale.
            \item \textbf{Pressions} : Exploitation industrielle (UFA) sélective mais intensive (60\% Ayous = surexploitation potentielle essence légère), pistes forestières (fragmentation), braconnage facilité par pistes (indices $\times 24$ en périphérie), agriculture itinérante sur brûlis aux abords.
            \item \textbf{État} : Déforestation faible en cœur (protection efficace UNESCO/MINFOF), modérée en UFA certifiées ($0,15\%$/an), forte en UFA non certifiées ($0,45\%$/an = 3$\times$ plus). Perte connectivité (corridors rompus).
            \item \textbf{Impacts} : Biodiversité : déclin grands mammifères (éléphant, gorille, bongo) en périphérie. Socio-économique : Privation PFNL (bush mango, okok) pour 12 000 personnes ; conflits homme-faune ; revenus précaires (braconnage, exploitation illicite).
            \item \textbf{Réponses} : Certification FSC (efficace : déforestation $\div 3$), Foresterie communautaire (naissante : 10 000 ha sur 1,6 M ha = 0,6\%), Surveillance (écogardes), Écotourisme (potentiel Dja).
        \end{itemize}
        \item \textbf{Évaluation de la durabilité (3 piliers) :}
        \begin{center}
            \begin{tabularx}{\linewidth}{|l|X|X|}\hline
            \rowcolor{popblueL}
            \textbf{Pilier} & \textbf{Points forts (Document + Cours)} & \textbf{Points faibles / Menaces} \\ \hline
            \textbf{Écologique} & Cœur Réserve intact (Patrimoine UNESCO). Certification FSC réduit impact UFA (RIL, zones de conservation HVC). & \textbf{Surexploitation essences prisées} (Ayous, Sapelli) : rotation 30 ans insuffisante pour gros diamants. \textbf{Fragmentation} par pistes. \textbf{Braconnage} massif périphérie (effet « vide sanitaire »). Perte connectivité paysages. \\ \hline
            \textbf{Économique} & Revenus export bois ($> 300$ Mrd FCFA/an national). Emplois UFA (qualifiés/non qualifiés). Certification = accès marchés premium (UE). & \textbf{Rente captée par États/Entreprises} (peu retombées locales). PFNL sous-valorisés (pas de transformation locale). Foresterie communautaire marginale (0,6\%). \\ \hline
            \textbf{Social} & Droits d'usage reconnus en théorie (Loi 94/01 Art. 26). Forêts communautaires créées (initiative positive). & \textbf{Exclusion réelle} : Baka/Bantu peu consultés (CLIP - Consultation Libre, Préalable, Éclairée souvent formelle). Conflits accès ressources. Braconnage = survie. Peu d'emplois durables locaux. \\ \hline
        \end{tabularx}
    \end{center}
    \item \textbf{Synthèse argumentée :}
    « L'exploitation dans le massif du Dja présente une \textbf{durabilité contrastée} : écologiquement viable dans le cœur protégé et les UFA certifiées (déforestation $0,15\%$/an), elle devient \textbf{non durable} dans les UFA non certifiées ($0,45\%$/an) et socialement \textbf{inéquitable} (populations riveraines exclues des bénéfices, dépendantes PFNL menacés par la fragmentation et le braconnage facilité par les pistes forestières). La certification FSC prouve son efficacité technique (RIL, HVC) mais ne règle pas le partage des bénéfices. »
    \item \textbf{Deux mesures pour améliorer la prise en compte des populations locales :}
    \begin{enumerate}
        \item \textbf{Extension effective de la foresterie communautaire et des forêts de terroir} : Viser 50 000 ha (3\% massif) en 5 ans. \textbf{Appui technique/financier} (MINADER, MINFOF, ONG, Fonds Forestier National) pour : inventaires PFNL, plans de gestion simples \textbf{intégrant PFNL et agriculture}, transformation locale (huile bush mango, séchage okok), contractualisation vente directe (équitable). \textit{Base légale : Loi 94/01 Art. 26, 37 ; Décret 95/531 (Forêts communautaires).}
        \item \textbf{Mise en œuvre rigoureuse du mécanisme de \cle{Partage des Bénéfices} (REDD+ / FSC / Loi 94/01 Art. 74)} : Conditionner le renouvellement des titres UFA à la signature de \textbf{Conventions de Partenariat Locales (CPL)} contraignantes : \% chiffre d'affaires (ex. 10\% taxe de reboisement + 5\% fonds communautaire) géré par COGERN (Comités de Gestion) élus (représentativité Baka/Bantu/Genre). Financement : écoles, santé, micro-projets agroforesterie, surveillance participative (brigades villageoises anti-braconnage). \textit{Indicateur : Nb CPL signées / Nb villages ; Montants versés / an / village.}
    \end{enumerate}
    \end{enumerate}
    \textbf{Conclusion :} La durabilité passe par l'élargissement de l'assise sociale de la gestion (foresterie communautaire, CPL contraignantes) et l'extension des pratiques certifiées (FSC) à l'ensemble des UFA, sous contrôle indépendant (Observatoire Indépendant).
\end{exoresolu}

\courssub{Savoir-faire complémentaire : Calculer un indicateur de pression ou d'état (Empreinte écologique, Taux d'érosion, Indice de biodiversité)}

\begin{methode}[Calcul et interprétation d'indicateurs quantitatifs de pression sur les ressources]
    \textbf{Objectif :} Maîtriser les formules simples exigibles au Probatoire pour quantifier une pression ou un état.
    \begin{enumerate}
        \item \textbf{Taux de déforestation annuel (\%/an) :}
        \[ \tau = \frac{S_{t_0} - S_{t_1}}{S_{t_0}} \times \frac{100}{t_1 - t_0} \]
        $S$ = surface forestière (ha) ; $t$ = année. \textbf{Attention :} « net » = déforestation – régénération/reboisement.
        \item \textbf{Taux d'érosion moyen (t/ha/an) - Équation universelle de perte de sol (USLE) simplifiée / RUSLE :}
        \[ A = R \times K \times LS \times C \times P \]
        $R$ = érosivité pluie (MJ mm/ha/h/an), $K$ = érodibilité sol, $LS$ = pente/longueur, $C$ = couverture végétale, $P$ = pratiques antiérosives.
        \textbf{Au Probatoire :} On donne souvent $A$ final ou on demande de comparer $C$ (ex. forêt $C=0,001$ vs champ nu $C=1$).
        \item \textbf{Empreinte écologique (ha équivalent global / hab) :} $EF = \sum \frac{Consommation_i}{Rendement_i} \times Facteur\_équivalence_i$. Comparer à la \textbf{Biocapacité} disponible. Si $EF > Biocapacité$ $\rightarrow$ \textbf{Déficit écologique} (surconsommation).
        \item \textbf{Indice de Shannon-Weaver (Diversité spécifique) :}
        \[ H' = -\sum_{i=1}^{S} p_i \ln p_i \quad \text{avec } p_i = \frac{n_i}{N} \]
        $S$ = richesse spécifique, $n_i$ = effectif espèce $i$, $N$ = effectif total. $H'_{max} = \ln S$. \textbf{Équitabilité :} $E = H' / H'_{max}$.
        \item \textbf{Indice de qualité de l'eau (ex. IBGN simplifié, ou concentration) :} Respect des normes (OMS, Arrêté 007/MINEPDED) : $[\ce{NO3-}] < 50$ mg/L, $[\ce{PO4^3-}] < 0,5$ mg/L, $DBO_5 < 25$ mg/L, $O_2 > 5$ mg/L.
    \end{enumerate}
    \textbf{Reflexe :} Toujours écrire la \textbf{formule}, \textbf{remplacer les symboles par les valeurs avec unités}, \textbf{calculer}, \textbf{conclure avec unité} et \textbf{interpréter} (seuil, tendance, comparaison référence).
\end{methode}

\begin{exoresolu}[Calcul du taux de déforestation et empreinte écologique - Cas du Mont Cameroun]
    \textbf{Énoncé :} Données pour la région du Mont Cameroun (zone de forêt dense humide de basse et moyenne altitude) :
    \begin{itemize}
        \item Surface forestière 2000 : $250\,000$ ha.
        \item Surface forestière 2020 : $212\,500$ ha.
        \item Population 2020 : $500\,000$ hab.
        \item Consommation moyenne bois énergie : $1,2$ m$^3$/hab/an (rendement charbonnage $20\%$).
        \item Productivité forêt naturelle : $5$ m$^3$/ha/an (accroissement naturel).
        \item Surface agricole utile (SAU) : $80\,000$ ha (rendement maïs $1,5$ t/ha/an).
        \item Besoin alimentaire : $200$ kg céréales/equ/hab/an.
    \end{itemize}
    \textbf{Questions :}
    \begin{enumerate}
        \item Calculez le taux de déforestation annuel moyen sur la période 2000-2020. Interprétez.
        \item Calculez la surface forestière nécessaire pour fournir durablement le bois énergie (approvisionnement = accroissement). Comparez à la surface restante. Concluez sur la durabilité de la filière bois énergie.
        \item Calculez l'empreinte écologique « terres agricoles » pour l'alimentation en céréales (équivalence terre agricole = 1). La biocapacité agricole locale est de $80\,000$ ha. Y a-t-il déficit ?
    \end{enumerate}
    \tcblower
    \textbf{Solution détaillée :}
    \begin{enumerate}
        \item \textbf{Taux de déforestation :}
        \[ \tau = \frac{250\,000 - 212\,500}{250\,000} \times \frac{100}{20} = \frac{37\,500}{250\,000} \times 5 = 0,15 \times 5 = \mathbf{0,75\% / an} \]
        \textbf{Interprétation :} Ce taux ($0,75\%$/an) est \textbf{très élevé} (moyenne nationale $\approx 0,1\%$/an, moyenne bassin du Congo $\approx 0,17\%$/an). Il traduit une pression anthropique intense (agriculture sur brûlis, exploitation bois d'œuvre/énergie, expansion urbaine) menant à une perte de $37\,500$ ha en 20 ans. Si ce taux persiste, la forêt disparaîtrait en $\approx 133$ ans (mais accélération probable).
        \item \textbf{Bilan bois énergie :}
        \begin{itemize}
            \item Demande bois énergie brut (bois vert) : $500\,000 \text{ hab} \times 1,2 \text{ m}^3/\text{hab/an} = \mathbf{600\,000 \text{ m}^3/an}$.
            \item Rendement charbonnage $20\%$ $\rightarrow$ Demande bois sur pied = $600\,000 / 0,20 = \mathbf{3\,000\,000 \text{ m}^3/an}$.
            \item Production durable (accroissement) = Surface $\times$ Productivité = $S_{forêt} \times 5 \text{ m}^3/\text{ha/an}$.
            \item Surface nécessaire pour production durable : $S_{nec} = \frac{3\,000\,000}{5} = \mathbf{600\,000 \text{ ha}}$.
        \end{itemize}
        \textbf{Comparaison :} Surface forestière restante (2020) = $212\,500$ ha $< 600\,000$ ha nécessaires.
        \textbf{Conclusion :} \textbf{Déficit massif} (facteur 2,8). La filière bois énergie est \textbf{structurellement non durable} : elle mine le capital forestier (prélèvement $>$ accroissement), accélérant la déforestation calculée en 1. Solutions : reboisement énergétique (eucalyptus, acacia), foyers améliorés (réduction 40-50\% conso), substitution (GPL, biogaz).
        \item \textbf{Empreinte écologique « Céréales » :}
        \begin{itemize}
            \item Demande totale : $500\,000 \text{ hab} \times 200 \text{ kg/hab/an} = 100\,000\,000 \text{ kg} = \mathbf{100\,000 \text{ t/an}}$.
            \item Production locale : $80\,000 \text{ ha} \times 1,5 \text{ t/ha/an} = \mathbf{120\,000 \text{ t/an}}$.
            \item Surface nécessaire (empreinte) : $\frac{100\,000 \text{ t}}{1,5 \text{ t/ha}} = \mathbf{66\,667 \text{ ha}}$.
        \end{itemize}
        \textbf{Comparaison :} Empreinte ($66\,667$ ha) $<$ Biocapacité agricole locale ($80\,000$ ha).
        \textbf{Conclusion :} \textbf{Excédent biophysique} pour les céréales ($+13\,333$ ha). \textbf{Mais} : cela ne compte pas les autres aliments (protéines, légumes, huiles), ni les terres pour le bois énergie ($600\,000$ ha !), ni l'empreinte carbone, ni l'urbanisation. L'empreinte écologique \textbf{globale} est donc très probablement en \textbf{déficit sévère}.
    \end{enumerate}
\end{exoresolu}

\courssub{Savoir-faire complémentaire : Argumenter une décision d'aménagement (Étude d'Impact Environnemental et Social - ÉIES)}

\begin{methode}[Structuration d'un avis d'ÉIES / Analyse coûts-bénéfices environnementaux]
    \textbf{Objectif :} Simuler l'instruction d'un dossier de projet (barrage, mine, plantation, route) selon la procédure camerounaise (Décret 2013/0171).
    \begin{enumerate}
        \item \textbf{Décrire le projet et son environnement de référence} : Localisation, composantes (ouvrages, infrastructures), emprise, alternatives (site zéro, site alternatif, technologie alternative). État initial (milieu physique, biologique, humain).
        \item \textbf{Identifier et qualifier les impacts} (Méthode Leopold modifiée ou grille) : Nature (positif/négatif), Portée (direct/indirect/cumulatif), Étendue (locale/régionale), Durée (temporaire/permanent), Réversibilité, Intensité (faible/moyenne/forte/critique). \textbf{Significativité = f(Intensité, Étendue, Durée, Réversibilité).}
        \item \textbf{Proposer la séquence ERC} (Éviter, Réduire, Compenser) pour \textbf{chaque impact significatif négatif}. Mesures techniques précises, coûts, responsable, planning.
        \item \textbf{Plan de Surveillance Environnementale (PSE) :} Indicateurs, fréquence, points de mesure, seuils d'alerte, responsable, budget.
        \item \textbf{Analyse Coûts-Bénéfices (économique élargie) :} Intégrer les \textbf{services écosystémiques} perdus/gagnés (valeur d'usage direct, indirect, option, existence). \textbf{Actualisation} (taux d'actualisation $t$).
        \item \textbf{Avis motivé :} « Avis favorable sous réserves » (mesures ERC + PSE obligatoires) / « Avis défavorable » (impacts critiques irréversibles non compensables, ex. extinction espèce endémique, déplacement population sans solution).
    \end{enumerate}
    \textbf{Reflexe Probatoire :} Citez le \textbf{Décret 2013/0171} (procédure ÉIES), la \cle{hiérarchie ERC}, la \cle{participation du public} (enquête publique), le \cle{Plan de Gestion Environnementale et Sociale (PGES)}. Distinguez \textbf{impact résiduel} (après mesures ERC) et \textbf{impact initial}.
\end{methode}

\begin{exoresolu}[Avis sur le projet de barrage hydroélectrique de Nachtigal (Sanaga) - Cas simplifié]
    \textbf{Énoncé :} Projet : Barrage de Nachtigal (420 MW), fleuve Sanaga, région Centre. Retenue : $4,6$ km$^2$. Déplacement : $250$ personnes. Zone : forêt dégradée, proximité parc national de la Mbam-Djerem. Espèces clés : poisson endémique \emph{Steatocranus} sp. (rheophile), tortue \emph{Trionyx triunguis}. Activités locales : pêche artisanale, agriculture, sable. Documents fournis : Étude d'impact du promoteur (NHPC), Avis ONG (WWF, Brainforest), Rapport MINEPDED.
    \textbf{Travail :} En tant qu'expert MINEPDED, rédigez les points clés de votre avis technique (Impacts significatifs, Mesures ERC, PSE, Avis final).
    \tcblower
    \textbf{Solution détaillée (Structure avis type) :}
    \begin{enumerate}
        \item \textbf{Impacts significatifs identifiés (analyse critique du promoteur) :}
        \begin{itemize}
            \item \textbf{Négatifs forts :}
            \begin{itemize}
                \item \textbf{Fragmentation fleuve / Perte habitat rheophile :} Barrage infranchissable pour \emph{Steatocranus} sp. (endémique Sanaga) $\rightarrow$ \textbf{risque extinction locale}. \textit{Irréversible sans dispositif franchissement efficace.}
                \item \textbf{Modification régime hydrologique aval :} Hydroélectrique = éclusées (variations rapides niveau/débit) $\rightarrow$ perturbation frayères, ripisylve, prélèvements eau potable (Edea, Yaoundé). \textit{Cumulatif avec barrages existants (Song Loulou, Edéa).}
                \item \textbf{Emission GES (réservoir) :} Décomposition biomasse forêt inondée ($4,6$ km$^2$) $\rightarrow$ $\ce{CH4}$, $\ce{CO2}$ (durée 10-20 ans). \textit{Sous-estimé par promoteur.}
                \item \textbf{Déplacement populations / Perte terres/sable :} 250 pers. + perte accès ressources (pêche, sable). Plan de Réinstallation (PAR) insuffisant (terres remplacement moins fertiles, pas de titre foncier).
            \end{itemize}
            \item \textbf{Positifs :} Production électricité bas carbone ($420$ MW, $\approx 2,9$ TWh/an) $\rightarrow$ substitution thermique (fuel/charbon), appui développement industriel, revenus État, emploi construction.
        \end{itemize}
        \item \textbf{Mesures ERC exigées (au-delà du promoteur) :}
        \begin{center}
            \begin{tabularx}{\linewidth}{|l|X|l|}\hline
            \rowcolor{popblueL}
            \textbf{Impact} & \textbf{Mesure ERC (Exigence MINEPDED)} & \textbf{Responsable / Délai} \\ \hline
            Fragmentation / \emph{Steatocranus} & \textbf{Réduire + Compenser} : \textbf{Passes à poissons performantes} (type échelle à bassins + attraction débit) \textbf{OBLIGATOIRES} (pas optionnelles). Suivi génétique populations amont/aval. \textbf{Compensation} : Sanctuarisation affluent non barré (ex. Mbam) pour population refuge. & NHPC (Concessionnaire) / Avant mise en eau. MINFOF / An 1. \\ \hline
            Régime hydrologique aval & \textbf{Réduire} : \textbf{Débit réservé écologique} dynamique (mimant crue naturelle), pas seulement $Q_{min}$ fixe. Gestion éclusées : vitesse montée/descente limitée ($< 10$ cm/min). Station limnimétrique temps réel aval. & NHPC / Continu. MINEA (Police eau) contrôle. \\ \hline
            GES Réservoir & \textbf{Éviter/Réduire} : \textbf{Déboisement total} zone inondable avant mise en eau (valorisation bois, réduction biomasse). \textbf{Compenser} : Reboisement $2 \times$ surface réservoir ($920$ ha) espèces locales à croissance rapide. & NHPC / Avant impoundment. MINEPDED / 5 ans. \\ \hline
            Populations déplacées & \textbf{Compenser (PAR renforcé)} : Terres remplacement \textbf{équivalentes} (potentiel agronomique) + \textbf{titres fonciers} (Certificats Fonciers). Indemnisation cultures/arbres à valeur marchande réelle. Programme restauration moyens d'existence (pisciculture, maraîchage irrigué, formation). Suivi socio-éco 5 ans. & NHPC / MINDAF (Domaine) / Avant déplacement. \\ \hline
        \end{tabularx}
    \end{center}
        \item \textbf{Plan de Surveillance Environnementale (PSE) - Indicateurs clés :}
        \begin{itemize}
            \item \textbf{Physico-chimie eau} : $O_2$, $T^\circ$, turbidité, nutriments (amont, aval immédiat, aval 10 km) - \textit{Mensuel + Crue/Étiage}.
            \item \textbf{Ichtiofaune} : Abondance \emph{Steatocranus} (pêche électrique standardisée), diversité, franchissement passes - \textit{Trimestriel}.
            \item \textbf{Hydrologie} : Niveau/débit continu (amont, aval), respect débit réservé - \textit{Continu (automate)}.
            \item \textbf{Social} : Niveau vie ménages réinstallés (revenus, sécurité alimentaire), conflits usage - \textit{Annuel 5 ans}.
            \item \textbf{GES} : Flux $\ce{CH4}/\ce{CO2}$ surface réservoir (chambres flottantes) - \textit{Annuel 10 ans}.
        \end{itemize}
        \item \textbf{Avis final motivé :}
        \textbf{« AVIS FAVORABLE SOUS RÉSERVES »}
        \begin{quote}
        Le projet Nachtigal répond à un besoin énergétique national majeur (PND 2030) avec une empreinte carbone faible \emph{in fine}. Toutefois, les impacts résiduels sur la biodiversité endémique (\emph{Steatocranus}) et les populations riveraines sont inacceptables sans renforcement majeur des mesures ERC. L'avis favorable est conditionné à :
        \begin{enumerate}
            \item La réalisation effective des passes à poissons \textbf{avant la première mise en eau} (preuve de fonctionnement).
            \item La validation par MINEPDED/MINEA du régime de débit réservé dynamique et de la gestion des éclusées.
            \item La finalisation du PAR conforme aux standards de la Banque Mondiale (OP 4.12) / BAD (OS 2) avec titres fonciers.
            \item Le dépôt d'une caution financière (garantie bancaire) couvrant le coût des mesures compensatoires (reboisement, sanctuarisation affluent).
        \end{enumerate}
        Le Plan de Surveillance Environnementale (PSE) joint doit être mis en œuvre par un bureau d'études indépendant, financé par NHPC, avec rapports publics trimestriels.
        \end{quote}
    \end{enumerate}
\end{exoresolu}

\courssub{Savoir-faire complémentaire : Évaluer un projet d'agroforesterie / restauration (Arbre à la parcelle)}

\begin{methode}[Conception et évaluation d'un système agroforestier (Diagnostic – Conception – Suivi)]
    \textbf{Objectif :} Proposer une association arbres-cultures-élevage adaptée au contexte pédoclimatique et socio-économique camerounais (zone forestière, soudanaenne, sahélienne).
    \begin{enumerate}
        \item \textbf{Diagnostic de la parcelle / terroir :} Type de sol (ferralsol, vertisol, arenosol), climat ( pluviométrie, saison sèche), culture principale (cacao, café, maïs, coton, mil), contraintes (érosion, fertilité déclinante, ravageurs, manque bois d'œuvre/énergie), main-d'œuvre disponible, tenure foncière (coutumier / titre).
        \item \textbf{Choix des espèces arborescentes (Critères) :}
        \begin{itemize}
            \item \textbf{Fonctionnelles :} Fertilité (légumineuses fixatrices N : \emph{Acacia}, \emph{Leucaena}, \emph{Gliricidia}, \emph{Faidherbia} - inversion cycle), Érosion (racines profondes, haies anti-érosives : \emph{Vetiver}, \emph{Gliricidia}), Ombrage (cacao/café : \emph{Albizia}, \emph{Erythrina}, \emph{Milicia}), Fourrage (saison sèche : \emph{Faidherbia}, \emph{Leucaena}, \emph{Moringa}), Bois/Énergie/Fruits/Médecine (valeur marchande).
            \item \textbf{Compatibilité :} Racines (éviter compétition eau/nutriments : pivotantes vs traçantes), Phénologie (déférence hydrique : \emph{Faidherbia} feuillue saison sèche = ombre + azote + fourrage SANS compétition culture pluviale), Allélopathie (tester \emph{Leucaena} / maïs).
            \item \textbf{Disponibilité plants :} Pépinières locales, ONG, IRAD, MINADER.
        \end{itemize}
        \item \textbf{Modèle spatial (Architecture) :}
        \begin{itemize}
            \item \textbf{Haies alignées / Bandes enherbées} (contour lignes de niveau) : anti-érosif, bois, fourrage. Espacement 10-20 m.
            \item \textbf{Parc à \emph{Faidherbia} / \emph{Acacia} / \emph{Vitellaria} (karité) / \emph{Parkia} (néré)} : arbres épars en plein champ (10-50 t/ha). Modèle sahélien / soudanien.
            \item \textbf{Système à étages (Forêt-jardin / Cacaoyère complexe)} : Strates multiples (arbres haut, moyen, bas, culture, couverture sol). Modèle zone forestière (Sud, Centre, Est).
            \item \textbf{Bordures / Haies vives / Brise-vent} : Délimitation, protection, biodiversité.
        \end{itemize}
        \item \textbf{Gestion technique (Itinéraire) :} Plantation (début saison des pluies, poquets $40 \times 40$ cm, amendement organique), Protection (pare-bêtes, paillage), Taille (formation, recépage, pollardage - \textbf{période : fin saison sèche} pour minimiser compétition), Récolte (bois, fruits, feuilles).
        \item \textbf{Évaluation multicritères (Tableau de bord) :}
        \begin{center}
            \begin{tabularx}{\linewidth}{|l|X|X|}\hline
            \rowcolor{popblueL}
            \textbf{Critère} & \textbf{Indicateur} & \textbf{Cible / Seuil acceptable} \\ \hline
            Productivité & Rendement culture principale (t/ha) & $\ge$ Témoin sans arbres (ou $+10\%$ à 5 ans) \\ \hline
            Fertilité & Taux MO sol (\%), Stock C (t/ha), pH, N/P/K dispo & $\uparrow$ MO $> 2\%$ ; Stock C $\uparrow$ 0,5 t/ha/an \\ \hline
            Eau & Humidité sol profil (sondes), Rendement eau (kg/mm/ha) & $\uparrow$ Disponibilité eau utile \\ \hline
            Biodiversité & Richesse spécifique (arbres, auxiliaires, vers de terre) & $H'$ Shannon $\uparrow$ ; Vers de terre $> 100$ ind/m$^2$ \\ \hline
            Économique & Revenu net / ha / an (cultures + arbres + travail) & $>$ Revenu système courant ; TRI $> 15\%$ \\ \hline
            Social / Genre & Charge travail (femmes/hommes), Accès ressources, Sécurité foncière & Réduction corvée bois/ eau (femmes) ; Accord clan/chef terre \\ \hline
        \end{tabularx}
    \end{center}
    \end{enumerate}
    \textbf{Reflexe Probatoire :} Citez des \textbf{espèces locales} (\emph{Faidherbia albida}, \emph{Acacia nilotica}, \emph{Gliricidia sepium}, \emph{Leucaena leucocephala}, \emph{Moringa oleifera}, \emph{Vitellaria paradoxa}, \emph{Parkia biglobosa}, \emph{Milicia excelsa}, \emph{Albizia} sp., \emph{Erythrina} sp.). Distinguez \textbf{zone forestière} (ombrage cacao, forêt-jardin) et \textbf{zone soudano-sahélienne} (parc à karité/néré/faidherbia, haies anti-érosives). Parlez de \textbf{sécurisation foncière} (certificat foncier, convention locale) comme condition \textit{sine qua non} de l'investissement arbre (long terme).
\end{methode}

\begin{exoresolu}[Conception d'un système agroforestier pour la zone soudano-sahélienne (Nord-Cameroun, culture Coton/Maïs)]
    \textbf{Énoncé :} Exploitation familiale de 5 ha dans le Diamaré (Maroua). Sol : Vertisol (argile gonflante, fissures saisonnières) / Sol ferrugineux tropical lessivé sur pente douce (3\%). Pluviométrie : $900$ mm (juin-sept). Cultures : Coton (rente), Maïs/Sorgho/Niébé (vivrier). Problèmes : Érosion hydrique (ravines), déclin fertilité (MO $< 1\%$), pénurie bois énergie/fourrage (saison sèche 7 mois), conflits agriculteurs/éleveurs (divagation). Tenure : Droit coutumier (chef de terre), pas de titre.
    \textbf{Travail :} Proposez un dispositif agroforestier complet (espèces, disposition, itinéraire technique, indicateurs de suivi). Justifiez chaque choix par le contexte.
    \tcblower
    \textbf{Solution détaillée :}
    \begin{enumerate}
        \item \textbf{Stratégie globale :} \textbf{Intensification écologique} par l'arbre pour : (1) stopper érosion (haies contour), (2) restaurer fertilité (légumineuses), (3) produire bois/énergie/fourrage/fruits (revenus, sécurité alimentaire), (4) pacifier relations agriculteurs/éleveurs (fourrage, délimitation).
        \item \textbf{Dispositif spatial (Plan de masse 5 ha) :}
        \begin{itemize}
            \item \textbf{Haies anti-érosives en courbes de niveau (Trait principal) :} Espacement $15$ m (pente 3\% $\rightarrow$ $D = 15 / \sin \alpha \approx 15 / 0,03 \approx 500$ m théorique, mais praticabilité $\rightarrow$ 15-20 m). \textbf{Espèces :} \emph{Gliricidia sepium} (fixation N, recépage facile, fourrage, piquets) + \emph{Vetiveria zizanioides} (vetiver) en bas de haie (racines profondes 3-4 m, barrière physique infranchissable). \textbf{Densité :} 1 plant / 0,5 m linéaire $\rightarrow$ 20 000 plants/ha haies $\approx$ 1000 plants/ha parcelle.
            \item \textbf{Parc à \emph{Faidherbia albida} (Gao) en plein champ :} 20 arbres/ha (espacement $22 \times 22$ m). \textbf{Justification :} \textbf{Phénologie inversée} (feuilles saison sèche) $\rightarrow$ Ombrage bétail/fourrage (gousses/feuilles riches protéines) + Apport N (chute feuilles décomposition début pluies) + Pas de compétition cultures pluviales (nu à la culture) + Bois d'œuvre/énergie. \textbf{Compatible coton/maïs.}
            \item \textbf{Bordure périphérique (Haie vive / Brise-vent) :} \emph{Acacia nilotica} (épines = clôture anti-divagation) + \emph{Moringa oleifera} (feuilles nutritives humain/bétail, graines huile, croissance rapide) + \emph{Ziziphus mauritiana} (jujubier, fruits vendables). Largeur 5 m (2 rangs décalés).
            \item \textbf{Bas-fonds / Zones humides (si existant) :} \emph{Acacia seyal} / \emph{Mitragyna inermis} (stabilisation berges, nectar apicole).
        \end{itemize}
        \item \textbf{Itinéraire technique \& Gestion (Calendrier) :}
        \begin{center}
            \begin{tabularx}{\linewidth}{|l|X|X|}\hline
            \rowcolor{popblueL}
            \textbf{Opération} & \textbf{Période (Nord)} & \textbf{Détails techniques} \\ \hline
            Préparation plants & Mars-Avril (pépinière) & Sacs plastiques, substrat sol + fumier (1/1). \emph{Faidherbia} : scarification graines (eau bouillante 1 min). \emph{Gliricidia} : boutures 30 cm ou graines. \\ \hline
            Marquage / Traçage courbes niveau & Avril-Mai (avant pluies) & Niveau à eau / niveau à fil. Piquets tous les 15 m. \\ \hline
            Plantation & \textbf{Juin (1ères grosses pluies)} & Poquets $40 \times 40 \times 40$ cm. Apport 2 kg fumier décomposé + 100 g NPK/poquet. Arrosage si pause pluies $> 10$ j. \\ \hline
            Protection & Immédiat & \textbf{Pare-bêtes} (grillage / épines / bidons coupés) \textbf{OBLIGATOIRE} (pression broutage forte). Paillage paille/feuilles (conserver humidité). \\ \hline
            Entretien (Désherbage) & Juillet, Août, Septembre & Manuel au pied (1 m rayon). Éviter herbicide (jeunes plants sensibles). \\ \hline
            Taille de formation / Recépage & \textbf{Mars-Avril (Fin saison sèche)} & \emph{Gliricidia} : recépage 0,5-1 m (production bois/feuilles). \emph{Faidherbia} : élagage bas (formation tronc unique, enlever gourmands). \emph{Acacia nilotica} : taille haie (forme trapézoïdale base large). \\ \hline
            Récoltes & Tout au long & \emph{Gliricidia} : bois/feuilles (recépage). \emph{Faidherbia} : gousses (fourrage), bois mort. \emph{Moringa} : feuilles (cuisine), gousses. \emph{Ziziphus} : fruits (vente). \\ \hline
        \end{tabularx}
    \end{center}
        \item \textbf{Sécurisation foncière (Condition de réussite) :}
        \begin{itemize}
            \item Négociation avec Chef de terre / Conseil villageois : \textbf{Convention locale écrite} (acte notarié ou délibération conseil) garantissant les droits de l'exploitant sur les arbres plantés (droit de plantation, récolte, transmission héréditaire) pour 25-30 ans minimum.
            \item Demande \textbf{Certificat Foncier} (Loi 74/1, 74/2, 2005/019) si possible, ou attestation de non-objection du chef de terre pour investissement long terme.
        \end{itemize}
        \item \textbf{Tableau de bord de suivi (Indicateurs SMART - 5 ans) :}
        \begin{center}
            \begin{tabular}{|l|l|l|}\hline
            \rowcolor{popblueL}
            \textbf{Indicateur} & \textbf{Cible An 3} & \textbf{Cible An 5} \\ \hline
            Taux de reprise plants (haies + parc) & $> 80\%$ & $> 90\%$ (remplacement échecs an 1-2) \\ \hline
            Hauteur moyenne \emph{Faidherbia} / \emph{Gliricidia} & $> 2,5$ m / $> 3$ m & $> 4$ m / $> 4$ m (recépage) \\ \hline
            Rendement Maïs (t/ha) & $1,8$ (vs $1,2$ témoin) & $2,5$ (vs $1,0$ témoin sans arbre) \\ \hline
            Rendement Coton (kg/ha) & $1200$ (vs $1000$) & $1500$ (vs $800$) \\ \hline
            Taux MO sol (0-20 cm) & $1,2\%$ & $1,8\%$ \\ \hline
            Volume bois énergie produit (stères/ha/an) & $2$ (haies) & $5$ (haies + \emph{Faidherbia} bois mort) \\ \hline
            Fourrage produit (kg MS/ha/an saison sèche) & $500$ (\emph{Gliricidia} + \emph{Faidherbia}) & $1500$ \\ \hline
            Conflits agriculteurs/éleveurs déclarés & $\downarrow 50\%$ & $0$ (haie vive fonctionnelle) \\ \hline
            Revenu net exploitation (kFCFA/ha/an) & $+ 20\%$ vs référence & $+ 50\%$ vs référence \\ \hline
        \end{tabular}
    \end{center}
        \item \textbf{Analyse Coût-Bénéfice simplifiée (An 5) :}
        \begin{itemize}
            \item Investissement initial (plants, pare-bêtes, main d'œuvre plantation) : $\approx 150\,000$ FCFA/ha.
            \item Coût entretien/an (taille, désherbage, surveillance) : $\approx 25\,000$ FCFA/ha/an.
            \item Bénéfices An 5 : Gain rendement maïs/coton ($\approx +50\,000$ FCFA/ha/an), Bois énergie ($5$ st $\times 5000 = 25\,000$), Fourrage (économie achat $\approx 30\,000$), Fruits Moringa/Jujube ($\approx 10\,000$). Total $\approx 115\,000$ FCFA/ha/an.
            \item \textbf{TRI (Taux de Rentabilité Interne) $> 25\%$} $\rightarrow$ \textbf{Très rentable.}
        \end{itemize}
    \end{enumerate}
    \textbf{Conclusion :} Ce dispositif « Haies \emph{Gliricidia/Vetiver} + Parc \emph{Faidherbia} + Bordure \emph{Acacia/Moringa/Ziziphus} » répond simultanément aux 4 problèmes (érosion, fertilité, bois/énergie/fourrage, conflits) avec des espèces adaptées au vertisol/sol ferrugineux et au climat soudano-sahélien. La \textbf{sécurisation foncière conventionnelle} est le levier indispensable pour l'adoption paysanne.
\end{exoresolu}

\courssub{Savoir-faire complémentaire : Analyser un conflit d'usage de l'eau (GIRE - Gestion Intégrée des Ressources en Eau)}

\begin{methode}[Analyse et résolution d'un conflit d'usage de l'eau (Cadre GIRE Cameroun)]
    \textbf{Objectif :} Appliquer les principes de la GIRE (Loi 98/005, Décret 2001/165) à un cas concret de partage de l'eau.
    \begin{enumerate}
        \item \textbf{Identifier les usages en concurrence} (Matrice Usages $\times$ Acteurs) :
        \begin{itemize}
            \item \textbf{Usages :} Eau potable (AEP), Irrigation, Élevage (abreuvement), Industrie / Mines, Hydroélectricité, Pêche / Aquaculture, Navigation, Loisirs / Tourisme, Écosystèmes (débit réservé).
            \item \textbf{Acteurs :} SNEC/CAMWATER, MINEA, MINADER (ONAD, SODECOTON), MINEPIA, Industries (Alucam, brasseries, cimenteries), EDC/NHPC, Communes, Associations d'usagers (ASUREP, Groupements irrigants, Éleveurs), ONG, Populations riveraines.
        \end{itemize}
        \item \textbf{Quantifier la demande et la ressource disponible :}
        \begin{itemize}
            \item Ressource : Module annuel, module d'étiage ($Q_{min}$), débit de crue, qualité (physico-chimie, bactériologie).
            \item Demande : Volumes prélevés (m$^3$/an), consommés (non restitués), restitués (qualité dégradée ?), horaires (pointe/jour/nuit).
        \end{itemize}
        \item \textbf{Appliquer la hiérarchie des usages (Loi 98/005 Art. 14) :}
        \textbf{1. Eau potable / Salubrité} (Priorité absolue) $\rightarrow$ \textbf{2. Écosystèmes (Débit réservé)} $\rightarrow$ \textbf{3. Autres usages} (Irrigation, Industrie, Énergie, Loisirs) selon critères : efficacité économique, valeur sociale, antériorité droits, impact environnemental.
        \item \textbf{Identifier le(s) conflit(s) :} Pénurie quantitative (prélèvements $>$ ressource étiage), Dégradation qualitative (pollution amont $\rightarrow$ inutilisable aval), Conflit temporel (hydroélectricité éclusées vs pêche/irrigation), Conflit spatial (amont/aval, transfrontalier).
        \item \textbf{Proposer des solutions GIRE (Boîte à outils) :}
        \begin{itemize}
            \item \textbf{Institutionnelles :} Comité de Bassin (Loi 98/005 Art. 38), Contrats d'objectif / Conventions d'usage, Police de l'eau (MINEA), Répartition par quotas (Plan d'Allocation de l'Eau - PAE).
            \item \textbf{Techniques :} Économie d'eau (goutte-à-goutte, réutilisation eaux usées traitées, réduction fuites réseau), Stockage (barrages collinaires, recharge nappe artificielle), Traitement pollution (lagunage, stations épuration).
            \item \textbf{Économiques :} Tarification eau (redevance prélèvement, redevance pollution) $\rightarrow$ incitation économie / pollueur-payeur. Paiement pour Services Environnementaux (PSE) amont (reboisement) $\rightarrow$ qualité/quantité aval.
            \item \textbf{Information / Participation :} Système d'Information sur l'Eau (SIE), Concertation continue, Gestion des crises (Comité de crise sécheresse).
        \end{itemize}
        \item \textbf{Rédiger l'avis / le plan d'action} : Prioriser les mesures « sans regret » (économie, protection zones de recharge), chiffrer les coûts/bénéfices, définir les responsabilités.
    \end{enumerate}
    \textbf{Reflexe Probatoire :} La \textbf{Loi 98/005} est la référence absolue. Citez l'\textbf{Art. 14} (Hiérarchie usages + Débit réservé), \textbf{Art. 38} (Comités de bassin), \textbf{Art. 44} (Police de l'eau). Distinguez \textbf{prélèvement} (volume pompé) et \textbf{consommation nette} (volume non restitué). Le \textbf{débit réservé} n'est pas un « reste » mais un \textbf{droit de l'environnement} à garantir en priorité après l'AEP.
\end{methode}

\begin{exoresolu}[Conflit d'usage sur le fleuve Bénoué (Région Nord) - Sécheresse 2025]
    \textbf{Énoncé :} Contexte : Année sèche exceptionnelle (pluviométrie $-30\%$). Fleuve Bénoué (affluent Niger). Usages : AEP Garoua (SNEC, $30\,000$ m$^3$/j), Irrigation riz SEMRY (Lagdo, $15$ m$^3$/s pointe), Barrage Lagdo (EDC, turbinage $200$ m$^3$/s pointe, éclusées), Élevage transhumant (abreuvement $500\,000$ TUV), Pêche artisanale ($5\,000$ pêcheurs), Écosystème (parc Bénoué, hippopotames, lamantins). Débit étiage observé (avril 2025) : $40$ m$^3$/s (vs moyenne $120$ m$^3$/s). Qualité : Turbidité $> 1000$ NTU (érosion amont), $[\ce{NO3-}]$ hausse.
    \textbf{Travail :} En tant qu'expert du Comité de Bassin Bénoué, proposez un \textbf{Plan d'Allocation de l'Eau (PAE) de crise} pour l'étiage 2025 (avril-juin). Hiérarchisez, quantifiez les quotas, proposez 3 mesures d'urgence et 2 mesures structurelles.
    \tcblower
    \textbf{Solution détaillée :}
    \begin{enumerate}
        \item \textbf{Bilan Ressource / Demande (Étiage critique - Avril) :}
        \begin{itemize}
            \item Ressource disponible (Débit étiage) : $Q_{dispo} = 40$ m$^3$/s = $3\,456\,000$ m$^3$/j.
            \item \textbf{Débit réservé écologique (Vie aquatique, parc, dilution pollution) :} Minimum scientifique $20\%$ module interannuel $\approx 24$ m$^3$/s. \textbf{Seuil critique survie faune (hippos, poissons) : $15$ m$^3$/s.} $\rightarrow$ \textbf{Fixé à $Q_{res} = 15$ m$^3$/s} (compromis survie / partage).
            \item \textbf{Ressource partageable :} $Q_{part} = 40 - 15 = \mathbf{25 \text{ m}^3/\text{s} = 2\,160\,000 \text{ m}^3/\text{j}}$.
        \end{itemize}
        \item \textbf{Demandes exprimées (Pointe journalière) :}
        \begin{itemize}
            \item AEP Garoua : $30\,000$ m$^3$/j = $0,35$ m$^3$/s (Priorité 1 - Incompressible).
            \item Irrigation SEMRY : $15$ m$^3$/s (Pointe) $\rightarrow$ \textbf{Demande $>$ Ressource partageable !} (Nécessite restriction).
            \item Barrage Lagdo (Hydro) : $200$ m$^3$/s (Turbinage) $\rightarrow$ \textbf{Impossible} (Doit turbinier sur stock réservoir, pas fil de l'eau). Éclusées interdites.
            \item Élevage / Pêche / Autre : Estimé $< 2$ m$^3$/s (diffus).
        \end{itemize}
        \item \textbf{Plan d'Allocation de Crise (Quotas journaliers max) :}
        \begin{center}
            \begin{tabularx}{\linewidth}{|l|l|l|X|}\hline
            \rowcolor{popblueL}
            \textbf{Usage (Priorité Loi 98/005)} & \textbf{Quota attribué (m$^3$/s)} & \textbf{Volume journalier (m$^3$/j)} & \textbf{Conditions / Restrictions} \\ \hline
            1. AEP Garoua (Potable) & \textbf{0,35} (Garanti 100\%) & 30\,000 & Surveillance qualité (turbidité, bactéries). Traitement renforcé. \\ \hline
            2. Débit Réservé (Écosystème) & \textbf{15,00} (Garanti 100\%) & 1\,296\,000 & \textbf{Non négociable.} Surveillance continue niveau/qualité aval barrage. \\ \hline
            3. Irrigation SEMRY (Riz) & \textbf{6,00} (Restriction 60\%) & 518\,400 & \textbf{Arrêt irrigation 12h/j} (nuit). Passage goutte-à-goutte / aspersion $\rightarrow$ économie 30\%. Priorité riz en floraison. \\ \hline
            4. Barrage Lagdo (Hydro) & \textbf{3,00} (Turbinage fil de l'eau) & 259\,200 & \textbf{Interdiction éclusées.} Turbinage continu « run-of-river » pour lisser débit. Production $\downarrow$ acceptée. \\ \hline
            5. Élevage / Pêche / Divers & \textbf{0,65} (Partage résiduel) & 56\,160 & Points d'eau aménagés (pompage solaire) hors lit mineur. \\ \hline
            \textbf{TOTAL} & \textbf{25,00} & \textbf{2\,160\,000} & \textbf{Équilibre respecté.} \\ \hline
        \end{tabularx}
    \end{center}
        \item \textbf{3 Mesures d'urgence (Avril-Juin 2025) :}
        \begin{enumerate}
            \item \textbf{Arrêté Préfectoral de Restriction d'Usage (Loi 98/005 Art. 44) :} Notification formelle quotas ci-dessus. Sanctions (amendes, suspension pompage) pour dépassement. Brigade de police de l'eau (MINEA + Forces de l'ordre) déployée sur points de prélèvement SEMRY / Lagdo.
            \item \textbf{Opération « Eau pour le Bétail » :} Déploiement de $20$ unités mobiles de pompage solaire + abreuvoirs (capacité $10$ m$^3$/j chacune) sur points stratégiques transhumance (loin du fleuve, zones pâturages résiduels). Budget urgence : $50$ M FCFA (Fonds National de Gestion des Risques / MINADER/MINEPIA).
            \item \textbf{Surveillance qualitative renforcée + Alerte :} 3 stations automatiques (Amont Garoua, Aval Lagdo, Parc Bénoué) : Niveau, $T^\circ$, $O_2$, Turbidité, $\ce{NO3-}$. Seuil alerte $O_2 < 3$ mg/L $\rightarrow$ Aération d'urgence / Arrêt rejets industriels. Information quotidienne radio locale (fulfulde, français) sur niveau fleuve / qualité / points d'eau bétail.
        \end{enumerate}
        \item \textbf{2 Mesures structurelles (Prévention crises futures) :}
        \begin{enumerate}
            \item \textbf{Plan d'Allocation de l'Eau (PAE) Pluriannuel du Bassin Bénoué :} Modélisation hydrologique (WEAP / SWAT) intégrant changement climatique. Définition de \textbf{règles de partage dynamiques} (quotas = f(Niveau barrage Lagdo, Prévision pluviométrie, Humidité sol). \textbf{Institutionnalisation Comité de Bassin Bénoué} (Loi 98/005 Art. 38) avec pouvoir de décision contraignant (représentants Etat, Régions, Communes, Usagers, ONG, Transfrontalier Niger Basin Authority).
            \item \textbf{Programme « Bénoué Vert » (Restauration bassin versant amont) :} Reboisement / Régénération Naturelle Assistée (RNA) sur $50\,000$ ha zones sources / versants dégradés (Régions Nord, Adamaoua, Nord-Ouest, RCA). \textbf{Objectif :} Augmenter infiltration $\rightarrow$ Recharge nappe / Alimentation base-flow étiage $\rightarrow$ Réduire érosion (turbidité) $\rightarrow$ Améliorer qualité eau. Financement : PSE (Payeurs : SNEC Garoua, EDC Lagdo, SEMRY, Brasseries) $\rightarrow$ Bénéficiaires : Communautés amont (récolte PFNL, carbone, bois). \textbf{Cible :} $+ 10\%$ débit étiage à horizon 15 ans.
        \end{enumerate}
    \end{enumerate}
    \textbf{Conclusion :} La GIRE impose de \textbf{quantifier, hiérarchiser (Loi 98/005), partager équitablement (PAE) et investir dans le capital naturel (bassin versant)}. La crise 2025 montre l'urgence d'un Comité de Bassin opérationnel et d'un PAE dynamique, au-delà de la gestion de crise réactive.
\end{exoresolu}

\courssub{Savoir-faire complémentaire : Valoriser les services écosystémiques (PSE - Paiement pour Services Environnementaux)}

\begin{methode}[Montage d'un mécanisme de Paiement pour Services Environnementaux (PSE) local]
    \textbf{Objectif :} Transformer la valeur économique des services écosystémiques (SE) en incitations financières pour les gestionnaires de terres (agriculteurs, communautés), selon l'approche « Beneficiary Pays » ou « Polluter Pays ».
    \begin{enumerate}
        \item \textbf{Identifier le Service Écosystémique (SE) ciblé et sa valeur :}
        \begin{itemize}
            \item \textbf{SE de Régulation :} Régulation hydrique (recharge nappe, prévention crues), Séquestration carbone, Contrôle érosion, Pollinisation, Régulation maladies.
            \item \textbf{SE d'Approvisionnement :} Eau de qualité, Bois, PFNL, Ressources génétiques.
            \item \textbf{SE Culturels :} Tourisme, Sacralité, Éducation, Paysage.
            \item \textbf{Valorisation économique :} Coût de remplacement (usine traitement eau vs forêt), Coût évité (dégâts crues), Marché carbone (crédits), Enquêtes consentement à payer (CVP).
        \end{itemize}
        \item \textbf{Identifier Acheteurs (Bénéficiaires) et Vendeurs (Fournisseurs) :}
        \begin{itemize}
            \item \textbf{Acheteurs :} Entreprises eau (SNEC, CAMWATER), Hydroélectriens (EDC, NHPC), Irrigants (SEMRY, SODECOTON), Industries, Communes, État (Budget / Fonds Environnement), ONG / Bailleurs (Banque Mondiale, AFD, GEF), Marché volontaire carbone.
            \item \textbf{Vendeurs :} Communautés villageoises (Terroirs), Exploitants forestiers (UFA), Agriculteurs (parcelles), Communes (Forêts communales), ONG gestion aires protégées.
        \end{itemize}
        \item \textbf{Définir l'Action de gestion (Conditionnalité) :} Qu'est-ce que le vendeur *doit faire* en plus du *business as usual* ? (Ex. : Reboisement $X$ ha/an, Arrêt brûlis, Agriculture de conservation, Surveillance anti-braconnage, Maintien jachère longue). \textbf{Additionnalité} prouvée.
        \item \textbf{Négocier le Contrat PSE :}
        \begin{itemize}
            \item \textbf{Prix :} Basé sur coût d'opportunité (revenu perdu agriculture) + coûts transaction + prime incitation. Ou basé sur valeur SE (coût évité).
            \item \textbf{Modalités :} Paiement à l'hectare/an, ou au résultat (ex. t $\ce{CO2}$ séquestré vérifié, m$^3$ eau infiltrée), ou forfaitaire.
            \item \textbf{Durée :} Minimum 5-10 ans (cycle arbre/climat).
            \item \textbf{Conditionnalité (MRV - Monitoring, Reporting, Verification) :} Indicateurs simples, vérifiables localement + tiers indépendant (satellite, terrain).
            \item \textbf{Gouvernance :} Comité de pilotage (Acheteurs, Vendeurs, Etat, Facilitateur ONG). Fonds fiduciaire / Compte séquestre.
        \end{itemize}
        \item \textbf{Assurer la pérennité :} Intégration dans plans communaux (PCD), sécurisation foncière vendeurs, renforcement capacités locales, diversification revenus (pas dépendance unique PSE).
    \end{enumerate}
    \textbf{Reflexe Probatoire :} Citez des exemples camerounais réels ou plausibles : \textbf{Projet PSE Lac Lagdo / Bénoué} (reboisement amont vs barrage), \textbf{Projet Carbone Mont Cameroun / Bakossi} (REDD+), \textbf{PSE Eau Yaoundé (Mfoundi)} (reboisement bassin versant vs traitement eau). Distinguez \textbf{PSE} (contractuel, volontaire, conditionnel) et \textbf{Taxes/Revenus} (obligatoires, non conditionnels à une action spécifique).
\end{methode}

\begin{exoresolu}[Conception d'un PSE « Eau » pour le bassin versant de la Mfoundi (Yaoundé)]
    \textbf{Énoncé :} La ville de Yaoundé (3,5 M hab.) puise $60\%$ de son eau brute dans le bassin de la Mfoundi (barrage d'Akomnyada, usine d'Etoudi). La déforestation / agriculture sur brûlis / urbanisation anarchique amont (communes de Mbankomo, Nkolafamba, Soa) dégradent la qualité (turbidité $> 500$ NTU, coliformes, pesticides) et la régularité du débit. Coût traitement eau (Etoudi) : $150$ FCFA/m$^3$ (produits chimiques, énergie, boues). Coût estimé si qualité « forêt intacte » : $80$ FCFA/m$^3$. Volume traité : $200\,000$ m$^3$/j. Surface bassin : $120\,000$ ha. Population rurale amont : $50\,000$ hab. Revenus moyens : $200\,000$ FCFA/an/ménage (agriculture itinérante, charbon).
    \textbf{Travail :} Montez un projet PSE « Eau Yaoundé ». Calculez la valeur du service, proposez le contrat (acheteurs, vendeurs, actions, prix, MRV, gouvernance).
    \tcblower
    \textbf{Solution détaillée :}
    \begin{enumerate}
        \item \textbf{Valorisation économique du service « Régulation qualité/quantité eau » :}
        \begin{itemize}
            \item Économie potentielle traitement : $(150 - 80)$ FCFA/m$^3$ $\times$ $200\,000$ m$^3$/j $\times$ $365$ j = $70 \times 73 \times 10^6$ = \textbf{5,11 Mrd FCFA/an}.
            \item Coût évité infrastructures nouvelles (barrage amont, transfert) : estimé $> 50$ Mrd FCFA (actualisé).
            \item \textbf{Valeur de référence du SE :} \textbf{5 Mrd FCFA/an} (économie traitement seule, conservative).
        \end{itemize}
        \item \textbf{Acteurs du contrat :}
        \begin{itemize}
            \item \textbf{Acheteurs (Payeurs) :} \textbf{CAMWATER / SNEC} (bénéficiaires directs économie traitement), \textbf{Communes de Yaoundé} (bénéficiaires santé/qualité vie), \textbf{MINEA / MINEPDED} (subvention incitative), \textbf{Bailleurs} (AFD, BEI - prêt/bonification). \textit{Objectif : mobiliser 1 à 2 Mrd FCFA/an (20-40\% valeur SE).}
            \item \textbf{Vendeurs (Fournisseurs) :} \textbf{Communautés villageoises} (groupements CVD - Comités Villageois de Développement) des communes amont (Mbankomo, Nkolafamba, Soa, etc.). \textbf{Communes rurales} (portage foncier, planification). \textbf{ONG facilitatrices} (WWF, CIFOR, ONG locales) pour montage/MRV.
        \end{itemize}
        \item \textbf{Actions conditionnelles (Additionnalité) - « Paquet PSE » par hectare/an :}
        \begin{center}
            \begin{tabularx}{\linewidth}{|l|X|l|}\hline
            \rowcolor{popblueL}
            \textbf{Action} & \textbf{Description technique} & \textbf{Indicateur MRV (Annuel)} \\ \hline
            \textbf{Reboisement / RNA} & Plants espèces locales (bois d'œuvre, PFNL, fertilité) sur terres dégradées / jachères courtes. RNA (Régénération Naturelle Assistée) favorisée. & Ha reboisés / RNA efectifs (survie $> 80\%$ an 3). \\ \hline
            \textbf{Agroforesterie / AC} & Haies \emph{Gliricidia/Leucaena} contour + \emph{Faidherbia} champ. Agriculture de Conservation (SCV : Semis Direct, Couverture, Rotation). & Ha en AC / Agroforesterie. Taux couverture sol $> 30\%$ permanence. \\ \hline
            \textbf{Arrêt brûlis / Feux contrôlés} & Zéro feu non maîtrisé. Plans de gestion feux villageois. & Nb feux détectés (satellite MODIS/VIIRS) = 0 sur zone contractée. \\ \hline
            \textbf{Protection berges / Zones humides} & Bande 10-20 m non cultivée, revégétalisée (\emph{Raphia}, \emph{Alchornea}). & Linéaire berges protégées (km). \\ \hline
            \textbf{Assainissement / Gestion déchets} & Latrines améliorées, compostage, collecte plastiques (réduction pollution fécale/chimique). & Nb latrines / foyers ; Volume déchets collectés. \\ \hline
        \end{tabularx}
    \end{center}
        \item \textbf{Calcul du prix / ha / an (Approche Coût d'Opportunité + Incitation) :}
        \begin{itemize}
            \item Revenu agriculture itinérante / charbon (net) : $\approx 150\,000$ FCFA/ha/an (main d'œuvre familiale non rémunérée).
            \item Coûts mise en œuvre actions (plants, formation, surveillance, entretien an 1-3) : $\approx 100\,000$ FCFA/ha/an (moyenné).
            \item Prime incitation (changement pratique, risque) : $50\,000$ FCFA/ha/an.
            \item \textbf{Prix PSE proposé :} \textbf{300\,000 FCFA/ha/an} (versé conditionnellement aux résultats MRV).
        \end{itemize}
        \item \textbf{Ciblage spatial et budget prévisionnel (5 ans) :}
        \begin{itemize}
            \item Zone prioritaire : $20\,000$ ha (têtes de bassin, versants fortes pentes, berges) sur $120\,000$ ha.
            \item Objectif : $4\,000$ ha/an mis sous contrat (reboisement/AC) $\rightarrow$ $20\,000$ ha an 5.
            \item Budget annuel moyen : $4\,000$ ha $\times$ $300\,000$ = \textbf{1,2 Mrd FCFA/an}.
            \item Coûts transaction / MRV / Facilitation (15\%) : $180$ M FCFA/an.
            \item \textbf{Total annuel :} \textbf{1,38 Mrd FCFA/an}.
            \item \textbf{Rentabilité pour CAMWATER :} Économie traitement $5,11$ Mrd FCFA/an $>$ Coût PSE $1,38$ Mrd FCFA/an $\rightarrow$ \textbf{Gain net $3,7$ Mrd FCFA/an}. Ratio Bénéfice/Coût $\approx 3,7$.
        \end{itemize}
        \item \textbf{Dispositif MRV (Monitoring, Reporting, Verification) :}
        \begin{itemize}
            \item \textbf{Monitoring continu (Communautés + ONG) :} Suivi survie plants, respect haies, absence feux (télédétection Sentinel-2 / alertes Global Forest Watch), qualité eau points fixes (kits terrain).
            \item \textbf{Reporting semestriel :} Rapports CVD $\rightarrow$ Commune $\rightarrow$ Comité Pilotage.
            \item \textbf{Verification annuelle (Tierce partie indépendante) :} Bureau d'études / Université (ENSET, IRAD) : Échantillonnage terrain $10\%$ ha, analyse images satellite, mesures qualité eau (turbidité, $\ce{NO3-}$, coliformes) exutoires sous-bassins. \textbf{Paiement déclenché seulement après validation rapport vérification.}
        \end{itemize}
        \item \textbf{Gouvernance du Fonds PSE « Mfoundi Vert » :}
        \begin{itemize}
            \item \textbf{Comité de Pilotage (CP) :} Présidence : Gouverneur Région Centre / MINEA. Membres : CAMWATER, Communes Yaoundé + Rurales, MINEPDED, MINADER, Représentants CVD, ONG Facilitatrice, Bailleurs. \textit{Décide : validation plans annuels, paiements, arbitrage conflits.}
            \item \textbf{Unité de Gestion (UG) :} ONG facilitatrice (ex. WWF-Cameroun) sous convention. \textit{Animation terrain, appui technique CVD, pré-MRV, secrétariat CP.}
            \item \textbf{Fonds fiduciaire :} Compte séquestre (Banque commerciale / BEAC) alimenté par contributeurs (CAMWATER, Communes, Subventions). Paiements directs vers comptes CVD / Communautés (mobile money / banque).
            \item \textbf{Sécurisation foncière :} Délibérations communes / chefs traditionnels actant « Zone PSE » = interdiction défrichement nouveau, priorité actions contrat. Intégration dans PCD (Plans Communaux Développement).
        \end{itemize}
    \end{enumerate}
    \textbf{Conclusion :} Ce PSE « Mfoundi Vert » est \textbf{efficace économiquement} (gain net 3,7 Mrd/an pour l'opérateur eau), \textbf{équitable} (revenu significatif $300\,000$ FCFA/ha/an pour ruraux pauvres vs $150\,000$ FCFA agriculture dégradante), \textbf{écologiquement additionnel} (restauration $20\,000$ ha bassin versant critique). La clé de succès : \textbf{MRV rigoureux indépendant} et \textbf{gouvernance multipartite transparente} (Fonds fiduciaire), inscrite dans la durée (10 ans minimum).
\end{exoresolu}

\begin{attention}[Les 5 erreurs qui coûtent des points au Probatoire]
    \begin{enumerate}
        \item \textbf{Confondre « Ressource naturelle » et « Service écosystémique ».} La ressource est le stock/flux physique (eau, bois, sol, espèce). Le service est le \textbf{bénéfice} que l'homme en tire (eau potable, régulation crues, pollinisation, valeur culturelle). \textit{Ex. : « La forêt est un service » $\rightarrow$ FAUX. La forêt \emph{fournit} des services.}
        \item \textbf{Oublier la hiérarchie légale des usages de l'eau (Loi 98/005 Art. 14).} Ne pas placer \textbf{Eau Potable} et \textbf{Débit Réservé (Écosystèmes)} en priorité absolue avant Irrigation/Industrie/Énergie. Proposer un partage « équitable » sans respecter cette hiérarchie = erreur juridique majeure.
        \item \textbf{Proposer des mesures « magiques » sans ancrage local (acteurs, foncier, coûts).} Ex. : « Reboiser 100 000 ha » sans dire : qui plante ? sur quel foncier (coutumier/titre) ? avec quels plants (pépinières locales) ? qui finance ? qui surveille ? \textbf{Toute mesure doit citer : Acteur responsable, Indicateur SMART, Cadre juridique camerounais, Source de financement.}
        \item \textbf{Ne pas quantifier / Ne pas utiliser les données du document.} Réponse qualitative vague (« pollution importante ») au lieu de « $[\ce{NO3-}] = 85$ mg/L $>$ norme 50 mg/L ». Au Probatoire, \textbf{chiffre = point}. Lire les axes des graphiques, les légendes des cartes, les tableaux fournis.
        \item \textbf{Confondre « Déforestation » et « Dégradation forestière ».} Déforestation = changement d'affectation (Forêt $\rightarrow$ Non-forêt : agriculture, urbain). Dégradation = perte de capacité productive/écologique \textbf{sans} changement d'affectation (exploitation sélective excessive, feux, prélèvement bois énergie, fragmentation). Les causes, indicateurs (taux couverture, biomasse, biodiversité) et solutions diffèrent.
    \end{enumerate}
\end{attention}

\newpage

\coursec{Exercices}

\courssub{Je vérifie mes connaissances}

\begin{exercice}[QCM : Ressources et pressions]
Parmi les affirmations suivantes, une seule est correcte. La recopier en justifiant brièvement votre choix.
\begin{enumerate}
\item L'eau douce est une ressource inépuisable à l'échelle humaine car le cycle de l'eau la renouvelle en permanence.
\item La déforestation pour l'agriculture sur brûlis augmente durablement la fertilité des sols tropicaux par l'apport de cendres.
\item La biodiversité désigne uniquement le nombre d'espèces animales présentes dans un milieu donné.
\item L'exploitation forestière sélective, accompagnée de reboisement et de respect des cycles de rotation, illustre une gestion raisonnée de la ressource forestière.
\end{enumerate}
\end{exercice}

\begin{exercice}[Vrai / Faux justifié]
Dire si chacune des affirmations suivantes est vraie (V) ou fausse (F) en justifiant votre réponse par une phrase précise.
\begin{enumerate}
\item L'érosion hydrique des sols est accélérée par le labour en courbes de niveau.
\item Le braconnage dans le parc national de Waza ne menace que les grandes espèces emblématiques (éléphants, lions) et non l'équilibre de l'écosystème.
\item Le développement durable repose sur trois piliers indissociables : l'efficacité économique, l'équité sociale et la préservation de l'environnement.
\item Les ressources minières (fer, bauxite, or) sont des ressources renouvelables si l'on attend assez longtemps.
\end{enumerate}
\end{exercice}

\begin{exercice}[Définitions contextualisées]
Donner une définition précise et rigoureuse de chacun des termes suivants en l'illustrant par un exemple concret pris au Cameroun :
\begin{enumerate}
\item Ressource naturelle renouvelable.
\item Gestion raisonnée (ou durable) d'une ressource.
\item Érosion des sols.
\item Aire protégée.
\end{enumerate}
\end{exercice}

\begin{exercice}[Calcul direct : Taux de déforestation]
Dans une zone forestière de l'Est-Cameroun, la superficie forestière était de \SI{120 000}{ha} en 2010. En 2020, elle n'est plus que de \SI{96 000}{ha}.
\begin{enumerate}
\item Calculer la superficie totale perdue sur cette période.
\item Calculer le taux moyen annuel de déforestation (en \% par an) en supposant une perte linéaire.
\item Exprimer cette perte moyenne annuelle en hectares par an.
\end{enumerate}
\end{exercice}

\courssub{Je m'entraîne}

\begin{exercice}[Analyse de document : Exploitation forestière dans l'Est-Cameroun]
On donne le tableau ci-dessous reprenant les données d'une concession forestière sur 5 ans.

\begin{center}
\begin{tabularx}{\linewidth}{|c|X|X|X|}\hline
\textbf{Année} & \textbf{Superficie exploitée (ha)} & \textbf{Superficie reboisée (ha)} & \textbf{Volume de bois prélevé (m$^3$)} \\ \hline
2018 & 2 500 & 200 & 45 000 \\ \hline
2019 & 2 800 & 300 & 50 000 \\ \hline
2020 & 3 000 & 250 & 55 000 \\ \hline
2021 & 3 200 & 400 & 58 000 \\ \hline
2022 & 3 500 & 350 & 62 000 \\ \hline
\end{tabularx}
\end{center}

\begin{enumerate}
\item Calculer le solde net de superficie forestière (exploitée $-$ reboisée) pour chaque année et le solde cumulé sur la période.
\item Tracer, sur un même graphique, l'évolution de la superficie exploitée et de la superficie reboisée (utiliser un papier millimétré ou un logiciel).
\item En vous basant sur ces résultats, la gestion de cette concession vous paraît-elle durable ? Justifiez en citant deux critères de la gestion durable non respectés.
\item Proposer deux mesures concrètes pour améliorer la durabilité de cette exploitation.
\end{enumerate}
\end{exercice}

\begin{exercice}[Construction de courbe : Pêche dans le fleuve Sanaga]
Les captures annuelles de poisson (en tonnes) au niveau de la retenue de la Lagdo et du cours aval de la Sanaga ont été relevées sur 15 ans.

\begin{center}
\begin{tabular}{|c|c|c|c|c|c|c|c|c|}\hline
\textbf{Année} & 2008 & 2009 & 2010 & 2011 & 2012 & 2013 & 2014 & 2015 \\ \hline
\textbf{Captures (t)} & 12 500 & 12 200 & 11 800 & 11 000 & 10 200 & 9 500 & 8 900 & 8 200 \\ \hline
\end{tabular}
\end{center}
\begin{center}
\begin{tabular}{|c|c|c|c|c|c|c|c|}\hline
\textbf{Année} & 2016 & 2017 & 2018 & 2019 & 2020 & 2021 & 2022 \\ \hline
\textbf{Captures (t)} & 7 800 & 7 300 & 7 000 & 6 800 & 6 500 & 6 200 & 6 000 \\ \hline
\end{tabular}
\end{center}

\begin{enumerate}
\item Tracer la courbe d'évolution des captures en fonction du temps (années en abscisses, captures en ordonnées).
\item Décrire l'allure de la courbe et calculer le pourcentage de baisse des captures entre 2008 et 2022.
\item Conclure sur l'état de la ressource halieutique. Parle-t-on de surexploitation ? Justifiez.
\item Proposer \textbf{deux} mesures de gestion durable applicables localement (ex: période de fermeture, taille minimale de capture, quotas, zones de frayères protégées) et expliquer le principe de chacune.
\end{enumerate}
\end{exercice}

\begin{exercice}[Étude de cas : Gestion de l'eau en zone urbaine - Yaoundé]
La ville de Yaoundé fait face à des pénuries d'eau récurrentes en saison sèche. Les causes principales sont : la croissance démographique rapide, la pollution des cours d'eau (Mfoundi, Mefou) par les déchets solides et liquides, la déforestation des bassins versants (collines de Nkolbisson, Nkolandom) et les fuites sur le réseau de distribution (rendement estimé à 60\%).
\begin{enumerate}
\item Identifier, dans ce texte, \textbf{quatre} pressions humaines distinctes exercées sur la ressource en eau.
\item Pour chaque pression identifiée, proposer une mesure technique ou organisationnelle de réduction de son impact.
\item Expliquer pourquoi la protection des bassins versants (reboisement, interdiction de l'habitat sur les versants raides) est une mesure plus durable que l'augmentation seule de la capacité de traitement de l'eau.
\end{enumerate}
\end{exercice}

\begin{exercice}[Synthèse rédactionnelle : Le brûlis et les alternatives]
Un chef de village de la région de l'Adamaoua pratique l'agriculture itinérante sur brûlis depuis des décennies. Il constate que les rendements baissent et que les jachères ne se reforment plus. Il vous demande conseil.
Rédiger un texte argumenté d'environ 15 lignes (environ 150 mots) dans lequel vous :
\begin{enumerate}
\item Expliquez les mécanismes par lesquels le brûlis systématique et répété dégrade la fertilité du sol (perte de matière organique, destruction de la faune du sol, érosion, lessivage).
\item Présentez \textbf{trois} alternatives agroécologiques adaptées au contexte (agroforesterie, compostage, jachère améliorée / plantes de couverture, cultures associées) en expliquant brièvement le principe de chacune.
\item Concluez sur l'intérêt économique et écologique de la transition pour sa communauté.
\end{enumerate}
\end{exercice}

\courssub{Je me prépare au Probatoire}

\begin{exercice}[Sujet type Probatoire : Le Parc National de Waza face aux changements globaux]
Le Parc National de Waza (Extrême-Nord, \SI{170 000}{ha}), classé réserve de biosphère UNESCO, subit de fortes pressions. Les données suivantes sont disponibles :
\begin{itemize}
\item \textbf{Climat :} pluviométrie moyenne passée de \SI{700}{mm/an} (années 70) à \SI{550}{mm/an} (2020) ; augmentation de la fréquence des sécheresses sévères.
\item \textbf{Hydrologie :} assèchement quasi total de la plaine d'inondation de Waza (zone de pâture clé) suite à la construction du barrage de Maga (amont) et à l'envasement du fleuve Logone.
\item \textbf{Activités humaines :} braconnage intensif (éléphants, kob, girafes) ; pression pastorale transhumante (zébus) en saison sèche dans le parc ; agriculture sur les marges (riziculture irriguée).
\item \textbf{Données fauniques :} effectifs estimés de l'éléphant : 1 200 (1980) $\to$ 450 (2010) $\to$ 200 (2022) ; Kob de Buffon : 30 000 (1980) $\to$ 5 000 (2022).
\end{itemize}

\begin{enumerate}
\item \textbf{Analyse des pressions (6 pts) :}
\begin{enumerate}
\item Classer les pressions citées en deux catégories : « Pressions climatiques/environnementales » et « Pressions anthropiques directes ».
\item Expliquer le lien de cause à effet entre la construction du barrage de Maga et la baisse des effectifs du Kob de Buffon.
\item Pourquoi le braconnage de l'éléphant a-t-il un impact disproportionné sur la structure de la savane (espèce clé / ingénieur écosystémique) ?
\end{enumerate}
\item \textbf{Évaluation de la durabilité (4 pts) :}
La gestion actuelle repose sur la surveillance par des écogardes (effectifs insuffisants) et l'interdiction totale d'accès aux populations riveraines. Discuter les limites de cette approche « forteresse » (exclusion sociale, conflits, coûts) au regard des trois piliers du développement durable.
\item \textbf{Proposition de plan d'action intégré (10 pts) :}
Vous êtes consultant pour le MINFOF. Proposer un plan d'action structuré en \textbf{trois axes stratégiques} associant \textbf{populations locales, État et écotourisme}. Pour chaque axe, donner \textbf{deux actions concrètes, mesurables et datées} (ex: « Créer 5 comités de gestion villageoise d'ici 2025 », « Former 20 jeunes guides éco-touristiques d'ici 2026 », « Réhabiliter 3 points d'eau artificiels dans le parc d'ici 2025 »).
\end{enumerate}
\end{exercice}

\begin{exercice}[Sujet type Probatoire : Gestion intégrée du bassin versant de la Mefou (Alimentation en eau de Yaoundé)]
Le barrage de la Mefou (mis en service 2024, capacité \SI{30 000}{m^3/jour}) alimente Yaoundé. Son bassin versant (\SI{350}{km^2}) est occupé par : agriculture (40\%), habitat non planifié (25\%), savane/brousse (20\%), forêt résiduelle (15\%).
Une étude de la qualité de l'eau brute à l'entrée de l'usine de traitement montre : Turbidité moyenne = \SI{150}{NTU} (norme < 50 NTU) ; Teneur en nitrates = \SI{45}{mg/L} (norme OMS < 50 mg/L, mais tendance à la hausse) ; Coliformes fécaux = \SI{2 500}{UFC/100 mL}.

\begin{enumerate}
\item \textbf{Diagnostic environnemental (6 pts) :}
\begin{enumerate}
\item Relier chaque usage du sol du bassin versant à la pollution observée (turbidité, nitrates, bactéries). Justifier les mécanismes physiques/chimiques/biologiques.
\item Calculer le coût annuel supplémentaire de traitement si la turbidité impose une surdose de coagulant (sulfate d'alumine) de \SI{15}{mg/L} pour \SI{30 000}{m^3/jour}, le coagulant coûtant \SI{300}{FCFA/kg}. (Indication : masse journalière = Volume $\times$ surdose ; coût annuel = masse journalière $\times$ 365 $\times$ prix/kg).
\end{enumerate}
\item \textbf{Stratégie de gestion raisonnée (8 pts) :}
Le ministère de l'Eau (MINEE) souhaite mettre en place un « Plan de Gestion Intégrée des Ressources en Eau » (PGIRE) pour ce bassin.
\begin{enumerate}
\item Définir le concept de PGIRE.
\item Proposer \textbf{quatre} mesures concrètes de protection du bassin versant (ex: reboisement berges, fossés d'infiltration, latrines écologiques, zonage agricole, paiement pour services environnementaux). Pour chaque mesure, préciser : l'acteur principal (État, ONG, Commune, Paysans), le levier financier possible, et l'indicateur de suivi.
\item Expliquer pourquoi l'implication des populations riveraines (comités de bassin) est une condition \textit{sine qua non} de l'efficacité de ce plan.
\end{enumerate}
\end{enumerate}
\end{exercice}

\begin{exercice}[Sujet type Probatoire : Exploitation minière et biodiversité au Sud-Cameroun]
Un projet d'exploitation de fer (type Mbalam/Nabeba) et de cobalt/nickel est envisagé dans la région du Sud, zone de forêt dense humide à haute biodiversité (proximité du Parc National de Campo Ma'an). L'exploitation prévue est à ciel ouvert (mine à ciel ouvert) sur \SI{500}{ha} avec usine de traitement sur site, route d'évacuation (\SI{150}{km}) et port minéralier.
\begin{enumerate}
\item \textbf{Impacts environnementaux (7 pts) :}
\begin{enumerate}
\item Lister \textbf{cinq} impacts négatifs majeurs de ce type de projet sur : l'eau (qualité/drainage acide), les sols (stériles, pollution), la forêt (fragmentation, perte d'habitat), la faune (couloirs de migration, braconnage induit), et l'air/paysage (poussières, bruit).
\item Expliquer le phénomène de \textbf{drainage minier acide (DMA)} : réaction chimique principale (pyrite + eau + oxygène), conséquences sur le pH et la solubilisation des métaux lourds (Fe, As, Pb), impact sur les cours d'eau aval.
\end{enumerate}
\item \textbf{Mesures d'atténuation et compensation (8 pts) :}
La réglementation camerounaise (Loi 96/12, Décret 2013/0171/PM) impose une Étude d'Impact Environnemental et Social (EIES) et un Plan de Gestion Environnementale et Sociale (PGES).
\begin{enumerate}
\item Différencier les mesures d'\textbf{évitement}, de \textbf{réduction}, de \textbf{réhabilitation} et de \textbf{compensation} (séquence ERC). Donner un exemple concret pour chacune appliqué à ce projet minier.
\item Le projet prévoit une compensation biodiversité : création d'une nouvelle aire protégée de \SI{10 000}{ha} en zone dégradée. Discuter la pertinence écologique de cette compensation (équivalence écologique, additionnalité, pérennité, fuite de pression).
\end{enumerate}
\item \textbf{Gouvernance et acceptabilité sociale (5 pts) :}
\begin{enumerate}
\item Citer trois droits des populations autochtones (Bagyeli/Bakola) et locales que le projet doit respecter (Consentement Libre, Préalable et Éclairé - CLPE, partage des bénéfices, droits coutumiers).
\item Proposer un mécanisme de suivi indépendant (comité de surveillance multi-acteurs) pour garantir le respect du PGES sur la durée de vie de la mine (25 ans).
\end{enumerate}
\end{enumerate}
\end{exercice}

\newpage

\coursec{Corrigés des exercices}

\begin{corrige}[Exercice 1 — QCM : Ressources et pressions]
La seule affirmation correcte est la \textbf{4}.

\textbf{Justifications de l'élimination :}
\begin{enumerate}
\item \textbf{Fausse.} Le cycle de l'eau renouvelle l'eau douce, mais à un \cle{rythme limité} : si les prélèvements (agriculture, industrie, villes) dépassent la vitesse de recharge des nappes et des cours d'eau, la ressource s'épuise localement. De plus, la pollution peut la rendre inutilisable. L'eau douce est donc une ressource \emph{renouvelable mais vulnérable}, non inépuisable.
\item \textbf{Fausse.} Les cendres n'apportent qu'une fertilité \emph{éphémère} (quelques saisons) : le brûlis détruit la matière organique et la faune du sol, puis l'érosion et le lessivage emportent les éléments nutritifs. La fertilité chute ensuite durablement.
\item \textbf{Fausse.} La biodiversité comprend \textbf{trois niveaux} : la diversité \emph{génétique} (au sein des espèces), la diversité \emph{spécifique} (espèces animales \textbf{et} végétales, champignons, micro-organismes) et la diversité \emph{écosystémique} (variété des milieux).
\item \textbf{Vraie.} L'exploitation sélective (prélèvement ciblé d'arbres matures), le reboisement et le respect des cycles de rotation garantissent que le prélèvement ne dépasse pas la capacité de renouvellement de la forêt : c'est la définition même de la \cle{gestion raisonnée}.
\end{enumerate}
\end{corrige}

\begin{corrige}[Exercice 2 — Vrai / Faux justifié]
\begin{enumerate}
\item \textbf{Faux.} Le labour en \cle{courbes de niveau} (perpendiculaire à la pente) \emph{ralentit} le ruissellement et donc \emph{réduit} l'érosion hydrique. C'est le labour dans le sens de la pente qui crée des chenaux d'écoulement et accélère l'érosion.
\item \textbf{Faux.} Le braconnage déséquilibre tout l'écosystème : la disparition des grands herbivores modifie la végétation, celle des prédateurs provoque la pullulation de proies, et les petites espèces (oiseaux, rongeurs, insectes) subissent aussi la pression. Chaque espèce joue un rôle dans les chaînes alimentaires.
\item \textbf{Vrai.} C'est la définition du développement durable (Sommet de Rio, 1992) : un développement qui répond aux besoins du présent sans compromettre ceux des générations futures, en conciliant \textbf{économie viable}, \textbf{société équitable} et \textbf{environnement préservé}.
\item \textbf{Faux.} Les minerais se forment par des processus géologiques sur des \emph{millions d'années} : à l'échelle humaine, les stocks sont fixes et s'épuisent. Ce sont des ressources \cle{non renouvelables}, quelle que soit la durée d'attente envisageable.
\end{enumerate}
\end{corrige}

\begin{corrige}[Exercice 3 — Définitions contextualisées]
\begin{enumerate}
\item \textbf{Ressource naturelle renouvelable} : élément du milieu naturel utilisable par l'Homme et dont le stock se reconstitue à l'échelle d'une vie humaine, à condition que les prélèvements n'excèdent pas la vitesse de renouvellement. \emph{Exemple :} le poisson de la Sanaga, l'eau du fleuve Logone, le bois des forêts de l'Est.
\item \textbf{Gestion raisonnée (durable)} : mode d'exploitation qui prélève la ressource à un rythme inférieur ou égal à sa capacité de régénération, en préservant les milieux et en associant les usagers, afin de garantir sa disponibilité pour les générations futures. \emph{Exemple :} les forêts communautaires de la région de l'Est avec plans d'aménagement et reboisement.
\item \textbf{Érosion des sols} : dégradation et déplacement des particules du sol sous l'action de l'eau (ruissellement), du vent ou des pratiques agricoles, entraînant la perte de l'horizon fertile superficiel. \emph{Exemple :} les ravins creusés sur les versants dénudés des collines de Yaoundé après les pluies.
\item \textbf{Aire protégée} : espace terrestre ou marin délimité par la loi, géré en vue de conserver la biodiversité, les paysages et les services écosystémiques. \emph{Exemple :} le Parc National de Waza, la réserve de faune du Dja (patrimoine mondial de l'UNESCO).
\end{enumerate}
\end{corrige}

\begin{corrige}[Exercice 4 — Calcul direct : Taux de déforestation]
\begin{enumerate}
\item \textbf{Superficie perdue :}
\[ 120\,000 - 96\,000 = \SI{24 000}{ha} \]
\item \textbf{Taux moyen annuel} (perte linéaire sur 10 ans, rapportée à la superficie initiale) :
\[ \text{Perte annuelle} = \dfrac{24\,000}{10} = \SI{2 400}{ha/an} \qquad \text{Taux} = \dfrac{2\,400}{120\,000} \times 100 = \SI{2}{\% / an} \]
\item \textbf{Perte moyenne annuelle :} \SI{2 400}{ha/an}, soit l'équivalent d'environ 3\,400 terrains de football perdus chaque année.
\end{enumerate}
\end{corrige}

\begin{attention}[Point de vigilance]
Le taux se calcule par rapport à la superficie \emph{initiale} (2010), pas à la superficie finale. Une erreur classique consiste aussi à diviser la perte totale par la superficie finale : $\frac{24\,000}{96\,000} = 25\,\%$ serait le taux cumulé rapporté à 2020, non le taux annuel demandé.
\end{attention}

\begin{corrige}[Exercice 5 — Analyse de document : Exploitation forestière dans l'Est-Cameroun]
\begin{enumerate}
\item \textbf{Soldes nets (exploitée $-$ reboisée) :}
\begin{center}
\begin{tabular}{|c|c|c|c|c|}\hline
\rowcolor{popblueL} \textbf{Année} & \textbf{Exploitée (ha)} & \textbf{Reboisée (ha)} & \textbf{Solde net (ha)} & \textbf{Solde cumulé (ha)} \\ \hline
2018 & 2 500 & 200 & 2 300 & 2 300 \\ \hline
2019 & 2 800 & 300 & 2 500 & 4 800 \\ \hline
2020 & 3 000 & 250 & 2 750 & 7 550 \\ \hline
2021 & 3 200 & 400 & 2 800 & 10 350 \\ \hline
2022 & 3 500 & 350 & 3 150 & \textbf{13 500} \\ \hline
\end{tabular}
\end{center}
Solde cumulé : $2\,300+2\,500+2\,750+2\,800+3\,150 = \SI{13 500}{ha}$ de forêt perdue nette en 5 ans.
\item \textbf{Graphique attendu :} deux courbes croissantes, celle des surfaces exploitées très au-dessus (pente forte) de celle des surfaces reboisées (presque plate) ; l'écart entre les deux courbes se creuse chaque année.
\item \textbf{Non, cette gestion n'est pas durable.} Deux critères non respectés :
\begin{itemize}
\item \textbf{Équilibre prélèvement/régénération :} le taux de reboisement ne couvre que 8 à 11\,\% des surfaces exploitées (ex. 2022 : $350/3\,500 = 10\,\%$), alors qu'une gestion durable impose de compenser l'intégralité des coupes.
\item \textbf{Prélèvement croissant sans limite :} le volume prélevé augmente chaque année ($+38\,\%$ entre 2018 et 2022) sans tenir compte de la capacité de renouvellement naturel de la forêt.
\end{itemize}
\item \textbf{Mesures :} (a) imposer un \cle{plan d'aménagement} avec reboisement obligatoire d'au moins 100\,\% des surfaces coupées et respect d'une rotation de 25--30 ans ; (b) plafonner les volumes prélevés (quota annuel fixé sur la base du taux de régénération) et passer à l'exploitation sélective des seuls arbres matures.
\end{enumerate}
\end{corrige}

\begin{corrige}[Exercice 6 — Construction de courbe : Pêche dans le fleuve Sanaga]
\begin{enumerate}
\item \textbf{Courbe :} nuage de points régulièrement décroissant, que l'on relie ; on peut modéliser la tendance par une droite descendante (régression approximative).
\item \textbf{Allure :} décroissance quasi continue et régulière sur 15 ans, sans aucune année de remontée. Pourcentage de baisse :
\[ \dfrac{12\,500 - 6\,000}{12\,500} \times 100 = \dfrac{6\,500}{12\,500} \times 100 = \SI{52}{\%} \]
Les captures ont été \textbf{divisées par plus de deux} en 15 ans.
\item \textbf{Conclusion :} oui, il s'agit d'une \cle{surexploitation} de la ressource halieutique : les prélèvements dépassent durablement la capacité de reproduction et de croissance des stocks de poissons. Une ressource gérée durablement montrerait des captures stables ou fluctuant autour d'une moyenne, pas un effondrement de 52\,\%.
\item \textbf{Mesures (deux au choix) :}
\begin{itemize}
\item \textbf{Période de fermeture (fermeture biologique) :} interdiction de pêche pendant la saison de reproduction, afin que les géniteurs se reproduisent et que les juvéniles grandissent avant d'être capturés.
\item \textbf{Taille minimale de capture :} remise à l'eau des poissons n'ayant pas atteint la taille de première maturité sexuelle, garantissant que chaque poisson s'est reproduit au moins une fois.
\item (Autres réponses acceptées : quotas de capture, maillage minimal des filets, protection des frayères.)
\end{itemize}
\end{enumerate}
\end{corrige}

\begin{corrige}[Exercice 7 — Étude de cas : Gestion de l'eau à Yaoundé]
\begin{enumerate}
\item \textbf{Quatre pressions humaines :}
\begin{itemize}
\item \textbf{Surconsommation} liée à la croissance démographique rapide (demande supérieure à l'offre) ;
\item \textbf{Pollution} des cours d'eau (Mfoundi, Mefou) par les déchets solides et liquides ;
\item \textbf{Déforestation des bassins versants} (Nkolbisson, Nkolandom), qui réduit l'infiltration et la recharge des nappes et accroît le ruissellement érosif ;
\item \textbf{Pertes techniques} : fuites du réseau (rendement de 60\,\%, soit 40\,\% de l'eau produite perdue).
\end{itemize}
\item \textbf{Mesures associées :}
\begin{itemize}
\item Surconsommation $\to$ campagnes d'économie d'eau, tarification progressive incitative, compteurs individuels ;
\item Pollution $\to$ collecte organisée des ordures, stations d'épuration, interdiction des rejets bruts, latrines améliorées ;
\item Déforestation $\to$ reboisement des versants, classement des collines en zones non constructibles ;
\item Fuites $\to$ programme de détection et réparation des fuites, renouvellement des conduites vétustes.
\end{itemize}
\item \textbf{Pourquoi protéger les bassins versants est plus durable :} augmenter la capacité de traitement ne fait que \emph{soigner le symptôme} (traiter une eau de plus en plus polluée, à coût croissant), sans reconstituer la ressource. Protéger le bassin versant agit sur la \emph{cause} : la végétation favorise l'infiltration (recharge des nappes, régulation des débits en saison sèche), retient les sédiments et filtre naturellement les polluants. C'est une solution pérenne, peu coûteuse en fonctionnement, qui fournit en outre des services annexes (biodiversité, lutte contre l'érosion, régulation du climat local) : elle satisfait les trois piliers du développement durable.
\end{enumerate}
\end{corrige}

\begin{corrige}[Exercice 8 — Synthèse rédactionnelle : Le brûlis et les alternatives]
\textbf{Corrigé type (environ 150 mots) :}

Le brûlis répété détruit la fertilité du sol par plusieurs mécanismes en chaîne. La combustion élimine la \cle{matière organique} et tue la faune du sol (vers de terre, micro-organismes) qui aère et décompose l'humus ; les cendres n'offrent qu'une fertilité éphémère d'une ou deux saisons. Le sol nu, sans couvert végétal, est ensuite exposé au \cle{ruissellement} qui emporte l'horizon fertile (érosion) et aux pluies qui \cle{lessivent} les éléments nutritifs en profondeur, hors de portée des racines. Les rendements s'effondrent et la jachère ne peut plus se régénérer.

Trois alternatives agroécologiques sont possibles. L'\textbf{agroforesterie} associe arbres (fruitiers, légumineuses fixatrices d'azote) et cultures sur la même parcelle : les arbres restituent de la matière organique et protègent le sol. Le \textbf{compostage} recycle les résidus de culture et déjections animales en engrais organique local et gratuit. La \textbf{jachère améliorée} avec plantes de couverture (mucuna, crotalaire) restaure l'azote et étouffe les adventices.

Cette transition sécurise les rendements à long terme, réduit les dépenses d'intrants et préserve le capital foncier de la communauté : c'est un gain à la fois économique et écologique.
\end{corrige}

\begin{attention}[Points de vigilance — rédaction]
L'examinateur attend : (1) au moins \textbf{trois mécanismes} de dégradation nommés précisément (matière organique, faune du sol, érosion, lessivage) ; (2) trois alternatives avec le \textbf{principe} de chacune (pas une simple liste) ; (3) une conclusion articulant intérêt \emph{économique} ET \emph{écologique}. Éviter les généralités non argumentées (« le brûlis c'est mauvais »).
\end{attention}

\begin{corrige}[Exercice 9 — Sujet type Probatoire : Le Parc National de Waza]
\textbf{1. Analyse des pressions (6 pts)}
\begin{enumerate}
\item \textbf{Classification (2 pts) :}
\begin{itemize}
\item \emph{Pressions climatiques/environnementales :} baisse de la pluviométrie (700 $\to$ \SI{550}{mm/an}), sécheresses récurrentes, envasement du Logone (processus naturel amplifié).
\item \emph{Pressions anthropiques directes :} barrage de Maga (modification hydrologique), braconnage, pression pastorale transhumante, riziculture irriguée sur les marges (prélèvement d'eau et emprise sur l'espace).
\end{itemize}
\item \textbf{Barrage de Maga $\to$ baisse du Kob (2 pts) :} le barrage retient les crues du Logone, supprimant les inondations saisonnières de la plaine de Waza. Or cette plaine inondée produisait, à la décrue, des pâturages herbeux riches (yaéré) dont dépend le Kob de Buffon en saison sèche. L'assèchement supprime cette ressource alimentaire clé : faute de nourriture et d'eau, les kobs meurent, migrent ou se reproduisent moins — d'où l'effondrement de 30\,000 à 5\,000 individus.
\item \textbf{Éléphant, espèce « ingénieure » (2 pts) :} l'éléphant façonne la savane : il abat des arbres (limitant l'embroussaillement et maintenant les prairies pour les herbivores), disperse les graines via ses fèces, creuse des points d'eau utilisés par d'autres espèces, et ouvre des pistes dans la végétation dense. Sa disparition provoque des modifications en cascade (effet trophique et paysager) touchant des dizaines d'espèces : l'impact est donc disproportionné par rapport à son seul effectif.
\end{enumerate}

\textbf{2. Évaluation de la durabilité (4 pts)}
L'approche « forteresse » (interdiction totale + surveillance armée insuffisante) échoue sur les trois piliers :
\begin{itemize}
\item \textbf{Pilier social :} exclusion des riverains de leurs zones de pêche, pâture et cueillette traditionnelles $\to$ appauvrissement, ressentiment, conflits homme-faune, braconnage de subsistance.
\item \textbf{Pilier économique :} coût élevé et permanent de la surveillance pour un État aux moyens limités ; aucun bénéfice local tiré du parc, donc aucune incitation à le protéger.
\item \textbf{Pilier environnemental :} malgré les moyens engagés, les effectifs fauniques s'effondrent : l'approche est inefficace car elle ne s'attaque pas aux causes (pauvreté, besoin de pâtures, demande d'ivoire).
\end{itemize}

\textbf{3. Plan d'action intégré (10 pts) — exemple de réponse attendue :}
\begin{itemize}
\item \textbf{Axe 1 — Gouvernance participative et populations locales :} (a) créer \textbf{10 comités villageois de gestion} de la périphérie du parc (COVAREF) d'ici fin 2025 ; (b) former et salarier \textbf{30 écogardes recrutés localement} d'ici 2026, et développer l'élevage de proximité hors du parc (2 ranchs communautaires d'ici 2027) pour réduire la transhumance interne.
\item \textbf{Axe 2 — Restauration écologique et action de l'État :} (a) réhabiliter la plaine d'inondation par des \textbf{lâchures d'eau concertées} depuis le barrage de Maga (protocole signé État-gestionnaire d'ici 2025, 3 lâchures pilotes en saison sèche 2026) ; (b) créer \textbf{5 points d'eau artificiels} et renforcer la lutte anti-braconnage (effectif d'écogardes doublé d'ici 2026, objectif : zéro éléphant braconné en 2027).
\item \textbf{Axe 3 — Écotourisme et partage des bénéfices :} (a) former \textbf{20 jeunes guides éco-touristiques} riverains d'ici 2026 et construire 2 écolodges gérés en partie par les communautés ; (b) instituer un \textbf{fonds de reversement} de 30\,\% des recettes touristiques aux villages (objectif : 5\,000 visiteurs/an en 2028), rendant la faune vivante plus rentable que le braconnage.
\end{itemize}
\end{corrige}

\begin{attention}[Barème indicatif et points de vigilance]
Barème : Q1a 2 pts ; Q1b 2 pts ; Q1c 2 pts ; Q2 4 pts (au moins 2 limites reliées aux piliers) ; Q3 10 pts (3 axes $\times$ 2 actions chiffrées et datées : 1 pt par action + cohérence globale). Exigences : les actions doivent être \textbf{mesurables et datées} (nombre, échéance) ; une action vague (« sensibiliser la population ») sans cible ni date ne vaut pas le point. Le lien causal barrage $\to$ plaine inondable $\to$ pâture $\to$ kob doit être explicite.
\end{attention}

\begin{corrige}[Exercice 10 — Sujet type Probatoire : Bassin versant de la Mefou]
\textbf{1. Diagnostic environnemental (6 pts)}
\begin{enumerate}
\item \textbf{Liens usage du sol $\to$ pollution (3 pts) :}
\begin{itemize}
\item \textbf{Agriculture (40\,\%) $\to$ nitrates et turbidité :} les engrais azotés non absorbés par les cultures sont \emph{lessivés} par les pluies vers la retenue (les nitrates, très solubles, s'infiltrent et ruissellent) ; les sols nus labourés s'érodent, chargeant l'eau en particules (turbidité de \SI{150}{NTU}).
\item \textbf{Habitat non planifié (25\,\%) $\to$ coliformes fécaux :} absence d'assainissement (latrines sommaires, rejets bruts) : les eaux usées chargées de bactéries d'origine fécale (\emph{Escherichia coli}...) ruissellent vers le cours d'eau ; les déchets solides aggravent la contamination.
\item \textbf{Perte de forêt et de couvert végétal (savane dégradée, forêt résiduelle réduite à 15\,\%) $\to$ turbidité :} sans racines ni litière, le ruissellement s'accélère et arrache les particules du sol (érosion hydrique), augmentant la turbidité et l'envasement de la retenue.
\end{itemize}
\item \textbf{Calcul du surcoût de traitement (3 pts) :}
\begin{align*}
\text{Masse journalière} &= 30\,000\ \text{m}^3/\text{j} \times 15\ \text{mg/L} = 30\,000\,000\ \text{L/j} \times 15\ \text{mg/L} = 450\,000\,000\ \text{mg/j} = \SI{450}{kg/j} \\
\text{Masse annuelle} &= 450 \times 365 = \SI{164 250}{kg/an} \\
\text{Coût annuel} &= 164\,250 \times 300 = \SI{49 275 000}{FCFA/an}
\end{align*}
Soit près de \textbf{49,3 millions de FCFA par an} de surcoût, uniquement à cause de la turbidité excessive — argument économique fort pour investir dans la protection du bassin versant.
\end{enumerate}

\textbf{2. Stratégie de gestion raisonnée (8 pts)}
\begin{enumerate}
\item \textbf{Définition du PGIRE (2 pts) :} la Gestion Intégrée des Ressources en Eau est une démarche de gestion \emph{concertée} de l'eau et des milieux à l'échelle du \cle{bassin versant}, qui coordonne tous les usages (potable, agricole, industriel, écologique) et tous les acteurs (État, communes, usagers) pour garantir une utilisation équitable et durable de la ressource, sans compromettre les écosystèmes aquatiques.
\item \textbf{Quatre mesures (4 pts) — exemple de tableau attendu :}
\begin{center}
\begin{tabularx}{\linewidth}{|X|l|X|X|}\hline
\rowcolor{popblueL} \textbf{Mesure} & \textbf{Acteur} & \textbf{Financement} & \textbf{Indicateur de suivi} \\ \hline
Reboisement des berges (bande riveraine de 30 m) & ONG + paysans & Subvention MINEE / PSE & ha reboisés ; turbidité (NTU) \\ \hline
Latrines écologiques dans l'habitat riverain & Commune & Budget communal + bailleurs & nombre de latrines ; UFC/100 mL \\ \hline
Zonage agricole : interdiction des cultures sur berges, promotion de l'agroforesterie & Paysans encadrés (MINADER) & Paiement pour services environnementaux (PSE) & nitrates (mg/L) ; ha en agroforesterie \\ \hline
Fossés d'infiltration et seuils de rétention sur les versants & État (MINEE) & Budget d'État & volume retenu ; débit de crue \\ \hline
\end{tabularx}
\end{center}
\item \textbf{Rôle des comités de bassin (2 pts) :} les riverains sont à la fois les \emph{auteurs} et les \emph{victimes} des pressions : aucune mesure (zonage, reboisement, assainissement) ne peut être appliquée durablement sans leur adhésion. Les comités de bassin associent les usagers aux décisions (légitimité), mobilisent les savoirs locaux, assurent la surveillance quotidienne (gratuite et permanente, contrairement aux contrôles étatiques ponctuels) et garantissent l'équité (pilier social), condition de la pérennité du plan. Sans eux, les mesures restent lettre morte ou sont contournées.
\end{enumerate}
\end{corrige}

\begin{attention}[Barème indicatif et points de vigilance]
Barème : Q1a 3 pts (1 pt par pollution correctement reliée à un mécanisme) ; Q1b 3 pts (conversion mg $\to$ kg : 1 pt ; calcul : 2 pts) ; Q2a 2 pts ; Q2b 4 pts (1 pt par mesure complète : acteur + financement + indicateur) ; Q2c 2 pts. Piège classique du calcul : oublier que $1\ \text{m}^3 = 1\,000\ \text{L}$, ou diviser au lieu de multiplier par 365. Vérification : $450 \times 365 = 164\,250$ et $164\,250 \times 300 = 49\,275\,000$.
\end{attention}

\begin{corrige}[Exercice 11 — Sujet type Probatoire : Exploitation minière au Sud-Cameroun]
\textbf{1. Impacts environnementaux (7 pts)}
\begin{enumerate}
\item \textbf{Cinq impacts majeurs (2,5 pts) :}
\begin{itemize}
\item \textbf{Eau :} drainage minier acide et pollution des cours d'eau par les métaux lourds ; modification du réseau hydrographique (détournements, envasement).
\item \textbf{Sols :} décapage et stockage de stériles et de résidus de traitement (pollution durable, stérilité des terres).
\item \textbf{Forêt :} défrichement de \SI{500}{ha} et surtout \emph{fragmentation} par la route de \SI{150}{km} (lisières, îlots forestiers isolés).
\item \textbf{Faune :} rupture des couloirs de migration (éléphants, grands singes de Campo Ma'an), mortalité routière, braconnage induit par l'afflux de travailleurs et l'ouverture des pistes.
\item \textbf{Air/paysage :} poussières et bruit des tirs et engins, dégradation paysagère du cratère minier.
\end{itemize}
\item \textbf{Drainage minier acide (2 pts) :} la mise au jour de la \cle{pyrite} (sulfure de fer, \ce{FeS2}) contenue dans les roches et stériles l'expose à l'air et à l'eau ; elle s'oxyde selon la réaction :
\[ \ce{4FeS2 + 15O2 + 14H2O -> 4Fe(OH)3 + 8H2SO4} \]
L'\textbf{acide sulfurique} produit abaisse fortement le pH des eaux de ruissellement (pH pouvant descendre sous 3). Cette acidité \emph{solubilise les métaux lourds} (Fe, As, Pb, Cd...) normalement piégés dans les roches : ils deviennent biodisponibles et toxiques. Les cours d'eau aval subissent alors une mortalité des poissons et invertébrés, une eau impropre à la consommation et à l'irrigation — pollution qui peut durer des \emph{décennies après la fermeture} de la mine.
\end{enumerate}

\textbf{2. Mesures d'atténuation et compensation (8 pts)}
\begin{enumerate}
\item \textbf{Séquence ERC (4 pts) :}
\begin{itemize}
\item \textbf{Évitement} : supprimer l'impact à la source. \emph{Ex. :} déplacer le tracé de la route pour ne pas traverser le couloir de migration des éléphants ni le noyau du parc.
\item \textbf{Réduction} : diminuer l'intensité de l'impact résiduel. \emph{Ex. :} confiner les stériles sous couverture étanche et traiter les eaux acides à la chaux avant rejet ; arroser les pistes contre les poussières.
\item \textbf{Réhabilitation} : restaurer le site après exploitation. \emph{Ex. :} remblayer le cratère, remettre la terre végétale stockée et reboiser avec des essences locales.
\item \textbf{Compensation} : contrebalancer les impacts résiduels irréductibles par un gain équivalent ailleurs. \emph{Ex. :} créer l'aire protégée de \SI{10 000}{ha} et financer la surveillance renforcée de Campo Ma'an.
\end{itemize}
\item \textbf{Discussion de la compensation (4 pts) :} sa pertinence est \emph{conditionnelle} :
\begin{itemize}
\item \textbf{Équivalence écologique :} une forêt secondaire sur zone dégradée ne remplace pas à court terme une forêt primaire à haute biodiversité ; le « comme pour comme » est difficile à atteindre (délai de plusieurs décennies, pertes irréversibles d'espèces endémiques).
\item \textbf{Additionnalité :} l'aire protégée ne compense réellement que si elle protège un milieu qui aurait été dégradé \emph{sans} le projet ; si la zone était déjà protégée de fait, il n'y a aucun gain net.
\item \textbf{Pérennité :} le financement de la gestion doit être garanti au-delà des 25 ans de la mine (fonds dédié, statut juridique solide).
\item \textbf{Fuite de pression :} protéger une zone peut déplacer le braconnage et la déforestation vers les forêts voisines non protégées — il faut un suivi à l'échelle du paysage.
\end{itemize}
\end{enumerate}

\textbf{3. Gouvernance et acceptabilité sociale (5 pts)}
\begin{enumerate}
\item \textbf{Trois droits des populations (3 pts) :}
\begin{itemize}
\item \textbf{CLPE} (Consentement Libre, Préalable et Éclairé) : les communautés Bagyeli/Bakola doivent être informées de façon compréhensible et donner leur accord \emph{avant} tout démarrage, sans pression ;
\item \textbf{Partage des bénéfices} : redevances, emplois locaux, infrastructures (écoles, centres de santé) et fonds de développement communautaire ;
\item \textbf{Respect des droits coutumiers} : accès aux terres, sites sacrés et ressources de subsistance (chasse, cueillette) ; indemnisation équitable en cas de déplacement.
\end{itemize}
\item \textbf{Mécanisme de suivi indépendant (2 pts) :} créer un \cle{comité de surveillance multi-acteurs} réunissant État (MINEE, MINFOF, MINEPDED), représentants élus des communautés (y compris autochtones), ONG environnementales, scientifiques indépendants et l'entreprise. Il disposerait d'un \emph{fonds de suivi abondé par le minier}, publierait des rapports publics semestriels (qualité de l'eau, reboisement, emplois locaux), aurait accès libre au site et pourrait saisir l'autorité administrative en cas de non-respect du PGES — avec un mandat garanti sur toute la durée de vie de la mine et après fermeture.
\end{enumerate}
\end{corrige}

\begin{attention}[Barème indicatif et points de vigilance]
Barème : Q1a 2,5 pts (0,5 par impact) ; Q1b 2 pts (équation 1 pt, explication pH/métaux 1 pt) ; Q2a 4 pts (1 pt par catégorie correctement définie ET illustrée) ; Q2b 4 pts (1 pt par critère discuté) ; Q3a 3 pts ; Q3b 2 pts. Points de vigilance : ne pas confondre \emph{réhabilitation} (restaurer le site) et \emph{compensation} (gain ailleurs) ; l'équation du DMA doit être équilibrée ; pour le CLPE, les trois adjectifs (libre, préalable, éclairé) sont exigés.
\end{attention}

\newpage

\coursec{Fiche bilan}

\begin{popfigure}
\begin{tikzpicture}[
  mindmap,
  grow cyclic,
  every node/.style={concept, text=white, font=\small\bfseries},
  root concept/.append style={concept color=popdark, font=\bfseries},
  level 1/.append style={level distance=3.6cm, sibling angle=51.4},
  level 2/.append style={level distance=2.1cm, sibling angle=40, font=\footnotesize},
]
\node[root concept]{Ressources naturelles et développement durable}
  child[concept color=popblue]{ node{Ressources naturelles}
    child{ node{renouvelables} }
    child{ node{non renouvelables} }
  }
  child[concept color=popgreen]{ node{Pressions sur l'eau et les sols}
    child{ node{pollution} }
    child{ node{érosion, épuisement} }
  }
  child[concept color=poporange]{ node{Pressions sur forêts et biodiversité}
    child{ node{déforestation} }
    child{ node{braconnage} }
  }
  child[concept color=poppurple]{ node{Développement durable}
    child{ node{3 piliers} }
    child{ node{gestion raisonnée} }
  }
  child[concept color=poppink]{ node{Exemples camerounais}
    child{ node{forêts, parcs} }
    child{ node{barrages, eau} }
  }
  child[concept color=popgold]{ node{Méthode}
    child{ node{analyser un document} }
    child{ node{proposer des mesures} }
  };
\end{tikzpicture}
\legende{Carte mentale de la leçon : les ressources naturelles, les pressions humaines qu'elles subissent, la notion de développement durable, les exemples camerounais et la démarche d'analyse d'une situation.}
\end{popfigure}

\begin{bilan}
\textbf{Définitions clés}
\begin{itemize}
  \item \cle{Ressource naturelle} : élément de la nature (eau, sol, forêt, espèces vivantes, minerais) utilisé par l'Homme pour satisfaire ses besoins.
  \item \cle{Ressource renouvelable} : ressource qui se régénère à l'échelle d'une vie humaine si son rythme d'exploitation ne dépasse pas son rythme de renouvellement (eau, forêt, biodiversité).
  \item \cle{Ressource non renouvelable} : ressource dont le stock est limité et qui ne se reconstitue pas à l'échelle humaine (pétrole, minerais).
  \item \cle{Développement durable} : développement qui répond aux besoins du présent sans compromettre la capacité des générations futures à répondre aux leurs (Rapport Brundtland, 1987).
  \item \cle{Gestion raisonnée} : exploitation planifiée et contrôlée d'une ressource de façon à la maintenir dans le temps.
  \item \cle{Biodiversité} : diversité des espèces vivantes, des gènes et des écosystèmes.
\end{itemize}

\textbf{Les trois piliers du développement durable}
\begin{center}
\begin{tabularx}{\linewidth}{|l|X|}\hline
\rowcolor{popblueL} \textbf{Pilier} & \textbf{Objectif} \\ \hline
Environnemental & protéger les ressources et les écosystèmes \\ \hline
Économique & assurer une activité économique viable et efficace \\ \hline
Social & garantir l'équité et le bien-être des populations \\ \hline
\end{tabularx}
\end{center}

\textbf{Pressions humaines et conséquences à connaître par cœur}
\begin{center}
\begin{tabularx}{\linewidth}{|l|X|X|}\hline
\rowcolor{popgreenL} \textbf{Ressource} & \textbf{Pressions} & \textbf{Conséquences} \\ \hline
Eau & rejets domestiques et industriels, engrais, pesticides, surpompage & pollution, eutrophisation, épuisement des nappes, maladies hydriques \\ \hline
Sols & déforestation, surpâturage, cultures intensives, produits chimiques & érosion, perte de fertilité, désertification \\ \hline
Forêts & exploitation abusive, feux de brousse, agriculture sur brûlis & déforestation, perte d'habitats, dérèglement climatique \\ \hline
Biodiversité & braconnage, destruction des habitats, pollution & disparition d'espèces, déséquilibre des écosystèmes \\ \hline
\end{tabularx}
\end{center}

\textbf{Mesures de gestion durable}
\begin{itemize}
  \item Eau : traitement des eaux usées, protection des sources, irrigation raisonnée (goutte-à-goutte).
  \item Sols : reboisement, jachère, rotation des cultures, agroforesterie, cordons de pierres.
  \item Forêts : reforestation, exploitation sélective avec permis, aires protégées.
  \item Biodiversité : parcs nationaux et réserves, lutte contre le braconnage, sensibilisation.
\end{itemize}

\textbf{Repères camerounais}
\begin{itemize}
  \item Le Cameroun possède environ \num{22} millions d'hectares de forêts (bassin du Congo) ; la déforestation y est liée à l'agriculture itinérante et à l'exploitation forestière.
  \item Aires protégées : parcs de Waza, de la Bénoué, de Korup, réserve de faune du Dja (patrimoine mondial de l'UNESCO).
  \item Barrages hydroélectriques (Songloulou, Édéa, Lagdo) : production d'énergie mais modification des écosystèmes aquatiques.
  \item Dégazage du lac Nyos (1986) : exemple de gestion d'un risque naturel lié à une ressource (eau chargée en \ce{CO2}).
\end{itemize}

\textbf{Méthode express : analyser une situation et proposer des solutions}
\begin{enumerate}
  \item Identifier la ressource concernée et l'activité humaine mise en cause.
  \item Décrire l'impact (nature, ampleur, durée) en s'appuyant sur les données du document.
  \item Relier l'impact à une conséquence sur l'environnement ou la santé.
  \item Proposer des mesures concrètes, adaptées au contexte local, en visant les trois piliers du développement durable.
\end{enumerate}
\end{bilan}

\begin{aretenir}[Checklist avant le Probatoire]
\begin{itemize}
  \item[$\square$] Je sais définir une ressource naturelle et distinguer ressource renouvelable et non renouvelable.
  \item[$\square$] Je connais la définition du développement durable (Rapport Brundtland, 1987).
  \item[$\square$] Je sais citer et expliquer les trois piliers du développement durable.
  \item[$\square$] Je sais identifier les pressions humaines sur l'eau (pollution, surpompage) et leurs conséquences.
  \item[$\square$] Je sais expliquer les causes de l'érosion et de la perte de fertilité des sols.
  \item[$\square$] Je sais citer au moins trois causes de la déforestation au Cameroun.
  \item[$\square$] Je sais expliquer le lien entre destruction des habitats et disparition des espèces.
  \item[$\square$] Je connais au moins deux exemples camerounais de conservation (parcs de Waza, Korup, réserve du Dja).
  \item[$\square$] Je sais proposer des mesures de gestion durable adaptées à une situation locale précise.
  \item[$\square$] Je sais analyser un document (texte, tableau, carte) sur l'exploitation d'une ressource et dégager une conclusion argumentée.
  \item[$\square$] Je rédige mes réponses avec le vocabulaire scientifique exact (érosion, eutrophisation, biodiversité, gestion raisonnée).
  \item[$\square$] Je vérifie que chaque solution proposée est réaliste et respecte l'équilibre environnement, économie, société.
\end{itemize}
\end{aretenir}

\findecours

\end{document}
